\documentclass[11pt]{amsart}
\usepackage{amssymb,amsmath,amsthm}
\usepackage{mathtools}
\usepackage{booktabs}
\usepackage[margin=1.15in]{geometry}
\usepackage[colorlinks=true,linkcolor=blue,citecolor=blue,urlcolor=blue]{hyperref}

\newtheorem{theorem}{Theorem}[section]
\newtheorem{lemma}[theorem]{Lemma}
\newtheorem{proposition}[theorem]{Proposition}
\newtheorem{corollary}[theorem]{Corollary}
\newtheorem{conjecture}[theorem]{Conjecture}
\newtheorem{problem}[theorem]{Problem}
\newtheorem{mainthm}{Theorem}
\renewcommand{\themainthm}{\Alph{mainthm}}
\theoremstyle{definition}
\newtheorem{definition}[theorem]{Definition}
\theoremstyle{remark}
\newtheorem{remark}[theorem]{Remark}
\numberwithin{equation}{section}
\numberwithin{table}{section}
\setlength{\emergencystretch}{3em}

\newcommand{\N}{\mathbb{N}}
\newcommand{\R}{\mathbb{R}}
\newcommand{\C}{\mathbb{C}}
\newcommand{\diff}{\,\mathrm{d}}
\newcommand{\kp}{\mathsf k_p}
\newcommand{\Lt}{\widetilde L}
\newcommand{\cD}{\mathfrak D}
\newcommand{\cA}{\mathfrak A}
\newcommand{\fe}{\mathfrak e}
\newcommand{\fj}{\mathfrak j}
\newcommand{\Nat}{N}
\newcommand{\br}{\mathbf r}
\newcommand{\binomc}[2]{\tbinom{#1}{#2}}
\DeclareMathOperator{\Res}{Res}
\DeclareMathOperator{\PV}{PV}
\DeclareMathOperator{\Var}{Var}
\renewcommand{\Re}{\operatorname{Re}}
\renewcommand{\Im}{\operatorname{Im}}

\title[Jacobi coefficients of measures on reciprocal spectra]{Jacobi coefficients of measures on reciprocal spectra: a
spectral dictionary and the natural index}
\author{Gordon Speagle}
\date{October 2026}

\subjclass[2020]{Primary 42C05; Secondary 47B36, 31A15, 30E15, 11M35}
\keywords{Orthogonal polynomials, Jacobi matrices, recurrence coefficients, discrete measures, constrained equilibrium,
Riemann--Hilbert problems, spectral zeta functions, Hurwitz zeta function, Mellin transforms, Toda hierarchy}

\begin{document}

\begin{abstract}
For an increasing sequence $\Lambda=\{\lambda_m\}$ with $\lambda_m\sim m^p$ we study the recurrence coefficients
$a_n,b_n$ of the discrete measure $\sum_m\lambda_m^{-s/p}\delta_{1/\lambda_m}$; for the lattice $\lambda_m=m^p$,
$n^pb_n\to c(p)$. For the lattice and every $p\ne1$ we prove an expansion of $n^pb_n/c$ to order $n^{-7}+n^{-2-3p}$, a
double series in $1/n$ and $\zeta(-jp)n^{-jp}$ computed by a finite residue recursion and exact computer algebra. In the
natural index $N=n+\frac\alpha2+\frac12$, $\alpha=(s-1-p)/p$, the integer layer is even in $1/N$ to all orders: separating
the roles of $\alpha$ as Bessel order and Szeg\H{o} exponent, we prove a reflection symmetry and, by steepest descent in a
second frame of non-integer index, that the Szeg\H{o} exponent enters only through $N$. The first-order zeta terms are
$\varepsilon_jN^{-1-jp}(1+g_jN^{-2})$, $\varepsilon_j=\binom{2j}{j}p\,c^j\zeta(-jp)$; the zeta layers are Toda flows of the
integer layers, and the leading term of each is an implicit shift of the index. For Hurwitz lattices $(m+\delta)^p$, $-1<\delta\le0$, the
offset only shifts the Szeg\H{o} exponent, so $N$ moves by $\delta$ and $\zeta(-jp)$ becomes $\zeta(-jp,1+\delta)$; for
$(m+\delta)^p+\kappa$, $p>1$, we prove the first pole term. We write the zeta layer as a Mellin--Barnes integral with kernel
$\pi[\cot(\pi wp/2)-\cot(\pi w)]$ and solve explicitly the constrained equilibrium problem with a constraint of infinite
mass at a hard edge; its soft edge carries every pole of the node zeta function with the weight $\pi\cot(\pi q)$. For general
spectra this becomes a conjectural dictionary between $b_n$ and the spectral zeta function of $\Lambda$, tested
numerically on perturbed lattices, Bessel zeros and spectra with multiplicities (with predictions fixed in advance), and on
planar drums and anharmonic oscillators.
\end{abstract}

\maketitle

% Section 1: Introduction.
\section{Introduction}\label{sec:intro}

\subsection{Setting}\label{sec:setting}
Let $p>0$, $s>1$, and let $\Lambda=\{\lambda_m\}_{m\ge1}$ be an increasing sequence of positive numbers with
$\lambda_m\sim m^p$. We call $\Lambda$ a \emph{node set} and study the discrete measure on its reciprocals
\begin{equation}\label{eq:sigmaL}
  \sigma_\Lambda:=\sum_{m\ge1}\lambda_m^{-s/p}\,\delta_{1/\lambda_m}.
\end{equation}
For the \emph{lattice} $\lambda_m=m^p$ this is the measure $\sigma_{p,s}=\sum_kk^{-s}\delta_{k^{-p}}$ of the earlier
papers of this series \cite{PaperI,PaperII,Paper3}. Let $P_n$ be the monic orthogonal polynomials of $\sigma_\Lambda$;
they satisfy
\begin{equation}\label{eq:rec}
  xP_n=P_{n+1}+a_nP_n+b_n^2P_{n-1}\qquad(n\ge0,\ b_0:=0,\ b_n>0),
\end{equation}
and $a_n,b_n$ are the entries of the Jacobi matrix whose spectrum is $\{1/\lambda_m\}\cup\{0\}$. With
\begin{equation}\label{eq:cp}
  I(p):=\frac{\sqrt\pi}{2}\,\frac{\Gamma(\frac1p+\frac12)}{\Gamma(\frac1p+1)},\qquad
  c(p):=\frac14\Bigl(\frac{\pi}{2I(p)}\Bigr)^{p},
\end{equation}
the lattice satisfies the pointwise law $n^pb_n\to c(p)$, $n^pa_n\to2c(p)$, with rate $O(1/n)$ and an explicit $1/n$
term for every $p>0$ \cite[Theorems~1.2 and~1.4]{PaperII}. We write $c=c(p)$,
\begin{equation}\label{eq:alpha-y}
  \alpha:=\frac{s-1-p}{p},\qquad y_n:=\frac{n^pb_n}{c}-1,
\end{equation}
so that $\alpha>-1$ is the Bessel order of the hard edge at $0$ in \cite{PaperII}. The questions of this paper are:
how does the expansion of $y_n$ continue, and how does it depend on $\Lambda$? The organising object is the
\emph{node zeta function}
\begin{equation}\label{eq:Z}
  Z_\Lambda(w):=\sum_{m\ge1}\lambda_m^{-w},
\end{equation}
continued meromorphically from $\Re w>1/p$; for the lattice $Z_\Lambda(w)=\zeta(pw)$. Two quantities recur: the
numbers
\begin{equation}\label{eq:W}
  W(q):=\binom{2q}{q}c^q=\frac{(4c)^q\,\Gamma(q+\frac12)}{\sqrt\pi\,\Gamma(q+1)}\qquad(q\in\C\setminus\{-\tfrac12,-\tfrac32,\dots\}),
\end{equation}
which satisfy $W(1/p)=1$ for every $p$ (this is the definition \eqref{eq:cp} of $c$), and the coefficients
\begin{equation}\label{eq:epsj}
  \varepsilon_j:=W(j)\,p\,\zeta(-jp)=\binom{2j}{j}p\,c^j\zeta(-jp)\qquad(j\ge1),
\end{equation}
so that $\varepsilon_1=\kp:=2pc\,\zeta(-p)$ is the constant of \cite[Theorem~1.4(ii)]{PaperII}.

\subsection{Results}\label{sec:results}
The first result is an expansion of the lattice coefficients far beyond the $1/n$ term. Its form is simplest in the
shifted index
\begin{equation}\label{eq:N}
  N:=n+\frac\alpha2+\frac12 .
\end{equation}
The rational functions $f_4,f_6$ of $(p,\alpha)$ and $g_{1,4}$ are given in \eqref{eq:f4}--\eqref{eq:g14}.

\begin{mainthm}[the lattice to order $n^{-7}+n^{-2-3p}$]\label{thm:A}
Let $p>0$, $p\ne1$, $s>1$. The recurrence coefficients of $\sigma_{p,s}$ satisfy
\begin{multline}\label{eq:thmA}
  \frac{n^pb_n}{c}=\Bigl(\frac nN\Bigr)^p\Bigl[1+\frac{f_2}{N^2}+\frac{f_4}{N^4}+\frac{f_6}{N^6}
  +\frac{\varepsilon_1}{N^{1+p}}\Bigl(1+\frac{g_1}{N^2}+\frac{g_{1,4}}{N^4}\Bigr)\\
  +\sum_{j\ge2}\frac{\varepsilon_j}{N^{1+jp}}\Bigl(1+\frac{g_j}{N^2}\Bigr)
  +\frac{3p+1}{2p}\,\frac{\varepsilon_1^2}{N^{2+2p}}\Bigr]
  +O\bigl(n^{-7}+n^{-2-3p}\bigr),
\end{multline}
where the sum is over the finitely many $j\ge2$ with $1+jp<\min(7,2+3p)$,
\[
  f_2=\frac{p+1}{12}-\frac{(p+2)\alpha^2}{8},\qquad g_j=A_j+B_j\alpha^2,\qquad
  B_j=\frac{(p+2)\bigl(j^2p^2+(j+1)p+2\bigr)}{8(2j-1)p},
\]
\[
  A_j=\frac{j^2(j^2+2j-2)p^3+2(j+1)(2j^2-1)p^2+(5j^2+2j-6)p+2(j-2)}{24(2j-1)p}.
\]
\end{mainthm}

Theorem~\ref{thm:A} is the case $K=6$, $B=2$ of Theorem~\ref{thm:N6}, which gives a double expansion of
$n^pb_n/c$ in $1/n$ and in the parameters $\mu_{n,j}=2\zeta(-jp)/(jn^{jp})$ of the band factor, to every order, through
a finite residue recursion. Its proof combines the third-order Riemann--Hilbert analysis of \cite{PaperII}, extended
to order $K$, with exact computer algebra over $\mathbb Q(p,\alpha)$ and, for the zeta terms, over
$\mathbb Q(p,\alpha,j)$ with $j$ a symbol. Three structural facts are visible in~\eqref{eq:thmA}: the shift
$n\mapsto N$ absorbs the whole $1/n$ term and the $(s-1)^2$ part of the $n^{-2}$ term; every layer computed is even in
$1/N$; and the first-order zeta terms are dressed by an even factor $1+g_jN^{-2}+\dotsb$ with $g_j$ quadratic in $\alpha$.
Theorem~\ref{thm:D} below explains the first two to all orders for the integer layer.
Not every displayed term is a coefficient of $b_n$ for every $p$: Table~\ref{tab:visible} lists which are.

The second result explains the zeta coefficients. The band factor of \cite{PaperII} enters the first-order
correction through one linear functional $e_n$ (Section~\ref{sec:kernel}), whose asymptotic series is
$\sum_j\varepsilon_jn^{-jp}$.

\begin{mainthm}[the zeta layer as a Mellin--Barnes integral]\label{thm:B}
Let $p>0$, $p\ne1$, and $0<c_0<1$. The band-factor functional $e_n$ of the lattice satisfies
\[
  e_n=\frac1{2\pi i}\int_{c_0-i\infty}^{c_0+i\infty}p\,W(w)\,\zeta(-wp)\,n^{-wp}\,
  \pi\Bigl[\cot\frac{\pi wp}2-\cot(\pi w)\Bigr]\diff w ,
\]
and, for every $J\ge1$, $e_n=\sum_{j\le J}\varepsilon_jn^{-jp}+O(n^{-(J+1)p})$. The kernel $-\pi\cot(\pi w)W(w)(4c)^{-w}$
is the Mellin transform of the truncated per-node law $(1-u)^{-1/2}\mathbf 1_{u<1}$ (Lemma~\ref{lem:K0}).
\end{mainthm}

The third result is on the side of potential theory. In the variable $x=n^p\cdot(\text{node})$, the $n$ nodes that
carry the polynomial $P_n$ have the continuum density $\sigma(y)=\frac1py^{-1/p-1}$, and the zero distribution of
$P_n$ is the equilibrium measure of unit mass constrained by $\sigma$ \cite{Rakhmanov,DragnevSaff,PaperI}. For constraints of finite
mass with compact support and no external field, the endpoint equation and the structure of the solution are known
\cite{BeckermannKuijlaars2001,BeckermannHelart} (Section~\ref{sec:literature}). We solve the problem in a new setting, a constraint of
infinite mass at a hard edge, for every constraint with $y\sigma(y)$ decreasing, and give the density in closed form.

\begin{mainthm}[constrained equilibrium]\label{thm:C}
Let $\sigma$ be a positive, locally H\"older continuous density on $(0,Y]$, $0<Y\le\infty$, with $y\sigma(y)$ strictly
decreasing and, if $Y=\infty$, $\int_1^\infty\sigma(y)\log y\diff y<\infty$. Then the soft-edge condition
$\int_a^Y\sigma(y)(1-a/y)^{-1/2}\diff y=1$ has exactly one root $a\in(0,Y)$, and the constrained equilibrium measure is
$\nu\diff x|_{[0,a]}+\sigma\diff y|_{[a,Y]}$ with the explicit density $\nu$ of~\eqref{eq:nu}, $0<\nu<\sigma$. For the
lattice $a=4c(p)$. For a counting function $\mathcal N(\lambda)\sim\sum_\rho c_\rho\lambda^\rho$, the root
$a_n=4c(1+\eta_n)n^{-p}$ of $\int(1-a_n\lambda)^{-1/2}\mathbf 1_{a_n\lambda<1}\diff\mathcal N(\lambda)=n$ satisfies
\[
  n\eta_n=p\sum_{j\ge0}W(j)Z_{\mathcal N}(-j)n^{-jp}+p\sum_qR_qW(q)\,\pi\cot(\pi q)\,n^{-pq}+O(n\eta_n^2),
\]
the second sum over the poles $q\ne-1/p$ of $Z_{\mathcal N}(-w)=\int\lambda^{w}\diff\mathcal N(\lambda)$, with residues
$R_q$ (Proposition~\ref{prop:M12}).
\end{mainthm}

The proof of the equilibrium statement uses no input from outside this series: the positivity of the logarithmic energy of
signed measures of mass zero, without compact support, is \cite[Lemma~3.3]{PaperI}.

The fourth result explains the natural index and the evenness. In the residue recursion the exponent $\alpha$ plays two roles,
the Bessel order $\nu$ of the hard edge and the Szeg\H{o} exponent $\beta$ of the weight; separating them gives formal series
$B(n;\beta,\nu)=\sum_kd_k(p,\beta,\nu)n^{-k}$, the integer layer of a decoupled model (Section~\ref{sec:shiftsec}), with
$B(n;\alpha,\alpha)$ the integer layer of the lattice.

\begin{mainthm}[shift and evenness]\label{thm:D}
Identically in $(p,\beta,\nu)$, $B(n;\beta,\nu)=(n/N)^pF_p(N;\nu)$ with $N=n+\frac{\beta+1}2$ and a formal series $F_p$ that does not
depend on $\beta$; moreover $F_p(N;\nu)=1+\sum_{k\ge1}f_{2k}(p,\nu^2)N^{-2k}$ is even in $N$ and in $\nu$. The zeta layers of weight at
most two ($\mu_1$, $\mu_2$, $\mu_1^2$) obey the same shift. In particular, for the lattice and every $p\ne1$, the integer layer of
$n^pb_n/c$ is $(n/N)^p(1+f_2N^{-2}+f_4N^{-4}+\dotsb)$ to every order, with no odd powers of $1/N$.
\end{mainthm}

This is Theorem~\ref{thm:shift} with Corollary~\ref{cor:even}. The evenness comes from an exact reflection
$d_k(p,-\beta-2,\nu)=(-1)^kd_k(p,\beta,\nu)$ (Theorem~\ref{thm:reflect}, by the grading argument of \cite{KuijlaarsTibboel}); the shift is proved by steepest descent of the same lattice
problem in a second frame, of non-integer index $n+\gamma$, in which the Szeg\H{o} exponent is $\alpha-2\gamma$ and the Bessel order is
still $\alpha$ (Lemma~\ref{lem:frame}), and by comparing the two expansions of $b_n$. In the same model the zeta layers are Toda flows
of the integer layers of $b_n$ and $a_n$, with the times $-2\zeta(-jp)/j$ (Theorem~\ref{thm:toda}): tilting the lattice by
$e^{\sum\tau_jx^j}$ changes only these times, and a Baire-category argument turns the exact Toda equations of the tilted measures into
identities between the layers; the method is that of \cite{EMP} for varying weights, here for a fixed measure. At the hard edge the
dispersionless limit of this flow is a family of Hopf equations, and their characteristics give the leading term of every zeta layer in
closed form, by the genus-zero mechanism of map enumeration \cite{Ercolani,EMP} in this new setting: to leading order in $1/n$ at each order in the times,
$n^pb_n/c$ is $(n/x)^p$ with $x=n-\frac1p\sum_j\varepsilon_jx^{-jp}$, so the zeta times act as an implicit shift of the index
(Theorem~\ref{thm:indexshift}). The diagonal integer layer of $a_n$ is centred at $N+\frac12$, even and independent of the Szeg\H{o}
exponent to order $N^{-4}$ (Proposition~\ref{prop:diag4}).

The fifth result treats the first node sets other than the lattice. For the Hurwitz lattice $\lambda_k=(k+\delta)^p$ the offset
$\delta$ changes only the Szeg\H{o} exponent, $\beta'=\alpha+2\delta$, at fixed Bessel order, and Theorem~\ref{thm:D} turns this into a
shifted natural index.

\begin{mainthm}[Hurwitz lattices]\label{thm:E}
Let $p\ne1$, $-1<\delta\le0$ ($\delta<0$ if $p<1$), $a=1+\delta$ and $1<s<1+2p$. For $\sigma_\Lambda$ with $\lambda_k=(k+\delta)^p$, the
double expansion of $n^pb_n/c$ is that of the decoupled model with $(\beta,\nu)=(\alpha+2\delta,\alpha)$ and band-factor parameters
$2\zeta(-jp,a)/(jn^{jp})$. Consequently expansion~\eqref{eq:thmA} holds with $N$ replaced by $N_\delta=n+\frac\alpha2+\frac12+\delta$ and
$\zeta(-jp)$ by $\zeta(-jp,a)$, with the same $f_{2k}$, $g_j$; the integer layer and the zeta layers of weight at most two are those of
the lattice in $N_\delta$ to all orders.
\end{mainthm}

For the shifted lattices $\lambda_k=(k+\delta)^p+\kappa$, $p>1$, whose node zeta function has a pole at $q=1-\frac1p$, we prove that
$n^pb_n/c$ is the Hurwitz expansion plus the pole term $-2(2-p)c\kappa n^{-p}=pR_qW(q)\pi\cot(\pi q)n^{-p}$ of Theorem~\ref{thm:C}, up
to $O(n^{-1-p})$ (Theorem~\ref{thm:kappa}).

\subsection{The spectral dictionary}\label{sec:dictionary}
For the lattice, Theorems~\ref{thm:A} and~\ref{thm:B} identify the coefficients of $b_n$ with values of $\zeta$, and
Theorem~\ref{thm:E} does the same for Hurwitz lattices with values of $\zeta(\cdot,a)$. For a general node set the
Riemann--Hilbert analysis has not been carried out, and we formulate what the theorems and the numerical evidence suggest
(Section~\ref{sec:conj}). In the natural index $N_\Lambda:=n+\frac\alpha2-Z_\Lambda(0)$
(which is~\eqref{eq:N} for the lattice, $\zeta(0)=-\frac12$), the expansion of $n^pb_n/c$ should read, at leading
order in each spectral datum,
\[
  \Bigl(\frac n{N_\Lambda}\Bigr)^{\!p}\Bigl[1+\frac{f_2}{N_\Lambda^2}+\dotsb
  +\sum_{j\ge1}\frac{W(j)\,p\,Z_\Lambda(-j)}{N_\Lambda^{1+jp}}
  +\sum_q\frac{p\,R_q\,W(q)\,\pi\cot(\pi q)}{N_\Lambda^{1+pq}}+\dotsb\Bigr],
\]
the second sum running over the poles $q$ of $Z_\Lambda(-w)-\zeta(-pw)$, with $w$-residues $R_q$. The integer residues
of $\pi\cot(\pi w)W(w)pZ_\Lambda(-w)$ are the zeta terms of Theorem~\ref{thm:B}; its residues at the poles of $Z_\Lambda$
are the pole terms of Theorem~\ref{thm:C}. In a refined form (Conjecture~\ref{conj:W}) the expansion is the decoupled model with Szeg\H{o} exponent $\alpha-2Z_\Lambda(0)-1$,
flowed along the Toda hierarchy with times read off from $Z_\Lambda(-j)$ and from the density and mass of the nodes, plus the pole
terms; it is proved for Hurwitz lattices and, at first order in the poles, for their constant shifts, and was tested on nine further
measures with predictions fixed before the measures were computed (\S\ref{sec:num-W}); one coefficient outside its claimed
precision failed and led to a refinement (Remark~\ref{rem:Wprime}). The dictionary is exact for the measures on reciprocal squared
Bessel zeros (Proposition~\ref{prop:bessel}), it reads the index offset of the spectrum from $Z_\Lambda(0)$, and it
predicts, for the eigenvalues of $-\frac{\diff^2}{\diff q^2}+q^N$ with $N$ even, an $n^{-2}$ shift equal to
$(N-1)/(6(N+2))$ (Section~\ref{sec:quartic}). For drums at $p=1$, outside the range of the theorems, the perimeter
appears at order $n^{-1/2}$ and the corners at order $n^{-1}$, as the kernel predicts (Section~\ref{sec:drums}). The
dictionary has a limit: subleading terms of each pole, one order lower, are not given by it
(Remark~\ref{rem:sublead}), and for spectra whose nodes deviate from a smooth node function by a fraction of the
spacing the expansion is replaced by noise of relative size $n^{-1/2}$, whose variance appears numerically to be set by the
spectral form factor (Section~\ref{sec:noise}).

\subsection{Organisation and notation}\label{sec:organisation}
Section~\ref{sec:prior} restates the results of the earlier papers that are used, together with two results of the
same project that extend them by one order. Section~\ref{sec:equilibrium} proves Theorem~\ref{thm:C},
Section~\ref{sec:kernel} proves Theorem~\ref{thm:B}, and Section~\ref{sec:natural} proves Theorem~\ref{thm:A} and its
general form. Section~\ref{sec:shiftsec} introduces the decoupled model and proves Theorem~\ref{thm:D}, the Toda
structure of the zeta layers and the index-shift law. Section~\ref{sec:hurwitz-thm} proves Theorem~\ref{thm:E} and the theorem on shifted Hurwitz
lattices; the analytic estimates of these two sections are transfers of those of \cite{PaperII}, which we outline and do not
reproduce. Section~\ref{sec:examples} collects the examples and the numerical evidence, including the pre-registered tests of
Conjecture~\ref{conj:W}, and Section~\ref{sec:conj} the conjectures. Section~\ref{sec:literature} relates the results to the
literature. Appendix~\ref{app:repro} lists the scripts.

Throughout, $p>0$ and $c=c(p)$; $\alpha$, $y_n$, $W$, $\varepsilon_j$, $N$ are as in \eqref{eq:alpha-y}--\eqref{eq:N}.
The Bessel order $\alpha$ is that of \cite{PaperII}, and $\kp=\varepsilon_1$; the constant $\kappa_p$ of
\cite{PaperI} does not occur. We write $\binom{2q}{q}=\Gamma(2q+1)/\Gamma(q+1)^2$ for complex $q$. The letter $N$
denotes the natural index, except in Section~\ref{sec:quartic}, where $q^N$ is a potential. Constants $C$ may change
from line to line. Numbers obtained by fitting floating-point data are quoted with the spread of the fit over
fit windows and nuisance exponents in parentheses; such a spread is not an error bound.

% Section 2: Setting and prior results.
\section{Setting and prior results}\label{sec:prior}

This section fixes the Riemann--Hilbert frame of \cite{PaperII} for the lattice $\sigma_{p,s}$ and restates the
expansions that the later sections refine. Theorem~\ref{thm:imp-onen} is proved in \cite{PaperII}.
Theorem~\ref{thm:E2E3} and Lemmas~\ref{lem:Fprime}--\ref{lem:E37} extend it by one order and to all band-factor
parameters; their proofs are third-order versions of \cite[Sections~8--10]{PaperII} and are not reproduced here.
% TODO (before submission): Theorem 2.3 and Lemmas 2.2, 2.4, 2.5 are proved in the repository notes
% research/G1/RH/second_order/E2_pgt1.md (Theorem E2.1) and E3.md (Theorems E3.1, E3.2, Corollaries E3.3, E3.4,
% Lemmas F.1', E3.4, E3.7), both reviewed, but they are not in any paper. Either write them into an appendix of this
% paper or into a separate paper, and replace this paragraph by a citation.

\subsection{The first two orders}\label{sec:imp}
\begin{theorem}[{\cite[Theorem~1.4]{PaperII}}]\label{thm:imp-onen}
Let $s>1$. If $p\ge1$, then $n^pb_n/c=1-\frac{s-1}{2n}+O(n^{-2})$ and $n^pa_n/(2c)=1-\frac{s-1+p}{2n}+O(n^{-2})$. If
$0<p<1$, both brackets carry the further term $\kp n^{-1-p}$ and the remainder is $O(n^{-1-2p}+n^{-2})$.
\end{theorem}

The $1/n$ coefficient of $n^pb_n/c$ is $d_1:=-(s-1)/2=-p(\alpha+1)/2$. Since $(1+\delta/n)^{-p}=1-p\delta/n+O(n^{-2})$,
it is the first term of the index shift $(n/N)^p$ with $N=n+\frac{\alpha+1}2$; this is the origin of~\eqref{eq:N}.

\subsection{The frame}\label{sec:frame}
We use the notation of \cite[Sections~4, 8 and~10]{PaperII}. After the rescaling $z=n^px$ the polynomials oscillate on the
band $[0,4c]$, with a hard edge at $0$ and a soft edge at $4c$. The weight on the band is
$W_0=\mathsf w_nz^\alpha(4c-z)H_n(z)$, where $\mathsf w_n>0$ is a constant and $H_n$ is the \emph{band factor}, which
collects the discreteness of the nodes. With $\mathcal R(z):=(z(z-4c))^{1/2}$ (branch $\sim z$ at $\infty$),
$\mathsf T:=\frac1{\sqrt2}\bigl(\begin{smallmatrix}1&1\\i&-i\end{smallmatrix}\bigr)$ and
$P_0^\infty=\mathsf T((z-4c)/z)^{\sigma_3/4}\mathsf T^{-1}$, the error matrix $\mathsf R$ of the steepest-descent
analysis has jumps $I+O(1/n)$ on two small circles around $0$ and $4c$ and exponentially small jumps elsewhere, and
\begin{equation}\label{eq:extract}
  n^{2p}b_n^2=(Q_1)_{12}(Q_1)_{21},\qquad Q_1=P_1+\mathsf R_1,
\end{equation}
where $P_1$ and $\mathsf R_1$ are the $1/z$ coefficients at infinity of the global parametrix and of $\mathsf R$
\cite[Lemmas~4.1, 4.2 and Proposition~5.6]{PaperII}. The subscript ``mod'' denotes quantities computed for the
\emph{model weight} $z^\alpha(4c-z)$, i.e.\ with $H_n$ replaced by a constant.

The band factor is $H_n=e^{-2G}$ with an explicit function $G$ \cite[Lemmas~8.1 and~10.1]{PaperII}. For $p>1$,
$G$ has a jump $i\fe$ on $(-\infty,0)$, with $|\fe(y)|\le\pi\exp(-2\pi\sin(\pi/p)n|y|^{-1/p})$, and
\begin{equation}\label{eq:G-pgt1}
  G(z)-G(1)=\frac1{2\pi}\int_{-\infty}^0\fe(y)\,\frac{z-1}{(y-z)(y-1)}\diff y\qquad(p>1)
\end{equation}
\cite[Lemma~8.1]{PaperII}. For $0<p<1$, $G$ jumps across the two seams $\arg z=\pm\vartheta$, $\vartheta=\pi p/2$, by
$\fj(|z|)$, where $\fj(r):=-\log(1-e^{-2\pi nr^{-1/p}})$, and
\begin{equation}\label{eq:G-plt1}
  G(z)-G(1)=\frac1{2\pi i}\int_0^\infty\fj(r)\bigl[e^{i\vartheta}\mathfrak K(re^{i\vartheta},z)
  -e^{-i\vartheta}\mathfrak K(re^{-i\vartheta},z)\bigr]\diff r,\qquad
  \mathfrak K(\zeta,z):=\frac{z-1}{(\zeta-z)(\zeta-1)}
\end{equation}
\cite[Lemma~10.1(ii)]{PaperII}. Expanding $\mathfrak K$ in powers of $z/\zeta$ gives the moments of the jump densities:

\begin{lemma}\label{lem:Fprime}
Let $p\ne1$ and $k\ge1$. The $k$-th moment of the jump density, $m_k:=\frac1{2\pi}\int_{-\infty}^0\fe(y)y^{-k-1}\diff y$
for $p>1$ and the corresponding seam integral for $p<1$, equals $\zeta(-kp)/(kn^{kp})$. For every $K\ge1$,
\[
  \log H_n(z)-\log H_n(1)=-\sum_{k\le K}\mu_{n,k}(z^k-1)+O(n^{-(K+1)p}),\qquad\mu_{n,k}:=\frac{2\zeta(-kp)}{kn^{kp}},
\]
uniformly for $|z|\le4c+2$, boundary values included, with one derivative.
\end{lemma}

For $p<1$ this is \cite[Lemma~10.1(iii)--(iv)]{PaperII}. For $p>1$ the proof is the same computation on the negative
axis: with $q(r)=\exp(-a_nr^{-1/p})$, $a_n=2\pi ine^{-i\pi/p}$, $\Re a_n=2\pi n\sin(\pi/p)>0$, one expands
$\fe=-2\Im\sum_mq^m/m$, integrates termwise by \cite[(5.9.1)]{DLMF}, sums by \cite[(25.2.1)]{DLMF} and applies the
reflection formula \cite[(25.4.1)]{DLMF}. At $p=2$ all $m_k$ vanish and $H_n$ is constant, in agreement with
\cite[Lemma~8.1]{PaperII}.

\subsection{Tilts and the gauge}\label{sec:gauge}
For $\mu=(\mu_1,\dots,\mu_K)\in\C^K$ put $q_\mu(x):=\exp(-\sum_k\mu_kx^k)$ and
\begin{equation}\label{eq:Qk}
  Q_k(z):=\sum_{j=0}^{k-1}\binom{2j}{j}c^jz^{k-1-j},\qquad M_\mu(z):=\sum_k\mu_kQ_k(z),
\end{equation}
so that $Q_k(4c)=c^{k-1}k\binom{2k}{k}/2$. The arcsine moments $\frac1\pi\int_0^{4c}x^j(x(4c-x))^{-1/2}\diff x=\binom{2j}{j}c^j$
give the Szeg\H{o} function of the tilt explicitly.

\begin{lemma}[gauge]\label{lem:gauge}
For the weight $z^\alpha(4c-z)q_\mu(z)$, the jump data on the two circles are $\Delta_k^\mu=G_\mu\Delta_k^{\mathrm{mod}}G_\mu^{-1}$,
where, with $h:=\mathcal RM_\mu/2$ and $S:=\sinh(h)/\mathcal R$,
\[
  G_\mu(z)=P_0^\infty e^{-\mathcal R M_\mu\sigma_3/2}(P_0^\infty)^{-1}
  =\mathsf T\begin{pmatrix}\cosh h&-(z-4c)S\\-zS&\cosh h\end{pmatrix}\mathsf T^{-1}
\]
is entire in $(z,\mu)$ with $\det G_\mu=1$. In particular $\Delta_k^\mu$ is meromorphic in each disc, with Laurent
coefficients entire in $\mu$.
\end{lemma}

By Lemma~\ref{lem:Fprime}, the band factor is the tilt $q_{\mu_n}$ with $\mu_n=(\mu_{n,1},\dots,\mu_{n,K})$, up to
$O(n^{-(K+1)p})$. The first-order effect of a tilt on $b_n$ is exactly linear:

\begin{lemma}[linear response]\label{lem:E37}
Let $e(\mu)$ be the $1/n$ coefficient of $\bigl((P_1+\mathsf R_1^{(1)}(\mu)/n)_{12}(P_1+\mathsf R_1^{(1)}(\mu)/n)_{21}\bigr)^{1/2}/c$
minus $d_1$, where $\mathsf R_1^{(1)}(\mu)=\Res_0\Delta_1^\mu+\Res_{4c}\Delta_1^\mu$. Then, for every $\mu\in\C^K$,
\[
  e(\mu)=cp\,M_\mu(4c)=cp\sum_k\mu_kQ_k(4c).
\]
\end{lemma}

The proof expands $G_\mu$ at $4c$ and at $0$. With $M_4:=M_\mu(4c)$, $G_\mu(0)$ commutes with the residue of
$\Delta_1^{\mathrm{mod}}$ at $0$, $M_\mu'(4c)$ drops out, and the residue at $4c$ moves by
\[
  M_4c^2p\begin{pmatrix}M_4c-\alpha-1&-i(M_4c-\alpha-2)\\-i(M_4c-\alpha)&-(M_4c-\alpha-1)\end{pmatrix},
\]
which is quadratic in $M_4$ but changes $i(A_{21}-A_{12})/(2c)$, the combination that enters the $1/n$ coefficient, by
exactly $cpM_4$. The identity $Q_k(4c)=c^{k-1}k\binom{2k}k/2$ follows from the recursion
$a_{k+1}=4a_k+\binom{2k}k$, $a_1=1$, satisfied by both $\sum_{j<k}4^{k-1-j}\binom{2j}j$ and $k\binom{2k}k/2$.

\subsection{The second order}\label{sec:E2E3}
\begin{theorem}\label{thm:E2E3}
Let $p>0$, $p\ne1$, $s>1$ and $K\ge1$. Then
\[
  \frac{n^pb_n}c=1+\frac{d_1}n+\frac{d_2}{n^2}+\sum_{j=1}^K\frac{\varepsilon_j}{n^{1+jp}}
  +O\bigl(n^{-3}+n^{-2-p}+n^{-1-(K+1)p}\bigr),
\]
\[
  d_2=d_2(p,s)=\frac{(p+1)(s-1)^2}{8p}-\frac{(p+2)(s-1-p)^2}{8p^2}+\frac{p+1}{12}.
\]
In particular, for $0<p<1$ with $1/p\notin\N$ the $n^{-2}$ coefficient is $d_2$; for $p=1/j_0$, $j_0\ge2$, it is
$d_2+\varepsilon_{j_0}=d_2-p/12$, because $W(j_0)=\binom{2j_0}{j_0}c(1/j_0)^{j_0}=1$.
\end{theorem}

For $p>1$ the terms with $j\ge2$ are $O(n^{-3})$, and the statement is $n^pb_n/c=1+d_1/n+d_2/n^2+\kp n^{-1-p}+O(n^{-3})$. For
$p=2$, $d_2=(s^2+6s-11)/16$ and $\kp=0$. The proof of Theorem~\ref{thm:E2E3} is the second Neumann iterate of the small-norm
problem, with the second asymptotic coefficients of the Bessel and Airy functions on the two circles and the band
factor entering through $G_{\mu_n}$; the $1/n^2$ term at $\mu=0$ is an exact Laurent computation, and the zeta terms are
Lemma~\ref{lem:E37} with $\mu_{n,j}\cdot cpQ_j(4c)=\varepsilon_jn^{-jp}$. Section~\ref{sec:natural} carries this out at
every order.

% Section 3: The constrained equilibrium (Theorem C).
\section{The constrained equilibrium}\label{sec:equilibrium}

At leading order the zeros of $P_n$, rescaled by $n^p$, follow the equilibrium measure of unit mass constrained by
the continuum density of the nodes. For the lattice this problem was solved in \cite[Theorem~1.3]{PaperI} in the index
variable. Here we solve it for a general constraint density in the edge variable, by a Cauchy-transform argument that
gives the density explicitly, and we show that the soft edge of the discrete analogue carries the poles of the node
zeta function with the weight $\pi\cot(\pi q)$. Nothing in this section concerns $b_n$ directly; the link is discussed
in Section~\ref{sec:conj}.

\subsection{Theorem C}\label{sec:thmQ}
Let $0<Y\le\infty$ and let $\sigma$ be a density on $(0,Y]$ such that
\begin{enumerate}
\item[(H1)] $\sigma$ is positive and locally H\"older continuous on $(0,Y)$ (the H\"older condition is used only in~(iii));
\item[(H2)] $y\mapsto y\sigma(y)$ is strictly decreasing on $(0,Y)$;
\item[(H3)] if $Y=\infty$, then $\int_1^\infty\sigma(y)\log y\diff y<\infty$.
\end{enumerate}
By (H2), $g(y):=\sqrt y\,\sigma(y)=y\sigma(y)/\sqrt y$ is strictly decreasing, and $\sigma(y)\ge y_0\sigma(y_0)/y$ for
$y<y_0$, so $\int_0\sigma=\infty$; if $Y=\infty$, (H3) gives $\int^\infty\sigma<\infty$ and hence $y\sigma(y)\to0$.
For a positive measure $\mu$ on $[0,Y]$ let $U^\mu(z):=\int\log\frac1{|z-x|}\diff\mu(x)$ and
$I(\mu):=\int U^\mu\diff\mu$. The \emph{constrained equilibrium problem} is to minimise $I(\mu)$ over positive $\mu$ on
$[0,Y]$ with $\mu([0,Y])=1$ and $\mu\le\sigma$. The \emph{soft-edge condition} is
\begin{equation}\label{eq:Q}
  F(a):=\int_a^Y\sigma(y)\Bigl(1-\frac ay\Bigr)^{-1/2}\diff y=1 .
\end{equation}

\begin{theorem}\label{thm:Q}
Assume \textup{(H1)--(H3)}.
\begin{enumerate}
\item[(i)] Equation~\eqref{eq:Q} has exactly one root $a\in(0,Y)$.
\item[(ii)] Put $s(y):=\sigma(y)\sqrt{y/(y-a)}$ on $(a,Y)$ and
\begin{equation}\label{eq:nu}
  \nu(x):=\frac1\pi\sqrt{\frac{a-x}x}\int_a^Y\frac{s(y)}{y-x}\diff y\qquad(0<x<a).
\end{equation}
Then $\mu:=\nu\diff x|_{[0,a]}+\sigma\diff y|_{[a,Y]}$ has unit mass, $0<\nu<\sigma$ on $(0,a)$, and
$\pi\sqrt{x(a-x)}\,\nu(x)=1-\int_a^Y\sigma(y)\sqrt{y(y-a)}\,(y-x)^{-1}\diff y$.
\item[(iii)] There is a constant $\ell$ with $U^\mu=\ell$ on $[0,a]$ and $U^\mu<\ell$ on $(a,Y]$.
\item[(iv)] $\mu$ is the unique solution of the constrained equilibrium problem.
\end{enumerate}
\end{theorem}

\begin{proof}
(i) Substituting $y=au$, $F(a)=\int_1^{Y/a}a\sigma(au)(1-1/u)^{-1/2}\diff u$. For fixed $u$, $a\sigma(au)=(au\,\sigma(au))/u$
is strictly decreasing in $a$ by (H2), and the upper limit decreases, so $F$ is strictly decreasing. On $[a_0,a_1]\subset(0,Y)$
we have $a\sigma(au)\le a_0\sigma(a_0u)$ by (H2), so $F$ is continuous by dominated convergence. As $a\to0$,
$F(a)\ge\int_a^Y\sigma\to\infty$. If $Y<\infty$, $F(a)\to0$ as $a\to Y$. If $Y=\infty$, split at $2a$: $\int_{2a}^\infty\le
\sqrt2\int_{2a}^\infty\sigma\to0$, and, since $\sigma(y)\le a\sigma(a)/y$ on $[a,2a]$, $\int_a^{2a}\le
a\sigma(a)\int_1^2(1-1/u)^{-1/2}\diff u\to0$.

(ii) Let $S(z):=\int_a^Ys(y)(z-y)^{-1}\diff y$ and $\tilde G(z):=((z-a)/z)^{1/2}S(z)$ (principal branch), analytic off
$[0,Y]$, and let $G(z):=\int(z-x)^{-1}\diff\mu(x)$. For fixed $y\in(a,Y)$ put
\[
  k_y(z):=\frac{((z-a)/z)^{1/2}-\sqrt{(y-a)/y}}{z-y},\qquad \rho_y(x):=\frac1\pi\sqrt{\frac{a-x}x}\,\frac1{y-x}>0\quad(0<x<a).
\]
The singularity of $k_y$ at $z=y$ is removable, so $k_y$ is analytic off $[0,a]$; it is $O(|z|^{-1/2})$ at $0$, bounded
at $a$, and $\sim(1-\sqrt{(y-a)/y})/z$ at infinity; since $((z-a)/z)^{1/2}$ takes the values $\pm i\sqrt{(a-x)/x}$ on
$(0,a)$ from above and below, $k_y$ jumps by $-2\pi i\rho_y$ there. As $[0,a]$ is compact, Liouville's theorem gives
$k_y(z)=\int_0^a\rho_y(x)(z-x)^{-1}\diff x$, and $\rho_y$ has mass $1-\sqrt{(y-a)/y}$. Since $\sigma=s\sqrt{(y-a)/y}$,
\[
  \tilde G(z)=\int_a^Ys(y)\Bigl[\frac{\sqrt{(y-a)/y}}{z-y}+k_y(z)\Bigr]\diff y
  =\int_a^Y\frac{\sigma(y)\diff y}{z-y}+\int_0^a\frac{\nu(x)\diff x}{z-x}=G(z),
\]
by Fubini's theorem (everything is positive, $|z-x|^{-1}\le|\Im z|^{-1}$, and $\int s(1-\sqrt{(y-a)/y})\diff y\le F(a)$),
because $\int s(y)\rho_y(x)\diff y=\nu(x)$. Hence $\mu$ has mass $\int_a^Y\sigma+\int s(1-\sqrt{(y-a)/y})=\int s=F(a)=1$;
$\Re G(x\pm i0)=0$ on $(0,a)$, where $S$ is real-analytic; and $\nu>0$ because $s>0$. With $R(y):=\sqrt{y(y-a)}$ one has $R\sigma=s\,(y-a)$, which gives
$\tilde G=R(z)^{-1}\bigl[F(a)+\int\sigma R\,(z-y)^{-1}\diff y\bigr]$; the difference of the two expressions for $\nu$ is
$(F(a)-1)/(\pi\sqrt{x(a-x)})$, which vanishes by~\eqref{eq:Q}. For the upper bound write $s(y)=g(y)(y-a)^{-1/2}$ and use
$\int_a^\infty(y-a)^{-1/2}(y-x)^{-1}\diff y=\pi(a-x)^{-1/2}$:
\[
  \nu(x)<\frac1\pi\sqrt{\frac{a-x}x}\,g(a)\,\frac\pi{\sqrt{a-x}}=\frac{g(a)}{\sqrt x}\le\frac{g(x)}{\sqrt x}=\sigma(x)
  \qquad(0<x<a).
\]

(iii) On $(0,a)$, $U^\mu$ is constant because $\Re G=0$ there. For $z\in(a,Y)$, $\varphi:=U^\mu-\ell$ satisfies
$\varphi(a)=0$ and $\varphi'(z)=-\Re G(z\pm i0)=-((z-a)/z)^{1/2}H(z)$ with $H(z):=\PV\!\int_a^Ys(y)(z-y)^{-1}\diff y$, which
exists since $\sigma$ is locally H\"older. Writing $s=g\cdot(y-a)^{-1/2}$,
\[
  H(z)=\int_a^Y\frac{(g(y)-g(z))(y-a)^{-1/2}}{z-y}\diff y+g(z)\,\PV\!\!\int_a^Y\frac{(y-a)^{-1/2}}{z-y}\diff y .
\]
The first integrand is positive since $g$ is decreasing. The principal value is $0$ if $Y=\infty$, and
$(z-a)^{-1/2}\log\frac{\sqrt{Y-a}+\sqrt{z-a}}{\sqrt{Y-a}-\sqrt{z-a}}>0$ if $Y<\infty$. Hence $H>0$, $\varphi'<0$ and
$\varphi<0$ on $(a,Y]$.

(iv) Since $\log^+|x-y|\le\log(1+x)+\log(1+y)$ and (H3) bounds the logarithmic moment of every $\mu'\le\sigma$,
$I(\mu')>-\infty$ for every competitor; let $\mu'\le\sigma$ have unit mass and $I(\mu')<\infty$, and $\tau:=\mu'-\mu$. Then
$I(\mu')-I(\mu)=I(\tau)+2\int(U^\mu-\ell)\diff\tau$. The integrand vanishes on $[0,a]$ (at every point, since the density
of $\mu$ is at most $g(a)x^{-1/2}$ on $(0,a)$ and bounded on $[a,Y]$, so $U^\mu$ is continuous); on $(a,Y]$, $\tau\le0$ and
$U^\mu-\ell<0$, so the integral is $\ge0$. For $I(\tau)$: $\tau$ is a finite signed measure with $\tau(\R)=0$, and
$\iint|\log|x-y||\diff|\tau|\diff|\tau|<\infty$ (the positive part of the logarithm by (H3), the negative part because $\mu$
has density at most $g(a)\max(x^{-1/2},a^{-1/2})$ and $I(\mu')<\infty$). So $I(\tau)\ge0$, with equality only if $\tau=0$,
by \cite[Lemma~3.3]{PaperI}.
\end{proof}

\begin{remark}\label{rem:Q}
(a) The proof of (ii)--(iii) uses only that $g$ is decreasing; (H2) enters through the monotonicity of $F$ in (i). So
(H2) may be replaced by ``$\sqrt y\,\sigma(y)$ strictly decreasing and $\int_0^Y\sigma>1$''; uniqueness of the root then
follows from (iv), since $a$ is the left end of the saturated set of the unique minimiser. Neither condition is necessary,
but they cannot simply be dropped: in a discretised energy minimisation, a constraint density with a dip below the edge gives
three roots of~\eqref{eq:Q} and a minimiser with two bands, and a shallower dip gives one root but a saturated interval
inside the band.
(b) Without~\eqref{eq:Q}, $\nu$ acquires a multiple of $(a-x)^{-1/2}$, which violates $\nu\ge0$ or $\nu\le\sigma$. Under
\eqref{eq:Q}, $\nu(a^-)=g(a)/\sqrt a=\sigma(a)$, so $\mu$ is continuous at $a$. If $\sigma$ is H\"older of order $>\frac12$
near $a$ (for instance $C^1$), then $\sigma-\nu=C\sqrt{a-x}\,(1+o(1))$ with $C=-h(a)/(\pi\sqrt a)>0$, where
$h(a):=\int_a^Y(g(y)-g(a))(y-a)^{-3/2}\diff y-2g(a)(Y-a)^{-1/2}\mathbf 1_{Y<\infty}$.
(c) In node language~\eqref{eq:Q} is the statement that $\frac1n\sum_{u_k<1}(1-u_k)^{-1/2}\approx1$, with
$u_k=a\lambda_k/n^p$: the edge is fixed by a sum of the per-node quantity $(1-u)^{-1/2}$ over the saturated nodes.
\end{remark}

\begin{corollary}[the lattice]\label{cor:Q1}
For $\sigma(y)=\frac1py^{-1/p-1}$ on $(0,\infty)$ the root of~\eqref{eq:Q} is $a=4c(p)$.
\end{corollary}

\begin{proof}
$y\sigma(y)=\frac1py^{-1/p}$ is strictly decreasing and (H3) holds. Equation~\eqref{eq:Q} reads
$a^{-1/p}\frac1pB(\frac1p,\frac12)=1$, i.e.\ $a^{1/p}=\sqrt\pi\,\Gamma(\frac1p+1)/\Gamma(\frac1p+\frac12)=\pi/(2I(p))$, and
$a=4c(p)$ by~\eqref{eq:cp}.
\end{proof}

The density $\sigma(y)=\frac1py^{-1/p-1}$ is the image of Lebesgue measure on the index variable $u$ under
$y=u^{-p}$, so Corollary~\ref{cor:Q1} concerns the problem of \cite[Theorem~1.3]{PaperI}; its edge $a=4c$ is the
right end of the interval on which the potential of the measure $\mu^*$ there is constant
\cite[Theorem~1.3(ii)]{PaperI}. The derivation above does not use the Riemann--Hilbert analysis of \cite{PaperII}.

\subsection{The soft edge carries the poles}\label{sec:M12}
Let $L(u):=(1-u)^{-1/2}\mathbf 1_{u<1}$; its Mellin transform $\Lt(w):=\int_0^\infty u^{-w-1}L(u)\diff u$ is computed in
Lemma~\ref{lem:K0}: $\Lt(w)=B(-w,\frac12)=-\pi\cot(\pi w)\,\Gamma(w+\frac12)/(\sqrt\pi\,\Gamma(w+1))$, so that
\begin{equation}\label{eq:Lt-W}
  -(4c)^q\,\Lt(q)=W(q)\,\pi\cot(\pi q).
\end{equation}
At $q\in-\frac12-\N_0$ the right side is a limit of the form $0\cdot\infty$ and is defined by the left side; for
instance $\Lt(-\frac12)=\pi$.

Let $\mathcal N$ be a $C^1$ counting function on $[\lambda_0,\infty)$, $\lambda_0>0$, with $\mathcal N'>0$, and suppose
that for exponents $\rho_0=1/p>\rho_1>\rho_2>\dotsb$ and constants $c_\rho$ with $c_{1/p}=1$,
\begin{equation}\label{eq:Nexp}
  \mathcal N'(\lambda)=\sum_{k<K}\rho_kc_{\rho_k}\lambda^{\rho_k-1}+O(\lambda^{\rho_K-1})\qquad(\lambda\to\infty)
\end{equation}
for every $K$ (if the list is finite, the remainder is $O(\lambda^{-M})$ for every $M$). Then
$Z_{\mathcal N}(s):=\int_{\lambda_0}^\infty\lambda^{-s}\diff\mathcal N(\lambda)$ continues meromorphically to $\C$, with
simple poles at $s=\rho$ of residue $\rho c_\rho$. In the variable $w=-s$, $w\mapsto Z_{\mathcal N}(-w)$ has poles at
$q=-\rho$ with residues $R_q=-\rho c_\rho$; the \emph{Weyl pole} is $q=-1/p$, $R_{-1/p}=-1/p$. We assume that
$\rho c_\rho=0$ whenever $\rho\in\{-1,-2,\dots\}$, i.e.\ no pole of $Z_{\mathcal N}(-w)$ sits at a positive integer.

\begin{proposition}\label{prop:M12}
Under these assumptions, as $A\to0^+$,
\begin{equation}\label{eq:M12a}
  \int_{\lambda_0}^\infty L(A\lambda)\diff\mathcal N(\lambda)\sim-\sum_qR_q\Lt(q)A^q
  +\sum_{j\ge0}\binom{2j}j4^{-j}Z_{\mathcal N}(-j)A^j,
\end{equation}
with remainder $O(A^{-\rho_K})$ after the terms with $q<-\rho_K$ and $j<-\rho_K$. Consequently the root of
$\int L(A\lambda)\diff\mathcal N(\lambda)=n$, written $A=4c(1+\eta)n^{-p}$, satisfies
\begin{equation}\label{eq:M12b}
  \eta=\frac pn\Bigl[\sum_{j\ge0}W(j)Z_{\mathcal N}(-j)n^{-jp}+\sum_{q\ne-1/p}R_qW(q)\,\pi\cot(\pi q)\,n^{-pq}\Bigr]
  +O(\eta^2),
\end{equation}
and the higher orders follow by reversion of the series.
\end{proposition}

\begin{proof}
For one power, $\diff\mathcal N=\rho c_\rho\lambda^{\rho-1}\diff\lambda$, the substitution $u=A\lambda$ and the Taylor
expansion of $(1-u)^{-1/2}$ on $[0,A\lambda_0]$ give
\[
  \int_{\lambda_0}^{1/A}(1-A\lambda)^{-1/2}\rho c_\rho\lambda^{\rho-1}\diff\lambda
  =\rho c_\rho\Bigl[A^{-\rho}B(\rho,\tfrac12)-\sum_{j\ge0}\binom{2j}j4^{-j}A^j\frac{\lambda_0^{\rho+j}}{\rho+j}\Bigr],
\]
first for $\rho>0$ and then for all $\rho\notin-\N$ by analytic continuation in $\rho$. The first term is
$-R_q\Lt(q)A^q$ with $q=-\rho$. For the remainder $r_K$ of~\eqref{eq:Nexp}, split at $\lambda=1/(2A)$: on
$[\lambda_0,1/(2A)]$ expand $(1-A\lambda)^{-1/2}$ in powers of $A\lambda$, and on $[1/(2A),1/A]$ use
$|r_K|\le C\lambda^{\rho_K-1}$; both give the integer terms up to $O(A^{-\rho_K})$. The coefficient of $A^j$ collects
$-\rho c_\rho\lambda_0^{\rho+j}/(\rho+j)$ over the powers and $\int\lambda^jr_K\diff\lambda$, i.e.\ the continued value
$Z_{\mathcal N}(-j)$. This proves~\eqref{eq:M12a}. The Weyl pole contributes $\frac1pB(\frac1p,\frac12)A^{-1/p}=n(1+\eta)^{-1/p}$
by Corollary~\ref{cor:Q1}, so the equation is $n(1+\eta)^{-1/p}+(\text{rest})=n$; by~\eqref{eq:Lt-W} and
$\binom{2j}j4^{-j}(4c)^j=W(j)$, the rest at $\eta=0$ is the bracket of~\eqref{eq:M12b}, and the implicit function
theorem gives~\eqref{eq:M12b}.
\end{proof}

The termwise proof avoids shifting the contour of the Mellin--Barnes form
\[
  \frac1{2\pi i}\int\Lt(w)Z_{\mathcal N}(-w)A^w\diff w
\]
of~\eqref{eq:M12a}, which would need bounds on $Z_{\mathcal N}$ on vertical lines. The numbers $Z_{\mathcal N}(-j)$ depend on the cutoff $\lambda_0$; they are not the values $Z_\Lambda(-j)$
of a node set, which enter through the band factor (Section~\ref{sec:kernel}). What~\eqref{eq:M12b} contributes for a
node set is the pole part: the weight $W(q)\pi\cot(\pi q)$ of each pole of $Z(-w)$ other than the Weyl pole.

\begin{remark}[exact root for a drum-type count]\label{rem:drum-root}
For $p=1$ ($c=\frac12$) and $\diff\mathcal N=(1-\frac\beta2\lambda^{-1/2})\diff\lambda$ on $(0,\infty)$, the condition is
$2/A-\frac{\pi\beta}2A^{-1/2}=n$. With $A=2(1+\eta)/n$ and $v:=(1+\eta)^{-1/2}$ it reads $v^2-\epsilon v-1=0$,
$\epsilon:=\pi\beta/(2\sqrt{2n})$, so $1+\eta=v^{-2}=1-\epsilon+\epsilon^2/2+O(\epsilon^3)$ and
\[
  n\eta=-\frac{\pi\beta}{2\sqrt2}\,n^{1/2}+\frac{\pi^2}{16}\beta^2+O(n^{-1/2}).
\]
The first term is the pole term of~\eqref{eq:M12b} at $q=-\frac12$, $R_q=\beta/2$, with $W(q)\pi\cot(\pi q)=-\pi(4c)^{-1/2}$;
the second is quadratic in the pole residue and is produced by the reversion. Both are compared with $b_n$ in
Section~\ref{sec:drums}.
\end{remark}

\begin{remark}[numerical checks]\label{rem:Q-checks}
For the lattice at $p=0.6$ and $p=1.5$, for the drum density of Remark~\ref{rem:drum-root} cut at the first node, and for a
fourth density, the root, the two forms of $\nu$, the mass, $0<\nu/\sigma<1$ and the outer inequality of
Theorem~\ref{thm:Q} were evaluated by quadrature, the outer inequality both through $H$ and by direct evaluation of
$U^\mu$, and $-((z-a)/z)^{1/2}H(z)$ was compared with a difference quotient of $U^\mu$; the root agrees with $4c$ to $10^{-12}$ in the lattice cases, and the
two forms of $\nu$ to $10^{-21}$, which tests the root and the quadrature rather than~(ii). The first-order pole term
of~\eqref{eq:M12b} was fitted from the root over $n=100$--$6400$ for $\lambda_k=(k+\gamma k^{-m})^p$ at
$(p,m)=(1.5,\frac12),(0.75,\frac14),(1.25,\frac12),(1.8,\frac34)$; it agrees with $pR_qW(q)\pi\cot(\pi q)$ to between
$10^{-8}$ (the printed digits) and $10^{-6}$ (scripts \texttt{m1\_theoremQ}, \texttt{m1\_edge}).
\end{remark}

% Section 4: The kernel (Lemma K0, Theorem B).
\section{The kernel}\label{sec:kernel}

By Lemma~\ref{lem:E37} and Lemma~\ref{lem:Fprime}, the zeta terms of Theorem~\ref{thm:E2E3} are the values at
$\mu=\mu_n$ of the linear functional $cpM_\mu(4c)$, truncated to $K$ terms. This section evaluates the same functional on
the band factor itself, without truncation, and writes it as one Mellin--Barnes integral (Theorem~\ref{thm:K}). Its
kernel contains $-\pi\cot(\pi w)$, which is also the Mellin transform of the truncated per-node law of
Section~\ref{sec:equilibrium} (Lemma~\ref{lem:K0}). The proofs are for the lattice; Section~\ref{sec:conj} discusses
other node sets.

\subsection{The truncated per-node law}\label{sec:K0}
\begin{lemma}\label{lem:K0}
Let $L(u):=(1-u)^{-1/2}\mathbf 1_{u<1}$. For $\Re w<0$,
\[
  \Lt(w):=\int_0^\infty u^{-w-1}L(u)\diff u=B\bigl(-w,\tfrac12\bigr)
  =-\pi\cot(\pi w)\,\frac{\Gamma(w+\frac12)}{\sqrt\pi\,\Gamma(w+1)} .
\]
$\Lt$ continues meromorphically to $\C$; its poles in $\Re w\ge0$ are simple, at $w=j\in\N_0$, with
$-\Res_{w=j}\Lt=\binom{2j}j4^{-j}$, the Taylor coefficients of $L$ at $0$. For $a=4c$,
$(a/n^p)^w\Gamma(w+\frac12)/(\sqrt\pi\Gamma(w+1))=W(w)n^{-pw}$.
\end{lemma}

\begin{proof}
The integral is $B(-w,\frac12)=\Gamma(-w)\sqrt\pi/\Gamma(\frac12-w)$ by \cite[(5.12.1)]{DLMF}. By the reflection formula
\cite[(5.5.3)]{DLMF}, $\Gamma(\frac12-w)\Gamma(\frac12+w)=\pi/\cos(\pi w)$ and $\Gamma(-w)\Gamma(1+w)=-\pi/\sin(\pi w)$, so
$\sqrt\pi\,\Gamma(-w)/\Gamma(\frac12-w)=\Gamma(-w)\Gamma(w+\frac12)\cos(\pi w)/\sqrt\pi=-\pi\cot(\pi w)\Gamma(w+\frac12)/(\sqrt\pi\,\Gamma(w+1))$.
The residue of $-\pi\cot(\pi w)$ at $j$ is $-1$, and $\Gamma(j+\frac12)/(\sqrt\pi\,j!)=\binom{2j}j4^{-j}$.
\end{proof}

The two analytic continuations of $(1-u)^{-1/2}$ to $u>1$ are $\mp i|1-u|^{-1/2}$, and their Mellin transforms carry
$\Gamma(-w)\Gamma(1+w)e^{\pm i\pi w}=-\pi e^{\pm i\pi w}/\sin(\pi w)$. Their average is $L$, and the average of the two
transforms is $-\pi\cot(\pi w)$: the cotangent is the signature of truncating the per-node law at the edge $u=1$.

\subsection{The band-factor functional}\label{sec:en}
For a function $f$ analytic near the jump contour of $G$ put, with the notation of~\eqref{eq:G-pgt1}--\eqref{eq:G-plt1},
\begin{gather*}
  \mathcal M_n[f]:=\frac1{2\pi}\int_{-\infty}^0\fe(y)f(y)\diff y\qquad(p>1),\\
  \mathcal M_n[f]:=\frac1\pi\int_0^\infty\fj(r)\,\Im\bigl[e^{i\vartheta}f(re^{i\vartheta})\bigr]\diff r\qquad(0<p<1),
\end{gather*}
so that $\mathcal M_n[y^{-k-1}]=m_k=\zeta(-kp)/(kn^{kp})$ by Lemma~\ref{lem:Fprime} (for $0<p<1$ the integrand is
real and $\mathcal M_n$ is the seam integral of \cite[Lemma~10.1(iii)]{PaperII}). Since
$\sum_{j\ge0}\binom{2j}jc^jt^j=(1-4ct)^{-1/2}$ and $\sum_l(4c)^lt^l=(1-4ct)^{-1}$,
\begin{equation}\label{eq:genQ}
  \sum_{k\ge1}Q_k(4c)\,y^{-k-1}=k_c(y):=y^{-2}\Bigl(1-\frac{4c}y\Bigr)^{-3/2}\qquad(|y|>4c).
\end{equation}
The \emph{band-factor functional} is
\begin{equation}\label{eq:en}
  e_n:=2cp\,\mathcal M_n[k_c],
\end{equation}
that is, $e_n=\frac{cp}\pi\int_0^\infty\fe(-r)\,r^{-1/2}(r+4c)^{-3/2}\diff r$ for $p>1$ and
$e_n=\frac{2cp}\pi\int_0^\infty\fj(r)\Im[e^{i\vartheta}k_c(re^{i\vartheta})]\diff r$ for $0<p<1$. On the seams
$\Im(1-4c/\zeta)=(4c/r)\sin\vartheta>0$, so the principal branch in $k_c$ is never crossed. By Lemma~\ref{lem:E37} and
Lemma~\ref{lem:Fprime}, $e(\mu_n)=cp\sum_{k\le K}\mu_{n,k}Q_k(4c)=2cp\sum_{k\le K}m_kQ_k(4c)$ is the truncation of $e_n$ obtained
by expanding $k_c$ in~\eqref{eq:genQ}.

\begin{theorem}\label{thm:K}
Let $p>0$, $p\ne1$, $s>1$, and $0<c_0<1$.
\begin{enumerate}
\item[(i)] $e_n=2cp\sum_{k\le K}m_kQ_k(4c)+O(n^{-(K+1)p})$ for every $K$.
\item[(ii)] $\displaystyle e_n=\frac1{2\pi i}\int_{c_0-i\infty}^{c_0+i\infty}p\,W(w)\,\zeta(-wp)\,n^{-wp}\,
  \pi\Bigl[\cot\frac{\pi wp}2-\cot(\pi w)\Bigr]\diff w$.
\item[(iii)] For every $J\ge1$, $e_n=\sum_{j\le J}\varepsilon_jn^{-jp}+O(n^{-(J+1)p})$.
\end{enumerate}
\end{theorem}

Part~(iii) follows from~(i) and $2cpm_jQ_j(4c)=\varepsilon_jn^{-jp}$; the content of the theorem is the exact
representation~(ii) and its kernel.

\begin{proof}
(i) Split the integral at $|y|=2(4c+2)$. On the far part, expand $k_c$ by~\eqref{eq:genQ} to $K$ terms with remainder
$O(|y|^{-K-2})$; the near part, and the near parts of the moments, are $O(e^{-n/C})$ because the jump densities are
\cite[Lemmas~8.1(i) and~10.1]{PaperII}. The remainder is bounded as in Lemma~\ref{lem:Fprime}.

(ii), $p>1$. For complex $w$ with $\Re w>0$, the computation of Lemma~\ref{lem:Fprime} with $k$ replaced by $w$ (the
Fubini bound depends only on $\Re w$) gives
\[
  I(w):=\int_0^\infty\fe(-r)r^{-w-1}\diff r=-2p\,\Gamma(wp)\zeta(1+wp)(2\pi n)^{-wp}\sin\Bigl(\pi w-\frac{\pi wp}2\Bigr).
\]
The strip is set by $\fe(-r)\to\pi(1-2/p)$ as $r\to\infty$ and by the exponential decay of $\fe(-r)$ as $r\to0$. By
\cite[(25.4.1)]{DLMF} at $1+wp$, $p\Gamma(wp)\zeta(1+wp)(2\pi)^{-wp}=-\pi\zeta(-wp)/(w\sin(\pi wp/2))$, so
\[
  I(w)=\frac{2\pi}w\,\zeta(-wp)\,n^{-wp}\Bigl[\sin(\pi w)\cot\frac{\pi wp}2-\cos(\pi w)\Bigr].
\]
Next, $r^{-1/2}(r+4c)^{-3/2}=r^{-2}(1+4c/r)^{-3/2}$, and the Mellin--Barnes representation
\begin{equation}\label{eq:MB32}
  (1+x)^{-3/2}=\frac1{2\pi i}\int_{(c_1)}\frac{\Gamma(-s)\Gamma(\frac32+s)}{\Gamma(\frac32)}\,x^s\diff s\qquad
  (x>0,\ -\tfrac32<c_1<0)
\end{equation}
is the Mellin inversion of the Beta integral \cite[(5.12.3)]{DLMF}, whose transform is absolutely integrable on the line.
Put $w=s+1$, so that $\Re w=c_0\in(0,1)$, and exchange the integrals (absolute convergence). Using
$\Gamma(1-w)/w=\pi/(\sin(\pi w)\Gamma(1+w))$ and $\Gamma(\frac32)=\sqrt\pi/2$, the factors collect into
$pW(w)\zeta(-wp)n^{-wp}\pi[\sin(\pi w)\cot(\pi wp/2)-\cos(\pi w)]/\sin(\pi w)$, which is the stated integrand.

(ii), $0<p<1$. The Mellin transform of the seam density is $\int_0^\infty\fj(r)r^{-w-1}\diff r=p(2\pi n)^{-wp}\Gamma(wp)\zeta(1+wp)$
for $\Re w>0$, by the computation of \cite[Lemma~10.1(iii)]{PaperII} with complex $w$; the strip comes from
$\fj(r)\sim\frac1p\log r$ at infinity and the super-exponential decay at $0$. Expand $(1+v)^{-3/2}$ with
$v=-4ce^{-i\vartheta}/r$ by~\eqref{eq:MB32} (now $\arg v=\pi-\vartheta\in(0,\pi)$, where the representation continues). With
$w=s+1$ the angular factor is $\Im[e^{i\vartheta}e^{-2i\vartheta}e^{-i\vartheta s}e^{i\pi s}]=-\sin((\pi-\vartheta)w)$, and with
$\vartheta=\pi p/2$ and the same two identities the integrand becomes
$pW(w)\zeta(-wp)n^{-wp}\pi\sin(\pi w-\frac{\pi wp}2)/(\sin(\pi w)\sin\frac{\pi wp}2)$. Finally
$\sin(\pi w-\frac{\pi wp}2)/(\sin(\pi w)\sin\frac{\pi wp}2)=\cot\frac{\pi wp}2-\cot(\pi w)$.

(iii) On vertical lines $|\cot(\pi wp/2)-\cot(\pi w)|\le Ce^{-\min(p,2)\pi|\Im w|}$, while $\zeta(-wp)$ and $W(w)$ grow at
most polynomially, so the contour can be moved to the right. The poles of $-\pi\cot(\pi w)$ at $w=j\ge1$ give
$-\Res=pW(j)\zeta(-jp)n^{-jp}=\varepsilon_jn^{-jp}$. The poles of $\cot(\pi wp/2)$ at $w=2l/p$ lie on the zeros
$\zeta(-2l)=0$ and are removable; where $2l/p=j$ is an integer the product is regular and $\varepsilon_j=0$.
\end{proof}

\begin{remark}[reading of Theorem~\ref{thm:K}]\label{rem:K}
(a) The second kernel $\pi\cot(\pi wp/2)$ comes only from writing $I(w)$ with $\zeta(-wp)$: in the form
$-2p\Gamma(wp)\zeta(1+wp)(2\pi n)^{-wp}\sin(\pi w-\pi wp/2)$ there are no poles at $w=2l/p$. So the band factor contributes
exactly the integer residues of $-\pi\cot(\pi w)$.
(b) The pole of $\zeta(-wp)$ at $w=-1/p$ (the Weyl pole) and the point $w=0$ lie to the left of the strip. The $j=0$ term
$pW(0)Z(0)=p\zeta(0)$, which is part of $d_1$, does not come from $e_n$ but from the index normalisation of the
frame.
(c) $e_n$ is a smooth function of $n$ with an expansion in powers of $n^{-p}$. The literal sum
$\frac pn\sum_{u_m<1}(1-u_m)^{-1/2}$, $u_m=4cm^p/n^p$, differs from its Mellin pole part by an oscillation of order
$n^{-1/2}$ set by the fractional part of $(n^p/4c)^{1/p}$, which is not present in $b_n$. Theorem~\ref{thm:K} says that
the band factor produces the pole part only.
(d) The negative-axis geometry ($p>1$) and the seam geometry ($p<1$) give the same kernel. At $p=1$ the two cotangents
coincide and their poles merge, which is where the expansions of \cite{PaperII} change form.
(e) The residue sum in~(iii) is asymptotic only: $\zeta(-jp)$ grows factorially in $j$.
\end{remark}

Numerically, the quadrature of $e_n$ and the Mellin--Barnes integral agree to all $16$ printed digits at
$(p,n)=(1.5,10),(1.25,20),(3,10),(0.6,10),(0.3,8),(0.75,20)$, and the partial residue sums approach them (scripts
\texttt{spec\_kernel}, \texttt{m2\_thmK\_plt1}).

\subsection{Perturbed index functions}\label{sec:perturbed}
For a node set other than the lattice, the analysis of \cite{PaperII} has not been carried out. The following
computation shows what the band factor would contribute if the construction carried over with the lattice's index
function $t(z)=nz^{-1/p}$ replaced by $t(z)=\kappa(n^p/z)$, $\kappa$ the inverse of $k\mapsto\lambda_k$.

\begin{proposition}\label{prop:pert}
Let $p>1$, $m>0$, and for $\lambda_k=(k+\gamma k^{-m})^p$ let $\delta t_+(r):=-\gamma n^{-m}(re^{i\pi})^{m/p}$ be the
first-order change of $t_+=nr^{-1/p}e^{-i\pi/p}$ on the negative axis, and
$\delta\fe:=2\Im[\mathcal F'(t_+)\delta t_+]$ with $\mathcal F'(t)=2\pi i/(e^{2\pi it}-1)$. Then, for $\Re w>(m+1)/p$,
\[
  \int_0^\infty\delta\fe(-r)r^{-w-1}\diff r=-2\pi\gamma p\,n^{-pw}\zeta(1+m-pw)
  \Bigl[\cos(\pi w)+\sin(\pi w)\tan\frac{\pi(pw-m)}2\Bigr].
\]
\end{proposition}

The proof is the computation of Lemma~\ref{lem:Fprime}: expand $\mathcal F'$ in powers of $e^{-2\pi it_+}$, integrate termwise
by \cite[(5.9.1)]{DLMF}, sum $\sum_{m'}m'^{\,m-pw}=\zeta(pw-m)$ by \cite[(25.2.1)]{DLMF} and apply \cite[(25.4.1)]{DLMF}; a
quadrature check agrees to $10^{-13}$ at $w=2.3$ and $3.1$ ($p=1.5$, $m=1$). With $\delta Z(-w)=pw\gamma\zeta(1+m-pw)$, the
first-order change of $Z_\Lambda(-w)$, the change of $e_n$ is
$\frac1{2\pi i}\int pW(w)\delta Z(-w)n^{-pw}\pi[-\cot(\pi w)-\tan(\pi(pw-m)/2)]\diff w$ on a line $\Re w>(m+1)/p$.

\begin{remark}\label{rem:pert}
Under that hypothetical construction, the integer residues give $pW(j)\delta Z(-j)n^{-jp}$, the zeta terms of the
perturbed node set, exactly as for the lattice. The pole of $\delta Z(-w)$ at $q=m/p$ lies to the left of the line, so
the band factor does not carry the pole term $pR_qW(q)\pi\cot(\pi q)n^{-1-pq}$ of Section~\ref{sec:conj}; in the
dictionary that term comes from the equilibrium, Proposition~\ref{prop:M12}. The $\tan$-pole at $w=(m+1)/p$ and the
change $\log|t'|$ of the node density sit at order $n^{-2-m}$, the order of the subleading pole terms of
Remark~\ref{rem:sublead}. None of this is a statement about $b_n$.
\end{remark}

% Section 5: The natural index and the expansion to n^{-6} (Theorem A = Theorem 5.7(iii)).
\section{The natural index and the expansion of the lattice}\label{sec:natural}

This section proves Theorem~\ref{thm:A}. The $1/n$ term and the $(s-1)^2$ part of the $n^{-2}$ term of the lattice are
one index shift (\S\ref{sec:shift}); for Bessel zeros the shift is exact to all orders (\S\ref{sec:bessel}). We then
extend the analysis of \cite{PaperII} and Theorem~\ref{thm:E2E3} to every order $K$ in $1/n$ and every order $B$ in the
band-factor parameters (\S\S\ref{sec:recursion}--\ref{sec:N6}). The analytic part is the order-$K$ version of
estimates already proved at order $2$; the algebraic part is a finite residue recursion, which we evaluate in exact
arithmetic. Strategy: the band factor is a tilt (Lemma~\ref{lem:Fprime}), tilts act on the jump data by an entire gauge
(Lemma~\ref{lem:gauge}), so every coefficient of the expansion is an entire function of the tilt parameters $\mu$ with
$n$-independent Taylor coefficients, and the Taylor coefficients are computed by the recursion.

\subsection{The index shift}\label{sec:shift}
Put $\delta_N:=\frac{\alpha+1}2$, so that $N=n+\delta_N$, and
\begin{equation}\label{eq:f2}
  f_2=f_2(p,\alpha):=\frac{p+1}{12}-\frac{(p+2)\alpha^2}8 .
\end{equation}

\begin{proposition}\label{prop:N1}
For all $p,s$, $d_1=-p\delta_N$ and $d_2(p,s)=\frac{p(p+1)}2\delta_N^2+f_2$. Equivalently,
$1+d_1/n+d_2/n^2=(n/N)^p(1+f_2N^{-2})+O(n^{-3})$.
\end{proposition}

\begin{proof}
With $s-1=p(\alpha+1)$ and $s-1-p=p\alpha$, the three terms of $d_2$ in Theorem~\ref{thm:E2E3} are
$p(p+1)(\alpha+1)^2/8=\frac{p(p+1)}2\delta_N^2$, $-(p+2)\alpha^2/8$ and $(p+1)/12$. Since
$(1+\delta/n)^{-p}=1-p\delta/n+\frac{p(p+1)}2\delta^2n^{-2}+O(n^{-3})$, the claim follows.
\end{proof}

So the lattice expansion to order $n^{-2}$ is an index shift by $\delta_N$ followed by a bare function with no $1/N$
term. The shift has two parts: $\frac\alpha2$, attached to the hard edge of Bessel order $\alpha$, and
$\frac12=-\zeta(0)$, the regularised count of the nodes (Section~\ref{sec:conj}).

\subsection{Bessel zeros}\label{sec:bessel}
For $\mu>-1$ let $j_{\mu,1}<j_{\mu,2}<\dotsb$ be the positive zeros of the Bessel function $J_\mu$ and
$\lambda_k:=(j_{\mu,k}/\pi)^2$. For $p=s=2$ the measure~\eqref{eq:sigmaL} is
$\sigma_\Lambda=\pi^2\sum_kj_{\mu,k}^{-2}\delta_{\pi^2/j_{\mu,k}^2}$, and $\mu=\frac12$ gives the lattice $\lambda_k=k^2$.

\begin{proposition}\label{prop:bessel}
The recurrence coefficients of $\sigma_\Lambda$ satisfy, for $n\ge1$, with $N:=n+\mu/2$,
\[
  \frac{n^2b_n}{c(2)}=\frac{4n^2}{(2n+\mu)\sqrt{(2n+\mu)^2-1}}=\Bigl(\frac nN\Bigr)^2\Bigl(1-\frac1{4N^2}\Bigr)^{-1/2}.
\]
The bare function $(1-\frac1{4N^2})^{-1/2}=1+\frac1{8N^2}+\frac3{128N^4}+\frac5{1024N^6}+\dotsb$ is even in $1/N$, and its
coefficients are $f_2(2,-\frac12)$, $f_4(2,-\frac12)$, $f_6(2,-\frac12)$ of~\eqref{eq:f2}, \eqref{eq:f4}, \eqref{eq:f6}.
\end{proposition}

\begin{proof}
The Lommel polynomials of order $\nu=\mu+1$ satisfy the monic recurrence $xh_n=h_{n+1}+\beta_n^2h_{n-1}$ with
$\beta_n^2=1/(4(n+\nu-1)(n+\nu))$, and their Jacobi operator has spectrum the closure of $\{\pm1/j_{\nu-1,k}\}$
\cite[\S3]{VA}; the orthogonality measure is symmetric with masses proportional to $j_{\nu-1,k}^{-2}$.
% TODO (citation): VA Sec. 3 is quoted only for the recurrence and the spectrum. The masses ~ j^{-2} follow from the
% partial fractions of the m-function -2 nu J_nu(1/x)/J_{nu-1}(1/x); this is done in PaperII Remark 9.6 only for
% nu = 3/2 (via DLMF 10.16.1, 4.22.3). For general nu quote the Mittag-Leffler expansion of J_nu/J_{nu-1}
% (DLMF 10.21-type) before submission; until then the identification is checked numerically (spec_lommel, 1e-15).
By the quadratic transformation of symmetric measures (as in \cite[Remark~9.6]{PaperII}), the image measure in
$y=x^2$ has $b_n^{(y)}=\beta_{2n-1}\beta_{2n}$, and $\sigma_\Lambda$ is its image under $y\mapsto\pi^2y$. Hence
$b_n=\pi^2/(4(2n+\mu)\sqrt{(2n+\mu)^2-1})$, and $c(2)=\pi^2/16$ gives the first identity. The second is
$(2N)\sqrt{4N^2-1}=4N^2(1-\frac1{4N^2})^{1/2}$.
\end{proof}

The first identity agrees with the Jacobi coefficients computed from the zeros (mpmath zeros for $k\le120$, McMahon's
expansion beyond, a continuous tail) to $1.6\cdot10^{-15}$ at $\mu=0$ and $\mu=1.5$, $n\le40$ (script
\texttt{spec\_lommel}). Here $\alpha=(s-1-p)/p=-\frac12$, and the large-$k$ form
$j_{\mu,k}=\pi(k+\frac\mu2-\frac14)+O(1/k)$ gives $Z_\Lambda(0)=\zeta(0,\frac\mu2+\frac34)=-\frac\mu2-\frac14$, so
$N=n+\frac\alpha2-Z_\Lambda(0)$: Proposition~\ref{prop:bessel} is the natural-index law of Section~\ref{sec:conj},
exactly.
% TODO (citation): McMahon's expansion (DLMF 10.21.19) is not quoted in the project's citation record.

\subsection{The residue recursion}\label{sec:recursion}
We use the frame of \S\ref{sec:frame} for $p\ne1$: the clockwise circles $\partial U_0$, $\partial U_r$ around $0$ and $4c$,
the local maps $f$ (hard edge) and $\tilde\varphi$ (soft edge), and the Szeg\H{o} quotients $\rho$, $\omega$ of
\cite[Sections~8--10]{PaperII}. Let $M:=\frac1{\sqrt2}\bigl(\begin{smallmatrix}1&i\\i&1\end{smallmatrix}\bigr)$, let $\circ$ be the
entrywise product, let $a_k=a_k(\alpha)$, $b_k=b_k(\alpha)$ be the coefficients of the Bessel expansions
\cite[(10.17.1), (10.17.8)]{DLMF} and $u_k,v_k$ those of the Airy expansions \cite[(9.7.2)]{DLMF}, and put
\[
  \Xi_k:=M^{-1}(M\circ C_k),\ C_k:=\begin{pmatrix}(-1)^ka_k&a_k\\(-1)^kb_k&b_k\end{pmatrix},\qquad
  X_k:=M^{-1}(M\circ C_k'),\ C_k':=\begin{pmatrix}u_k&(-1)^ku_k\\v_k&(-1)^kv_k\end{pmatrix}.
\]

\begin{lemma}[the jumps to every order]\label{lem:N6a}
Let $p\ne1$ and $K,J\ge1$. Uniformly on $\partial U_0\cup\partial U_r$, boundary values included, the jump of the error
matrix on the circles is $I+\sum_{k\le K}\Delta_kn^{-k}+O(n^{-K-1})$, where $\Delta_k$ equals
\begin{align*}
  \Delta_k^0&=(2f^{1/2})^{-k}P_0^\infty\rho^{\sigma_3}(\sigma_3\Xi_k\sigma_3)\rho^{-\sigma_3}(P_0^\infty)^{-1}&&\text{on }\partial U_0,\\
  \Delta_k^r&=(-1/\tilde\varphi)^kP_0^\infty(-\omega)^{-\sigma_3/2}X_k(-\omega)^{\sigma_3/2}(P_0^\infty)^{-1}&&\text{on }\partial U_r,
\end{align*}
and $\Delta_k=G_{\mu_n}\Delta_k^{\mathrm{mod}}G_{\mu_n}^{-1}+O(n^{-(J+1)p})$.
\end{lemma}

\begin{proof}[Proof sketch]
For $k=1$ this is \cite[Lemmas~9.1, 9.2 and~10.6]{PaperII}, and for $k=2$ the third-order version cited in
\S\ref{sec:E2E3}. The proof is the same with $K+1$ terms of the expansions \cite[(10.40.1)--(10.40.4), (10.17.5),
(10.17.6), (10.17.11), (10.17.12)]{DLMF} of the Bessel model \cite[\S6]{KMVV} and \cite[(9.7.5)--(9.7.6)]{DLMF} of the
Airy model, whose Poincar\'e remainders are $O(n^{-K-1})$ uniformly because the arguments are of size $n$ on the circles.
The sign pattern $(-1)^k$ is the same in every sector, for every $k$: each column of the Airy model is a constant
times $\mathrm{Ai}(Z)$ or $\mathrm{Ai}'(Z)$ with $Z=e^{2\pi im/3}\zeta$, $|\mathrm{ph}\,Z|\le\pi-\delta$, and
$Z^{3/2}=\pm\zeta^{3/2}$, so $(-1)^k$ sits on the column that carries $e^{-\xi}$ whichever branch occurs; the column
carrying $e^{\pm\xi}$ is fixed at $k=0$. The Bessel model is the same with $Y=\pm w$. The $n$-dependence of the weight is
Lemma~\ref{lem:Fprime} at order $J$, and the conjugation by $G_{\mu_n}$ is Lemma~\ref{lem:gauge}, whose proof commutes
diagonal factors and does not depend on $k$.
\end{proof}

\begin{lemma}[meromorphy]\label{lem:N6b}
For every $k$, $\Delta_k^{\mathrm{mod}}$ is meromorphic in $U_0$ and in $U_r$, and $\Delta_k^\mu:=G_\mu\Delta_k^{\mathrm{mod}}G_\mu^{-1}$
is meromorphic with Laurent coefficients that are entire functions of $\mu$.
\end{lemma}

\begin{proof}[Proof sketch]
For $p>1$: $\hat P^{\mathrm{mod}}(P_0^\infty)^{-1}$ is analytic in $U_e\setminus\{e\}$ and has the expansion of
Lemma~\ref{lem:N6a} uniformly on compact subsets, so $\Delta_k=\lim n^k(\cdots)$ is analytic and single-valued there; its
explicit formula grows at most polynomially at $e$, so $e$ is a pole. For $0<p<1$, where the glued hard-edge parametrix
carries seam jumps inside $U_0$, one checks directly that $\Delta_k$ has no jump: from the definitions,
$X_k=\frac12(u_k+v_k)I-\frac12(u_k-v_k)\sigma_2$ for even $k$ and $X_k=\frac12(u_k+v_k)\sigma_3+\frac12(v_k-u_k)i\sigma_1$ for odd
$k$, so $JX_kJ^{-1}=(-1)^kX_k$ for $J=i\sigma_2$; on $(4c-\epsilon,4c)$, $\tilde\varphi_+=-\tilde\varphi_-$ contributes the same
$(-1)^k$, and $P^\infty_{0,+}=P^\infty_{0,-}J$, $(-\omega)^{1/2}_+(-\omega)^{1/2}_-=1$. The hard edge is the same with
$\sigma_3\Xi_k\sigma_3$, $f^{1/2}_+=-f^{1/2}_-$ and the relations for $\rho$ of \cite[Section~10]{PaperII}. The gauge
$G_\mu$ is entire in $(z,\mu)$ and satisfies the same relations.
\end{proof}

For $\mu\in\C^J$ define, near each $e\in\{0,4c\}$, $\mathsf R_{\mathrm{in}}^{(0)}:=I$ and for $m\ge1$
\begin{equation}\label{eq:recursion}
\begin{gathered}
  \mathsf F_m(\mu):=\sum_{j=1}^m\mathsf R_{\mathrm{in}}^{(m-j)}(\mu)\Delta_j^\mu,\qquad
  \mathsf R_{\mathrm{out}}^{(m)}(\mu):=\sum_{e\in\{0,4c\}}\bigl(\text{principal part of }\mathsf F_m(\mu)\text{ at }e\bigr),\\
  \mathsf R_{\mathrm{in}}^{(m)}(\mu):=\mathsf R_{\mathrm{out}}^{(m)}(\mu)-\mathsf F_m(\mu),\qquad
  \mathsf R_1^{(m)}(\mu):=\Res_0\mathsf F_m(\mu)+\Res_{4c}\mathsf F_m(\mu).
\end{gathered}
\end{equation}
By Lemma~\ref{lem:N6b} these are well defined, $\mathsf R_{\mathrm{in}}^{(m)}$ is analytic in each disc, and all entries are
entire in $\mu$. For $m=1,2$ and $\mu=0$ they are the objects of \cite[Proposition~9.5]{PaperII} and its third-order
version.

\begin{proposition}[small norm to order $K$]\label{prop:N6c}
Let $p\ne1$ and $K,J\ge1$. Then $\mathsf R_1=\sum_{m\le K}\mathsf R_1^{(m)}(\mu_n)n^{-m}+O(n^{-K-1}+n^{-1-(J+1)p})$, and the
same expansion holds for $\mathsf R$ uniformly at positive distance from its jump contour.
\end{proposition}

\begin{proof}[Proof sketch]
Write $\eta_n:=n^{-K-1}+n^{-1-(J+1)p}$. By Lemma~\ref{lem:N6a} the jump is $I+\sum_{k\le K}\Delta_k^{\mu_n}n^{-k}$ on the circles
up to $O(\eta_n)$, and off the circles it is $I$ up to errors that are $O(\eta_n)$ in $L^1\cap L^2\cap L^\infty$, as in
\cite[Propositions~9.5 and~10.8]{PaperII}; the Cauchy operator is bounded uniformly in $n$. One proves by induction on
$m=1,\dots,K+1$ that the boundary value $\mathsf R_-$ equals $\sum_{k<m}\mathsf R_-^{(k)}n^{-k}+O_{L^2}(n^{-m}+n^{-1-(J+1)p})$, where
$\mathsf R^{(k)}$ is $\mathsf R^{(k)}_{\mathrm{in}}(\mu_n)$ inside the discs and $\mathsf R^{(k)}_{\mathrm{out}}(\mu_n)$ outside. In the
step, insert the claim into $\mathsf R=I+\mathcal C(\mathsf R_-(J-I))$: the new error is
$O(n^{-m-1}+n^{-1-(J+1)p})$, and the term of order $n^{-m}$ is the Cauchy integral over the clockwise circles of
$\mathsf F_m(\mu_n)$, which by the residue theorem is $\mathsf R_{\mathrm{out}}^{(m)}$ outside the discs and
$\mathsf R_{\mathrm{out}}^{(m)}-\mathsf F_m=\mathsf R_{\mathrm{in}}^{(m)}$ inside. The coefficient $\mathsf R_1$ is
$-\frac1{2\pi i}\int\mathsf R_-(J-I)$, and $-\frac1{2\pi i}\oint_{\mathrm{cw}}\mathsf F_m=\Res_0\mathsf F_m+\Res_{4c}\mathsf F_m$.
\end{proof}

Define $\cD_K(n;\mu):=((Q_1)_{12}(Q_1)_{21})^{1/2}/c$, expanded to order $n^{-K}$, where
$Q_1=P_1+\sum_{m\le K}\mathsf R_1^{(m)}(\mu)n^{-m}$, and write $\cD_K(n;\mu)=\sum_{\br}\mu^{\br}\cD_{\br,K}(n)$ for its
Taylor series in $\mu$. Each $\cD_{\br,K}$ is a polynomial in $1/n$ of degree at most $K$. A polynomial in $1/n$ is
\emph{written in natural form} when it is re-expanded in $1/N$ by $1/n=N^{-1}/(1-\delta_NN^{-1})$; a layer $\mu^{\br}$ carries
$n^{-|\br|_wp}=N^{-|\br|_wp}(N/n)^{|\br|_wp}$, with $|\br|_w:=\sum_jjr_j$.

\begin{lemma}[exact algebra]\label{lem:N6d}
The recursion~\eqref{eq:recursion} and the extraction $\cD_K$ give, identically in the symbols:
\begin{enumerate}
\item[(a)] $\mu=0$, $K=6$: $\cD_{0,6}(n)=(n/N)^p[1+f_2N^{-2}+f_4N^{-4}+f_6N^{-6}+O(N^{-7})]$, with no $N^{-1},N^{-3},N^{-5}$
terms, $f_2$ as in~\eqref{eq:f2} and $f_4,f_6$ as in~\eqref{eq:f4}, \eqref{eq:f6}; in particular $d_1=-p(\alpha+1)/2$ and $d_2$ is
that of Theorem~\ref{thm:E2E3};
\item[(b)] first order in $\mu_j$, for every integer $j\ge1$ ($j$ a symbol), $K=4$: the coefficient of $n^{-1}$ per unit
$\binom{2j}j\mu_j$ is $pj/2$ (with $c=1$), and in natural form the layer is a multiple of
$N^{-1-jp}[1+g_jN^{-2}+O(N^{-4})]$, with no $N^{-1}$, $N^{-3}$ term in the bracket and $g_j=A_j+B_j\alpha^2$ as in
Theorem~\ref{thm:A}; the same holds at $K=5$ for the integers $j=1,\dots,5$, and for $j=1$ the $N^{-4}$ coefficient of
the bracket is $g_{1,4}$ of~\eqref{eq:g14};
\item[(c)] second order in $\mu_1$, $K=5$: per unit $\mu_1^2$ the natural form has no $N^{-1-2p}$ term, the coefficient
$p(3p+1)/2$ at $N^{-2-2p}$, no $N^{-3-2p}$ and no $N^{-5-2p}$ term, and $p^2h_2$ at $N^{-4-2p}$, with $h_2$ as in~\eqref{eq:h2}.
\end{enumerate}
\end{lemma}

\begin{proof}
With $c=1$ every input is a rational function of $(p,\alpha)$: the local maps have convergent series with rational
coefficients, $\pi$ cancels in $f$, the square roots are of series whose leading coefficient is $1$ or a perfect square,
and the factors $2^{-1/2}$ in $\mathsf T$ and $M$ occur in conjugate pairs. The recursion is carried out on truncated
Laurent series over $\mathbb Q(p,\alpha)[i]$; every series carries the index of its highest exactly known coefficient, and
every principal part checks that its coefficients down to the residue are known, otherwise the run stops. For (b),
$\mu_j$ is normalised by $\binom{2j}j$; the gauge to first order is $\partial_{\mu_j}\Delta_k^\mu|_{\mu=0}=[\mathcal G_j,\Delta_k^{\mathrm{mod}}]$
with $\mathcal G_j=\mathsf T\bigl(\begin{smallmatrix}0&-(z-4c)Q_j/2\\-zQ_j/2&0\end{smallmatrix}\bigr)\mathsf T^{-1}$, and it needs only the
Taylor data of $Q_j/\binom{2j}j$ at $0$ and at $4c=4$, which are rational in $j$: $\prod_{i\le m}(j-i)/(2(2j-2i-1))$ at $0$
and $j(j-1)\cdots(j-m)/(4^{m+1}(\frac12)_{m+1})$ at $4$, from $\sum_jQ_j(z)t^j=t/((1-zt)\sqrt{1-4t})$. So (b) is a
computation over $\mathbb Q(p,\alpha,j)[i]$. For (c) the gauge is taken to second order,
$\cosh h=1+z(z-4)\mu^2Q_1^2/8+O(\mu^4)$, $S=\mu Q_1/2+O(\mu^3)$. The scripts and their outputs are listed in
Appendix~\ref{app:repro}; each output begins with the SHA-256 hashes of the scripts that produced it. Since $(z-4c)Q_j$ and
$zQ_j$ scale as $c^j$ under $z=c\zeta$, $\mu_j$ enters only as $c^j\mu_j$; the integer layer and the ratios $g_j$, $h_2$
are independent of $c$.
\end{proof}

The closed forms are
\begin{equation}\label{eq:f4}
  f_4=\frac{(p+1)(6p^3+46p^2+55p+3)}{1440\,p(3p+2)}-\frac{5(p+1)(p+2)\alpha^2}{96(3p+2)}-\frac{p(p+2)^2\alpha^4}{128(3p+2)},
\end{equation}
\begin{equation}\label{eq:f6}
  f_6=\frac{\phi_0+\phi_2\alpha^2+\phi_4\alpha^4+\phi_6\alpha^6}{(3p+2)^2(5p+2)},
\end{equation}
\begin{align*}
  \phi_0&=-\frac{p+1}{362880\,p^2}\bigl(450p^8+1590p^7-316p^6-23761p^5-66864p^4-65507p^3\\
  &\qquad\qquad-18220p^2-1632p-40\bigr),\\
  \phi_2&=\frac{(p+1)(p+2)(90p^6+270p^5+664p^4-165p^3-1649p^2-288p+28)}{11520\,p},\\
  \phi_4&=-\frac{(p+1)(p+2)^2(20p^4+45p^3+162p^2+72p+16)}{1536},\\
  \phi_6&=\frac{p(p+2)^3(16p^3+7p^2-4p-4)}{3072},
\end{align*}
\begin{multline}\label{eq:g14}
  g_{1,4}=\frac{p+1}{5760\,p^2(3p+2)}\bigl(9p^6+117p^5+161p^4-569p^3-1550p^2-976p+48\bigr)\\
  +\frac{(p+1)(p+2)(3p^4+33p^3+124p^2+182p+108)\alpha^2}{192\,p(3p+2)}\\
  +\frac{(p+2)^3(p^3-4p-4)\alpha^4}{128\,p(3p+2)},
\end{multline}
\begin{equation}\label{eq:h2}
  h_2=\frac{(p+1)(p+2)(3p^2+2p+1)}{8p^2}+\frac{(p+2)(2p^3-p^2-7p-6)\alpha^2}{16p^2}.
\end{equation}
For $j=1$, $g_1=(p+1)(p^2+3p-2)/(24p)+(p+2)(p^2+2p+2)\alpha^2/(8p)$. At the Bessel point $(p,\alpha)=(2,-\frac12)$,
$f_4=\frac3{128}$ and $f_6=\frac5{1024}$, in agreement with Proposition~\ref{prop:bessel}.

\subsection{The theorem}\label{sec:N6}
\begin{theorem}\label{thm:N6}
Let $p>0$, $p\ne1$, $s>1$, and $K,J,B\ge1$.
\begin{enumerate}
\item[(i)] \emph{Double expansion.} $\displaystyle\frac{n^pb_n}c=\sum_{|\br|\le B}\mu_n^{\br}\cD_{\br,K}(n)
  +O\bigl(n^{-K-1}+n^{-2-(B+1)p}+n^{-1-(J+1)p}\bigr)$.
\item[(ii)] \emph{Natural forms of the layers.} $\cD_{0,6}$, $\mu_{n,j}\cD_{e_j,4}$ and $\mu_{n,1}^2\cD_{2e_1,5}$ satisfy
\begin{gather*}
  \cD_{0,6}(n)=\Bigl(\frac nN\Bigr)^p\Bigl[1+\frac{f_2}{N^2}+\frac{f_4}{N^4}+\frac{f_6}{N^6}+O(N^{-7})\Bigr],\\
  \mu_{n,j}\cD_{e_j,4}(n)=\Bigl(\frac nN\Bigr)^p\frac{\varepsilon_j}{N^{1+jp}}\Bigl[1+\frac{g_j}{N^2}+O(N^{-4})\Bigr]\qquad(j\ge1),\\
  \mu_{n,1}^2\cD_{2e_1,5}(n)=\Bigl(\frac nN\Bigr)^p\frac{\varepsilon_1^2}{N^{2+2p}}\Bigl[\frac{3p+1}{2p}+\frac{h_2}{N^2}+O(N^{-4})\Bigr],
\end{gather*}
with the brackets free of odd powers of $1/N$ to the orders shown; for $j=1$ the second bracket is
$1+g_1N^{-2}+g_{1,4}N^{-4}+O(N^{-5})$.
\item[(iii)] \emph{The coefficients of $b_n$.} With $K=6$, $B=2$ and $J\ge6/p-1$, expansion~\eqref{eq:thmA} holds.
\end{enumerate}
\end{theorem}

\begin{proof}
(i) By~\eqref{eq:extract} and Proposition~\ref{prop:N6c}, $n^{2p}b_n^2=(Q_1)_{12}(Q_1)_{21}$ with
$Q_1=P_1+\sum_{m\le K}\mathsf R_1^{(m)}(\mu_n)n^{-m}+O(n^{-K-1}+n^{-1-(J+1)p})$. The $n^0$ term of $(Q_1)_{12}(Q_1)_{21}$ is
$c^2$, since $P_1$ does not depend on $\mu$, so each coefficient of $\cD_K(n;\mu)$ is a polynomial in the entries of the
$\mathsf R_1^{(i)}(\mu)$, $i\le K$, which are entire in $\mu$ and independent of $n$. Taylor's theorem in $\mu$ on $|\mu|\le1$
applies coefficientwise. The $n^0$ coefficient is $1$ for all $\mu$, and the $n^{-1}$ coefficient is $d_1+cpM_\mu(4c)$,
exactly linear in $\mu$ by Lemma~\ref{lem:E37}; hence $\cD_{\br,K}=O(n^{-2})$ for $|\br|\ge2$, and the Taylor remainder
after order $B$ is $O(|\mu|^{B+1}n^{-2})$. Since $|\mu_n|=O(n^{-p})$, this is $O(n^{-2-(B+1)p})$.

(ii) This is Lemma~\ref{lem:N6d} in natural form, with the normalisations
$\mu_{n,j}\binom{2j}jpc^jj/2=\varepsilon_jn^{-jp}$ and $\mu_{n,1}^2p(3p+1)/2\cdot c^2=\varepsilon_1^2n^{-2p}(3p+1)/(2p)$.

(iii) In~(i) keep $\br=0$ at $K=6$; $\br=e_j$ for every $j$ with $1+jp<\min(7,2+3p)$, at $K=5$ ($j\le5$) or $K=4$; and
$\br=2e_1$ at $K=5$. The other terms with $|\br|=2$ ($\mu_1\mu_2,\mu_2^2,\dots$) are $O(n^{-2-3p})$ because
$\cD_\br=O(n^{-2})$ and $|\br|_w\ge3$, and those with $|\br|=3$ are $O(n^{-2-3p})$. The truncation errors are bounded as
follows. The $j=1$ bracket is known to relative order $N^{-5}$, an error $O(n^{-6-p})$, which is $O(n^{-2-3p})$ for $p<2$
and $O(n^{-8})$ for $p\ge2$. For $j\ge2$ the bracket error is $O(n^{-5-jp})$, and $5+2p>2+3p$ for $p<3$, $5+2p\ge11$ for
$p\ge3$. The $h_2$ term of the $\varepsilon_1^2$ layer is $O(n^{-4-2p})$, which is $O(n^{-2-3p})$ for $p<2$ and $O(n^{-7})$ for
$p\ge\frac32$. Finally $n^{-1-(J+1)p}=O(n^{-7})$.
\end{proof}

A term $CN^{-e}$ of~\eqref{eq:thmA} is a coefficient of the expansion of $n^pb_n/c$ (in powers of $N$, or equivalently,
after re-expansion, of $n$) exactly when $e<\min(7,2+3p)$. Table~\ref{tab:visible} lists the resulting ranges. The terms
outside their range are true identities of the layers in Theorem~\ref{thm:N6}(ii) but are not determined as coefficients
of $b_n$, because the uncomputed layers $\mu_1\mu_2$, $\mu_1^3$ enter at $n^{-2-3p}$.

\begin{table}[ht]
\centering
\small
\begin{tabular}{lll}
\toprule
term & exponent $e$ & a coefficient of $n^pb_n/c$ for\\
\midrule
$f_2N^{-2}$ & $2$ & all $p\ne1$\\
$f_4N^{-4}$ & $4$ & $p>\frac23$\\
$f_6N^{-6}$ & $6$ & $p>\frac43$\\
$\varepsilon_jN^{-1-jp}$ & $1+jp$ & $(j-3)p<1$ and $1+jp<7$\\
$\varepsilon_1g_1N^{-3-p}$ & $3+p$ & $\frac12<p<4$\\
$\varepsilon_2g_2N^{-3-2p}$ & $3+2p$ & $1<p<2$\\
$\varepsilon_1g_{1,4}N^{-5-p}$ & $5+p$ & $\frac32<p<2$\\
$\frac{3p+1}{2p}\varepsilon_1^2N^{-2-2p}$ & $2+2p$ & $p<\frac52$\\
$\varepsilon_jg_jN^{-3-jp}$ ($j\ge3$), $\varepsilon_1^2h_2N^{-4-2p}$ & & never at $B=2$\\
\bottomrule
\end{tabular}
\caption{Which terms of~\eqref{eq:thmA} are coefficients of $b_n$ ($p\ne1$). Coincident exponents merge: at $p=1/j_0$,
$\varepsilon_{j_0}N^{-2}$ merges with $f_2N^{-2}$, and $f_2+\varepsilon_{j_0}=f_2-p/12$, as in Theorem~\ref{thm:E2E3}.}
\label{tab:visible}
\end{table}

\begin{remark}\label{rem:N6}
(a) Theorem~\ref{thm:N6} does not prove evenness in $1/N$ to all orders, nor any closed form of the bare function; the
computed layers are listed in~(ii), and the layers $\mu_1\mu_2$, $\mu_2^2$, $\mu_1^3$ and higher are not computed. Evenness of the
integer layer to all orders is Corollary~\ref{cor:even} below. Nothing
is claimed at $p=1$, where the hard-edge and band-factor constructions of \cite{PaperII} change.
(b) The same recursion evaluated numerically at $40$ digits, with truncation depths $40$--$64$, gives $f_1=f_3=f_5=f_7=0$ to
$10^{-35}$ at every sampled $(p,\alpha)$ with $p\in[\frac13,12]$, each $f_{2k}$ even in $\alpha$, and
$f_8(3,0)=1311857/632311515$, $f_8(2,0)=3899/2^{21}$ after rational recognition. The denominators of $f_4$, $f_6$ are
$3p+2$ and $(3p+2)^2(5p+2)$; they come from the factors $1/(a'+k)$, $a'=\frac1p+\frac12$, of the hard-edge local map,
$a'+k=((2k+1)p+2)/(2p)$, and $632311515=3^5\cdot5\cdot11^3\cdot17\cdot23$ is $3^5\cdot5\,(3p+2)^3(5p+2)(7p+2)$ at $p=3$
(scripts \texttt{m4\_rh4}, \texttt{m4\_f8\_denominator}). These statements are numerical, except $f_7=0$ and the evenness in
$\alpha$, which hold to all orders by Corollary~\ref{cor:even}.
(c) Heuristically, evenness in $\alpha$ after the shift reflects the fact that the Bessel coefficients $a_k(\alpha)$,
$b_k(\alpha)$ depend on $\alpha$ only through $\alpha^2$, while elsewhere $\alpha$ enters only through the power of $z$ at the
hard edge, which the shift by $\frac\alpha2$ absorbs. Section~\ref{sec:shiftsec} makes this precise: the two roles of $\alpha$
separate, the Szeg\H{o} exponent is absorbed by the natural index (Theorem~\ref{thm:shift}), and the evenness in $1/N$ is a
reflection symmetry (Theorem~\ref{thm:reflect}, Corollary~\ref{cor:even}).
(d) The factor $1/(2j-1)$ in $g_j$, read as a function of a continuous Mellin variable $w$, is a simple pole at
$w=\frac12$, the Mellin signature of $(1-u)^{+1/2}$, against $(1-u)^{-1/2}$ for the leading kernel
(Lemma~\ref{lem:K0}). Partial fractions give
$B_j=\frac{p+2}{8p}\bigl[\frac{p^2j}2+\frac{p(p+2)}4+\frac{(p+2)(p+4)}{4(2j-1)}\bigr]$. This suggests a second-order edge kernel
(Problem~\ref{prob:kernel2}).
\end{remark}

% Section 6: Reflection, shift and the Toda structure (from research/G1/RH/spectral/JZC_M4.md Sec. 7-8, JZC_SIGMA.md, JZC_TODA.md, JZC_SHIFT.md;
% updated 2026-10-23 after JZC_TODA / JZC_SHIFT review; attributions added (Rule L, 2026-10-23) from the drafted passages of
% research/G1/RH/spectral/NOVELTY_TODA.md Sec. 5, all sources second-hand (WebFetch), Rule C reading still due).
\section{Reflection, shift and the Toda structure}\label{sec:shiftsec}

Theorem~\ref{thm:N6} computes the layers of the lattice expansion to fixed orders and observes that they are even in $1/N$.
This section explains the evenness and the natural index to all orders. The exponent $\alpha$ plays two roles in the
residue recursion, and separating them gives a two-parameter family of formal problems (\S\ref{sec:decoupled}). On this
family there are two exact symmetries: a reflection, which is algebra (Theorem~\ref{thm:reflect}), and a shift, which we
prove by steepest descent of the same lattice problem in a second, non-integer frame (Theorem~\ref{thm:shift}). Together they
give the evenness of the integer layer to all orders (Corollary~\ref{cor:even}). The zeta layers are Toda flows of the integer
layers of $b_n$ and $a_n$, and on the hard edge their dispersionless limit gives the leading term of every zeta layer in closed form
(\S\S\ref{sec:toda}--\ref{sec:indexshift}).

\subsection{The decoupled model}\label{sec:decoupled}
The exponent $\alpha$ enters the jump data of Lemma~\ref{lem:N6a} in two ways: as the Bessel order $\nu$ in the coefficients
$a_k(\nu),b_k(\nu)$, which are polynomials in $\nu^2$, and as the Szeg\H{o} exponent $\beta$ in $\rho^2=\phi^{2(\beta+1)}$ at the hard
edge and in the factor $\phi^{-2(\beta+1)}$ at the soft edge, where $\phi(z)=(z^{1/2}+(z-4c)^{1/2})/(2c^{1/2})$ is the exterior
conformal map of $[0,4c]$. For the lattice $\nu=\beta=\alpha$. Keeping $\nu$ and $\beta$ independent, the residue
recursion~\eqref{eq:recursion} and the extraction of \S\ref{sec:recursion}, with the gauge of Lemma~\ref{lem:gauge} for the
band-factor parameters $\mu$, define formal series
\begin{equation}\label{eq:decoupled}
  \frac{n^p\,b(n;\beta,\nu;\mu)}{c}=\sum_{\br}\mu^{\br}\,\cD_{\br}(n;\beta,\nu),\qquad
  B(n;\beta,\nu):=\cD_0(n;\beta,\nu)=\sum_{k\ge0}d_k(p,\beta,\nu)\,n^{-k},
\end{equation}
whose truncations at order $n^{-K}$ are the polynomials $\cD_{\br,K}$ of \S\ref{sec:recursion} when $\beta=\nu=\alpha$. We call
this the \emph{decoupled model} $M(\beta,\nu)$ and $B$ its \emph{integer layer}. Each $d_k$ is a rational function of
$(p,\beta,\nu)$, $d_0=1$: with $c=1$ the recursion and the extraction are field operations on Laurent data that are rational
in $p$ (Lemma~\ref{lem:N6d}), polynomial in $\nu$ through $a_k(\nu),b_k(\nu)$, and polynomial in $\beta$, since $\beta$ enters
only through binomial series of $(1+\text{series without constant term})^{2(\beta+1)}$.

\begin{theorem}[reflection]\label{thm:reflect}
As formal series, $b(n;\beta,\nu;0)=b(-n;-\beta-2,\nu;0)$; that is, $d_k(p,-\beta-2,\nu)=(-1)^kd_k(p,\beta,\nu)$ for all $k$.
\end{theorem}

\begin{proof}
With $T=\frac1{\sqrt2}\left(\begin{smallmatrix}1&1\\ i&-i\end{smallmatrix}\right)$ we have $\sigma_2T=T\sigma_3$, so
$P_0^\infty=T\mathfrak a_\infty^{\sigma_3}T^{-1}$ and $P_1=-c\sigma_2$ commute with $\sigma_2$. The core matrices $X_k$ and
$\sigma_3\Xi_k\sigma_3$ lie in the span of $I,\sigma_2$ for even $k$ and of $\sigma_3,i\sigma_1$ for odd $k$, so conjugation by
$\sigma_2$ multiplies them by $(-1)^k$, while $\sigma_2d^{\sigma_3}\sigma_2=d^{-\sigma_3}$ for scalars $d$. Hence
$\sigma_2\Delta_k(\beta)n^{-k}\sigma_2=\Delta_k(-\beta-2)(-n)^{-k}$, the recursion is covariant under conjugation by a constant
matrix, and $(Q_1)_{12}(Q_1)_{21}$ is invariant under conjugation by $\sigma_2$.
\end{proof}

The grading of the jump corrections into $\operatorname{span}\{I,\sigma_2\}$ and $\operatorname{span}\{\sigma_3,i\sigma_1\}$ is the device
that Kuijlaars and Tibboel \cite[Lemma~4.1]{KuijlaarsTibboel} use to show that the odd powers of $1/n$ in the off-diagonal recurrence
coefficient vanish for one-cut varying weights; here it gives a covariance under $(n,\beta)\mapsto(-n,-\beta-2)$.

\begin{lemma}[self-similar Buyarov--Rakhmanov identity]\label{lem:BR}
Let $\Phi$ denote the phase of the lattice equilibrium in the scaled frame, with hard-edge branch $2f^{1/2}$ and soft-edge
branch $\tilde\varphi$. Then $(1+pz\partial_z)\Phi=-2\log\phi$ up to the additive constants of the branches; in particular
$(1+pz\partial_z)(2f^{1/2})=2\operatorname{arsinh}(v/2)$ at $z=-v^2$ ($c=1$).
\end{lemma}

\begin{proof}
At the hard edge $2f^{1/2}=\pi\sqrt{-z}\,\mathfrak a(z)$, and the series of $\mathfrak a$ has the factors $1/(a'+k)$ with
$a'=\frac1p+\frac12$. Since $(1+pz\partial_z)[\sqrt{-z}\,z^k]=p(a'+k)\sqrt{-z}\,z^k$, these factors cancel term by term and
leave the arcsine series. The soft-edge branch is the analytic continuation of the same phase.
\end{proof}

Under the self-similar scaling of the lattice problem, $1+pz\partial_z$ is the derivative in $n$ of $n\Phi(xn^p)$; so
Lemma~\ref{lem:BR} says that the $n$-flow of the phase is generated by $\log\phi$, which also generates the $\beta$-flow
($\partial_\beta\rho^2=2\log\phi\cdot\rho^2$). This suggested the following theorem; its proof does not use the lemma.

\subsection{The shift theorem}\label{sec:shift-thm}
\begin{theorem}[shift]\label{thm:shift}
Identically in $(p,\beta,\nu)$, as formal series:
\begin{enumerate}
\item[(i)] $B(n;\beta,\nu)=(n/N)^pF_p(N;\nu)$ with $N:=n+\frac{\beta+1}2$, where $F_p(N;\nu)$ is a formal power series in $1/N$
that does not depend on $\beta$; equivalently, for every real $\gamma$ (the \emph{frame identity})
\begin{equation}\label{eq:FI}
  B(n+\gamma;\beta-2\gamma,\nu)=\Bigl(1+\frac\gamma n\Bigr)^pB(n;\beta,\nu),
\end{equation}
the left side re-expanded in $1/n$;
\item[(ii)] for $\br\in\{e_1,e_2,2e_1\}$ the same holds with the exponent $p$ replaced by $(1+|\br|_w)p$: the series
$(N/n)^{(1+|\br|_w)p}\cD_{\br}(n;\beta,\nu)$, re-expanded in $1/N$, does not depend on $\beta$.
\end{enumerate}
\end{theorem}

The infinitesimal form of~\eqref{eq:FI} is the flow identity $\partial_nB-2\partial_\beta B=(p/n)B$. The equivalence of~(i)
and~\eqref{eq:FI} is algebra: if (i) holds, then $(n+\gamma)+\frac{\beta-2\gamma+1}2=N$ and
$B(n+\gamma;\beta-2\gamma,\nu)=((n+\gamma)/N)^pF_p(N)=(1+\gamma/n)^p(n/N)^pF_p(N)$. Conversely, given~\eqref{eq:FI}, fix
$\beta_0$ and put $\gamma:=(\beta_0-\beta)/2$; then $(N/n)^pB(n;\beta,\nu)=(N/(n-\gamma))^pB(n-\gamma;\beta_0,\nu)$, and
$n-\gamma=N-\frac{\beta_0+1}2$ does not involve $\beta$. All re-expansions are between formal series in $1/n$, $1/(n-\gamma)$,
$1/N$, which differ by constants.

\begin{corollary}[evenness]\label{cor:even}
$F_p(N;\nu)=F_p(-N;\nu)$, and $F_p$ depends on $\nu$ only through $\nu^2$. Hence $F_p(N;\nu)=1+\sum_{k\ge1}f_{2k}N^{-2k}$ with
each $f_{2k}$ rational in $p$ and polynomial in $\nu^2$. For the lattice and every $p\ne1$, the integer layer of
Theorem~\ref{thm:N6} is, for every $K$,
\[
  \cD_{0,K}(n)=\Bigl(\frac nN\Bigr)^p\Bigl[1+\sum_{2k\le K}\frac{f_{2k}(p,\alpha^2)}{N^{2k}}\Bigr]+O(N^{-K-1}),\qquad N=n+\frac\alpha2+\frac12 ;
\]
in particular all odd coefficients $f_{2k+1}$ vanish, and the natural index of the lattice is the Szeg\H{o} shift $\frac{\beta+1}2$ at
$\beta=\alpha$.
\end{corollary}

\begin{proof}
By Theorem~\ref{thm:reflect}, $B(n;\beta,\nu)=B(-n;-\beta-2,\nu)$. The natural index of the right side is
$-n+\frac{-\beta-1}2=-N$, so Theorem~\ref{thm:shift}(i) gives $(n/N)^pF_p(N)=((-n)/(-N))^pF_p(-N)$. The coefficients $a_k(\nu)$,
$b_k(\nu)$ depend on $\nu$ only through $\nu^2$, and $\beta$ no longer occurs in $F_p$. Each $d_k$ is polynomial in $\nu$, hence
so is each $f_{2k}$.
\end{proof}

Corollary~\ref{cor:even} settles the integer-layer part of Remark~\ref{rem:N6}(a)--(c): $f_7=0$ and the evenness in $\alpha$ hold to
all orders, while the values $f_8,f_{10},\dots$ and the observed degree $k$ of $f_{2k}$ in $\alpha^2$ remain numerical. The $\alpha$-even
part has the heuristic explanation of Remark~\ref{rem:N6}(c): $\alpha$ enters the integer layer through $\nu^2$ and through $\beta$,
and the shift absorbs $\beta$.

Evenness in a half-shifted index is known for varying weights with a string equation \cite[Theorem~5.2]{BleherIts}, \cite{EMP,KuijlaarsTibboel}.
In the semiclassical Laguerre case the index and the hard-edge exponent enter the discrete Painlev\'e system through $n+\alpha/2$ and
$\alpha^2$ \cite[Theorem~2.1]{FilipukVanAsscheZhang}. What is new here, as far as we know, is the decoupling of the Szeg\H{o} exponent
from the Bessel order and the frame identity~\eqref{eq:FI} for a fixed discrete measure.

\subsection{Proof of the shift theorem: a second frame}\label{sec:shift-proof}
The decoupled model is not related to itself by any transformation of the formal recursion that we know. In particular,
the Christoffel shift $x\sigma_{p,s}=\sigma_{p,s+p}$ of \cite{PaperII} moves the Bessel order with the Szeg\H{o} exponent,
$(\beta,\nu)\mapsto(\beta+1,\nu+1)$, and any polynomial modification of the measure that is regular and nonzero at $0$ moves
neither. The proof instead analyses the same exact object twice: the degree-$n$ Riemann--Hilbert problem $Y$ of the lattice
is steepest-descended in the frame of an arbitrary real index $m=n+\gamma$, and the two resulting expansions of the same
number $b_n$ are compared.

Let $U_m(z):=m^{pn\sigma_3}Y(zm^{-p})$, the rescaling of \cite{PaperII} with $m^p$ in place of $n^p$; the degree stays $n$. Put
$t_m:=mz^{-1/p}$, so that the nodes $X_k=(m/k)^p$ satisfy $t_m(X_k)=k$, and let $g$, $\phi$, $\tilde\varphi$ be the lattice
$g$-function and phases of \cite{PaperII}, which do not depend on the index. With the principal branch of $z^\gamma$ (cut
$(-\infty,0]$) and the constant $C=m^{p\gamma\sigma_3}$, consider
\[
  \tilde U:=C\,U_mz^{\gamma\sigma_3}C^{-1},\qquad \tilde Uz^{-m\sigma_3}=C(U_mz^{-n\sigma_3})C^{-1}\to I\quad(z\to\infty).
\]
The conjugation by a constant diagonal matrix leaves the product of the off-diagonal entries of the $1/z$ coefficient unchanged,
so $\tilde U$ still encodes $m^{2p}b_n^2$.

\begin{lemma}[the $m$-frame]\label{lem:frame}
(i) The residue conditions of $\tilde U$ are those of the lattice measure $\sigma_{p,\tilde s}$, $\tilde s=s-2p\gamma$, in the frame
of index $m$, so that its band weight has the exponent $\tilde\alpha:=\alpha-2\gamma$. After Steps~1--4 of \cite[Section~5]{PaperII} with
$n$ replaced by $m$ in every place except the degree, the jumps are those of \cite{PaperII} with weight exponent $\tilde\alpha$, there
is no jump on $(-\infty,0)$, and the $(2,1)$ entry of the lip jumps in the hard-edge disc is $\omega^\natural e^{-4\zeta^{1/2}}$,
$\zeta=m^2f$, with $\omega^\natural=\omega^{(\tilde\alpha)}e^{\mp2\pi i\gamma}$ in $\C_\pm$, whose boundary values on the band are
$e^{\mp i\pi\alpha}$. That is, the $m$-frame problem has Bessel order $\alpha$ and Szeg\H{o} exponent $\alpha-2\gamma$.

(ii) Let $p>1$, $|\alpha|<1$ and $0\le\gamma<(1-|\alpha|)/2$. Then for all $K,J,B\ge1$, with $\mu_{m,j}=2\zeta(-jp)/(jm^{jp})$,
\[
  \frac{m^pb_n}c=\sum_{|\br|\le B}\mu_m^{\br}\,\cD_{\br,K}(m;\alpha-2\gamma,\alpha)+O\bigl(m^{-K-1}+m^{-2-(B+1)p}+m^{-1-(J+1)p}\bigr).
\]
\end{lemma}

\begin{proof}[Proof sketch]
(i) The residue of $\tilde U_{j2}$ at $X_k$ acquires the factor $(c_1/c_2)X_k^{-2\gamma}=k^{2p\gamma}$, which turns the mass $k^{-s}$ into
$k^{-\tilde s}$. The computation of the jumps in \cite[Section~5]{PaperII} uses the integrality of the index only on $(-\infty,0)$,
where $e^{m\phi\sigma_3}$ jumps by $e^{-2\pi im\sigma_3}$ and the cut of $z^{\gamma\sigma_3}$ contributes $e^{2\pi i\gamma\sigma_3}$; the
product is $e^{-2\pi in\sigma_3}=I$ because the degree $n$ is an integer. In the hard-edge disc $\phi=\mp i\pi-2f^{1/2}$ in $\C_\pm$
\cite{PaperII}, so $e^{2m\phi}=e^{\mp2\pi i\gamma}e^{-4m f^{1/2}}$; the Szeg\H{o} quotient $\omega^{(\tilde\alpha)}$ tends to
$e^{\mp i\pi\tilde\alpha}$ on the band, and $e^{\mp i\pi(\tilde\alpha+2\gamma)}=e^{\mp i\pi\alpha}$. At the soft edge every jump is written
with $\tilde\varphi$, which does not depend on the index. The normalisation at infinity is that of \cite{PaperII} with $m$, since
$g-\operatorname{Log}z$ has no jump on $(-\infty,0)$ for every real $m$.

(ii) These are the jumps of the Hurwitz problem of Section~\ref{sec:hurwitz-thm} under the dictionary
$(n,\delta,\alpha+2\delta)\mapsto(m,-\gamma,\alpha-2\gamma)$, with the lattice band factor instead of the Hurwitz one and without the
moved nodes. The analytic estimates of that problem (the outer scalar, the a-priori bound, the comparison at the hard edge with
Bessel order $\alpha$ and arbitrary real Szeg\H{o} exponent, the band-factor gauge and the order-$K$ small-norm expansion, steps
(H-b)--(H-e) of \S\ref{sec:H-proof}) are statements for a fixed index in which the index enters only as a real scale, through
$t_m$, the Bessel variable $\zeta=m^2f$ and the moments $\mu_{m,j}$; none of them uses its integrality. The range of $\gamma$ is the
range $-1<\delta\le0$ of those estimates, with $|\tilde\alpha|<1$. The band factor transforms homogeneously:
$\mu_{n,j}(n/m)^{jp}z^j=\mu_{m,j}z^j$. The expansion then follows as in the proof of Theorem~\ref{thm:N6}(i), for the decoupled model at index $m$.
\end{proof}

The proof of Lemma~\ref{lem:frame}(ii) is a transfer of estimates, written out step by step against \cite{PaperII} and the
Hurwitz analysis in the notes of the repository named in the data availability statement; we do not repeat it. There is no
circularity: Section~\ref{sec:hurwitz-thm} uses Theorem~\ref{thm:shift} only for the natural form (Theorem~\ref{thm:H}(ii)), not
for the estimates.

\begin{proof}[Proof of Theorem~\ref{thm:shift}]
\emph{Two expansions of one number.} Fix $p\in(1,2)\setminus\mathbb Q$, $\alpha\in(-\frac12,\frac12)$, $\gamma\in(0,\frac14)$. As $n\to\infty$
through the integers, Theorem~\ref{thm:N6}(i) and Lemma~\ref{lem:frame}(ii), the latter multiplied by $(n/m)^p$, both express
$n^pb_n/c$ as a finite sum $\sum_ec_en^{-e}$ plus $o(n^{-E})$, for any given $E$ (choose $K,B,J$ large); in the $m$-frame the factors
$(n/m)^{(1+|\br|_w)p}$ and $m^{-k}$ are convergent series in $1/n$. In both, the exponents are $e=k+|\br|_wp$, $k\in\N_0$, and since $p$
is irrational, the exponents $k+jp$ are pairwise distinct. Two finite sums that agree up to $o(n^{-E})$ have the same
coefficients below $E$; at $e=k\in\N_0$ only $\br=0$ contributes. Hence, coefficientwise,
$B(n+\gamma;\alpha-2\gamma,\alpha)=(1+\gamma/n)^pB(n;\alpha,\alpha)$.

\emph{From an open set to an identity.} The coefficient of $n^{-k}$ in the difference of the two sides is a rational function
$R_k(p,\alpha,\gamma)$, which vanishes on $((1,2)\setminus\mathbb Q)\times(-\frac12,\frac12)\times(0,\frac14)$. Writing $R_k=N_k/D_k$ with
polynomials, $N_k(p,\cdot,\cdot)$ vanishes on an open set of $(\alpha,\gamma)$ for infinitely many $p$, so each of its coefficients, a
polynomial in $p$, has infinitely many zeros; thus $R_k\equiv0$. Renaming $\alpha\mapsto\nu$ and putting $\gamma=(\nu-\beta)/2$ gives
$B(n_0+\frac{\nu-\beta}2;\beta,\nu)=((n_0+\frac{\nu-\beta}2)/n_0)^pB(n_0;\nu,\nu)$ for every $\beta$. With
$F_p(N;\nu):=(N/n_0)^pB(n_0;\nu,\nu)$, $n_0:=N-\frac{\nu+1}2$, this is~(i).

\emph{The zeta layers.} The same comparison identifies the coefficient of $n^{-k-Jp}$ for each weight $J=|\br|_w$:
$\sum_{|\br|_w=J}\hat\mu^{\br}[X_{\br}]_k=0$, where $\hat\mu_j=2\zeta(-jp)/j$ and $X_{\br}:=(n/m)^{(1+J)p}\cD_{\br}(m;\alpha-2\gamma,\alpha)-\cD_{\br}(n;\alpha,\alpha)$.
For $J=1$ only $\br=e_1$ occurs, and $\zeta(-p)\ne0$ for $p\in(1,2)$, so $X_{e_1}=0$. For $J=2$ the relation reads
$\zeta(-2p)[X_{e_2}]_k+4\zeta(-p)^2[X_{2e_1}]_k=0$; the brackets are rational in $p$ (after removing a common factor $c^2$), and if
$[X_{e_2}]_k\not\equiv0$, then $\zeta(-2p)/\zeta(-p)^2$ would be a rational function of $p$. It is not, since it vanishes at every odd
integer $p\ge3$ ($\zeta(-2p)=0$, $\zeta(-p)\ne0$). So $X_{e_2}=X_{2e_1}=0$, and the identity argument extends both to all
$(p,\alpha,\gamma)$; (ii) follows as~(i) did.
\end{proof}

For weights $J\ge3$ the separation would need the linear independence of the monomials $\prod_j\zeta(-jp)^{r_j}$, $|\br|_w=J$, over
$\mathbb Q(p)$, which we do not prove; the shift of the layers of weight $\ge3$ is open. Exact computer algebra confirms the frame
identity~\eqref{eq:FI} identically in $(p,\alpha,\beta,\gamma)$ to order $n^{-4}$, and $40$-digit evaluation of the recursion to order
$n^{-8}$ at $32$ points, including points with $p<1$ and with $\gamma$ outside the range of Lemma~\ref{lem:frame}; controls with a
half shift $\beta=\alpha-\gamma$, the wrong sign, no frame factor, no index shift, or the Bessel order moved with $\beta$, all fail
(scripts \texttt{m4\_shift\_exact}, \texttt{sigma\_frame\_exact}, \texttt{sigma\_frame\_num}). These checks confirm the algebraic
consequence; they do not test the analysis of Lemma~\ref{lem:frame}(ii).

\subsection{The zeta layer as a Toda flow}\label{sec:toda}
In the unscaled variable the band factor $\exp(-\sum_k\mu_{n,k}z^k)$, $z=n^px$, is $\exp(\sum_kt_kx^k)$ with
\begin{equation}\label{eq:todatimes}
  t_k:=-\frac{2\zeta(-kp)}k ,
\end{equation}
a weight factor independent of $n$; these are the times of the Toda hierarchy. For a weight $w_t=e^{\sum_jt_jx^j}w$ with monic
orthogonal polynomials $\pi_n$, $h_n=\int\pi_n^2w_t$ and Jacobi matrix $\mathcal J$ (orthonormal basis), $\partial_{t_j}\pi_n$ has degree
$<n$, so $\partial_{t_j}h_n=(\mathcal J^j)_{nn}h_n$ and, since $b_n^2=h_n/h_{n-1}$,
\begin{equation}\label{eq:toda}
  \partial_{t_j}\log b_n^2=(\mathcal J^j)_{nn}-(\mathcal J^j)_{n-1,n-1},\qquad \partial_{t_1}a_n=b_{n+1}^2-b_n^2 .
\end{equation}

\begin{lemma}[Toda hierarchy for tilts]\label{lem:todatilt}
Let $\mu$ be a positive measure on $\R$ with compact support and infinitely many support points, $P_t(x):=\sum_{j\le J}t_jx^j$,
$t\in\R^J$, and let $a_n(t),b_n(t),h_n(t)$ and $\mathcal J(t)$ belong to $e^{P_t}\mu$. Then $a_n,b_n,h_n$ are real-analytic in $t$, and
for every $j\le J$ and $n\ge0$ ($n\ge1$ in the second identity)
\[
  \partial_{t_j}\log h_n=(\mathcal J^j)_{nn},\qquad
  \partial_{t_j}b_n=\tfrac12b_n\bigl[(\mathcal J^j)_{nn}-(\mathcal J^j)_{n-1,n-1}\bigr],\qquad
  \partial_{t_j}a_n=F_j(n)-F_j(n-1),
\]
with $F_j(n):=b_{n+1}(\mathcal J^j)_{n+1,n}$ and $F_j(-1):=0$.
\end{lemma}

\begin{proof}
The moments $\int x^ke^{P_t}\diff\mu$ are entire in $t$ and the Hankel determinants are positive, so the coefficients of the monic
orthogonal polynomials $\pi_n$, and $h_n=\int\pi_n^2e^{P_t}\diff\mu$, are real-analytic in $t$. Since $\partial_{t_j}\pi_n$ has degree $<n$,
differentiating $\int\pi_n\pi_me^{P_t}\diff\mu=0$ ($m<n$) gives $\partial_{t_j}\pi_n=-\sum_{m<n}h_m^{-1}\bigl(\int x^j\pi_n\pi_me^{P_t}\diff\mu\bigr)\pi_m$,
and $\partial_{t_j}h_n=\int x^j\pi_n^2e^{P_t}\diff\mu=h_n(\mathcal J^j)_{nn}$; with $b_n^2=h_n/h_{n-1}$ this is the second identity. If
$\pi_n=x^n+\gamma_nx^{n-1}+\dotsb$, then $a_n=\gamma_n-\gamma_{n+1}$, and the formula for $\partial_{t_j}\pi_n$ gives
$\partial_{t_j}\gamma_n=-(h_n/h_{n-1})^{1/2}(\mathcal J^j)_{n,n-1}=-F_j(n-1)$.
\end{proof}

The Toda hierarchy for exponential tilts is classical (Moser, Flaschka; see the survey \cite{VanAsscheToda}, Lemma~2.1 and
Theorem~2.4 there); we include the short proof for the identities in the form used below. For $j=1$ these are the equations $\partial_t\log b_n^2=a_n-a_{n-1}$, $\partial_ta_n=b_{n+1}^2-b_n^2$ of the Toda lattice. To state the
consequence for the layers, let $\cA_{\br}(n)$ be the analogues of $\cD_{\br}$ for the diagonal coefficient, obtained from the second
moment $\mathsf R_2=\Res_0(z\mathsf F)+\Res_{4c}(z\mathsf F)$ of the same residue recursion and the extraction
$n^pa_n=(Q_2)_{12}/(Q_1)_{12}-(Q_1)_{22}$ of \cite{PaperII}, and $A:=\cA_0$; the $n^0$ coefficient of $\cA_K(n;\mu)$ is $1$ for every
$\mu$. Substitute $\mu_j=-u_jn^{-jp}$ with times $u_j$ that do not depend on $n$, and put
\[
  \hat b:=\sum_{\br}(-u)^{\br}n^{-|\br|_wp}\cD_{\br}(n),\qquad\hat a:=\sum_{\br}(-u)^{\br}n^{-|\br|_wp}\cA_{\br}(n),
\]
formal series in $1/n$ and the $u_j$. Let $\hat{\mathcal J}$ be the tridiagonal matrix indexed by $i\in\mathbb Z$ with
$\hat{\mathcal J}_{ii}=2(1+\frac in)^{-p}\hat a(n+i)$ and $\hat{\mathcal J}_{i,i-1}=\hat{\mathcal J}_{i-1,i}=\hat b_i:=(1+\frac in)^{-p}\hat b(n+i)$,
where $n\mapsto n+i$ acts at fixed $u$ and is re-expanded in $1/n$; these are the formal values of $n^pa_{n+i}/c$ and $n^pb_{n+i}/c$.
Using the Toda equations, with shifted indices re-expanded in $1/n$, to generate the coefficients of an asymptotic expansion is the
method of Ercolani, McLaughlin and Pierce \cite[Section~4]{EMP} for varying weights $e^{-N(x^2/2+tx^{2\nu})}$. Here there is no large
parameter other than $n$, the measure is fixed, the times enter as $u_jn^{-jp}$ with base values~\eqref{eq:todatimes}, and the
identities are transferred from the exact flows by a category argument rather than derived from a string equation.

\begin{theorem}[Toda structure]\label{thm:toda}
Let $p>0$, $p\ne1$ and $1<s<1+2p$.
\begin{enumerate}
\item[(i)] \emph{Tilted lattices.} For every $J$ and $\tau\in\R^J$, the measure $e^{\sum_{j\le J}\tau_jx^j}\sigma_{p,s}$ satisfies
Theorem~\ref{thm:N6}(i) with the same layers $\cD_{\br,K}$ and with $\mu_{n,j}$ replaced by $\mu_{n,j}-\tau_jn^{-jp}$, and
$n^pa_n/(2c)=\sum_{|\br|\le B}\mu_n^{\br}\cA_{\br,K}(n)+O(n^{-K-1}+n^{-1-(B+1)p}+n^{-1-(J+1)p})$, likewise with the shifted $\mu_n$.
\item[(ii)] \emph{Formal Toda hierarchy.} For every $j$, as identities of formal series,
\[
  \partial_{u_j}\hat b=\tfrac12c^jn^{-jp}\,\hat b\,\bigl[(\hat{\mathcal J}^j)_{00}-(\hat{\mathcal J}^j)_{-1,-1}\bigr],\qquad
  \partial_{u_j}(2\hat a)=c^jn^{-jp}\bigl[\hat b_1(\hat{\mathcal J}^j)_{10}-\hat b_0(\hat{\mathcal J}^j)_{0,-1}\bigr].
\]
\item[(iii)] The identities of (ii) hold for the layers of the decoupled model identically in $(p,\beta,\nu)$, and they determine every
$\cD_{\br}$ and $\cA_{\br}$ from the two integer layers $B$ and $A$.
\end{enumerate}
In particular, since $\partial_{\mu_j}=-n^{jp}\partial_{u_j}$, at $\mu=0$
\begin{gather*}
  \partial_{\mu_1}(n^pb_n)=-\tfrac12n^p(a_n-a_{n-1})\,n^pb_n,\qquad
  \tfrac12\partial_{\mu_1}^2(n^pb_n)=\tfrac12n^{2p}\partial_{t_1}^2(n^pb_n),\\
  \partial_{\mu_2}(n^pb_n)=-\tfrac12n^{2p}\bigl[(\mathcal J^2)_{nn}-(\mathcal J^2)_{n-1,n-1}\bigr]n^pb_n,\qquad
  (\mathcal J^2)_{nn}=a_n^2+b_n^2+b_{n+1}^2,
\end{gather*}
with $\partial_{t_1}^2\log b_n^2=b_{n+1}^2-2b_n^2+b_{n-1}^2$, the shifted entries being those of the formal series.
\end{theorem}

The relation $\partial_{\mu_j}=-n^{jp}\partial_{u_j}$ is the chain rule for the substitution $\mu_j=-u_jn^{-jp}$. Its content is that the
shift $n\mapsto n+i$ acts at fixed unscaled time, $\mu_j(n+i)=\mu_j(n)(1+i/n)^{-jp}$; treating $\mu$ as independent of $n$ gives the
right first-order layers and wrong second-order ones.

\begin{proof}[Outline]
(i) The tilted lattice is $\sigma_{p,s,\ell}$ of \cite{PaperII} with $\ell(k)=\exp\Lambda(\log k)$, $\Lambda(\xi)=\sum_j\tau_je^{-jp\xi}$. This
$\Lambda$ is entire, real on $\R$ and bounded on $\{-c_0\le\Re\xi\le0\}$, and $\sup_{\Re\xi\ge T}|\Lambda'|\le p\sum_jj|\tau_j|e^{-jpT}\to0$, so $\ell$
is analytic of every angle \cite[Definition~1.1]{PaperII}, of either sign and any size of $\tau$, and the frame of \cite{PaperII} applies. Its
band weight is that of the lattice times $r_n(z)=\exp(\sum_j\tau_jn^{-jp}(z^j-1))$, an entire polynomial tilt; by Lemma~\ref{lem:gauge}
it changes the band-factor parameters to $\mu_n-\tau n^{-p\cdot}$ exactly and nothing else. Every estimate behind Theorem~\ref{thm:N6}
uses the weight only through $W_0$, its Szeg\H{o} function and bounds on them, which change by factors depending on $\tau$ only; at
the hard edge $r_n^{1/2}$ is entire, so the jumps of the local parametrix are unchanged. All constants depend on $\tau$; no uniformity in
$\tau$ is needed. The statement for $a_n$ is the same residue recursion with the second moment $\mathsf R_2$.
(ii) Write $u=\tau+t$ with $t_j=-2\zeta(-jp)/j$, so that $u$ ranges over $\R^J$, and fix an order $L$. By (i), $b_m$ and $a_m$ of the tilted
measures equal polynomials in $u$ up to $O_u(m^{-p-L})$, with constants depending on $u$. These errors are continuous in $u$, and by the
Baire category theorem there is an open ball on which they are bounded uniformly for all large $m$. On this ball integrate
Lemma~\ref{lem:todatilt} along $u_j$ from $u$ to $u+he_j$ and insert the expansions: both sides are polynomials in $(u,h)$ up to
$O(n^{-p-L})$, uniformly, so each coefficient is $O(n^{-p-L})$. The coefficient of $u^\gamma h^m$ is $n^{-p(1+|\gamma|_w+mj)}$ times a
function analytic in $1/n$; as $L$ is arbitrary, the integrated identity holds as an identity of formal series, and its coefficient of $h$ is
(ii). No asymptotic expansion is differentiated term by term.
(iii) The same argument for the Hurwitz lattices of Section~\ref{sec:hurwitz-thm}, with Theorem~\ref{thm:H}(i) in place of
Theorem~\ref{thm:N6}(i), gives (ii) for the layers of $M(\alpha+2\delta,\alpha)$, $-1<\delta<0$. The layer coefficients are rational in
$(p,\beta,\nu)$ (with $c=1$; the model is homogeneous under $z\mapsto z/c$, $\mu_j\mapsto\mu_jc^j$), and $\{(\alpha+2\delta,\alpha)\}$ is open,
so the identities hold identically. The coefficient of $u^{\gamma+e_j}$ on the left of (ii) is $\gamma_j+1$ times a new layer, while the
right side involves only layers of total order $\le|\gamma|$, since the shifts preserve the degree in $u$.
\end{proof}

The details of (i), the transfer of the estimates of Section~\ref{sec:natural} to the tilted measures row by row, are written out in the notes of the
repository (Appendix~\ref{app:repro}). Exact computer algebra confirms both equations of (ii) on the layers of the residue recursion,
identically in $(p,\alpha)$ on the lattice line and in $(p,\alpha,\beta)$ for the decoupled model, for $j\le3$ and to second order in
the times (third order at rational points, including $p<1$); the variant in which the shift acts at fixed scaled time $\mu$ fails at
second order. On the finite truncations of the lattice and Hurwitz measures Lemma~\ref{lem:todatilt} holds exactly in rational
arithmetic (scripts \texttt{toda\_measures}, \texttt{toda\_layers}). In this reading the zeta layer is not new data: it is a fixed, $n$-independent Toda flow of the two integer
layers, with times~\eqref{eq:todatimes} given by the zeta values.

\subsection{The index-shift law}\label{sec:indexshift}
Theorem~\ref{thm:toda} determines the leading term of every layer in closed form, by the genus-zero mechanism of \cite{Ercolani}. Let the decoupled model carry band-factor
parameters $\mu_1,\dots,\mu_r$ with exponents $j_1,\dots,j_r\in\N$, write $\mu_a=\hat\mu_an^{-pj_a}$ and
$\hat\varepsilon_a:=\frac{pj_a}2\binom{2j_a}{j_a}c^{j_a}\hat\mu_a$ (for the lattice, $\hat\varepsilon_j=\varepsilon_j$), and expand
$n^pb(n;\beta,\nu;\mu)/c=\sum_kn^{-k}b_k(\mu)$.

\begin{theorem}[index-shift law]\label{thm:indexshift}
Let $p\ne1$. Identically in $(p,\beta,\nu)$:
\begin{enumerate}
\item[(i)] $b_k$ is a polynomial of degree at most $k$ in $\mu$; that is, a layer $\mu^{\br}$ starts at $n^{-|\br|}$;
\item[(ii)] the part of $b_k$ of degree exactly $k$ is the $n^{-k}$ coefficient of $(n/x)^p$, where $x$ is the formal solution of
\[
  x=n-\frac1p\sum_a\hat\varepsilon_a\,x^{-pj_a}.
\]
\end{enumerate}
Equivalently, in natural form the layer $\mu^{\br}$ is
$(n/N)^p\hat\varepsilon^{\br}N^{-|\br|-S}\bigl[(p+1+S)_{|\br|-1}/(\br!\,p^{|\br|-1})+O(N^{-1})\bigr]$, $S:=p\sum_ar_aj_a$, with the
rising factorial $(y)_m=y(y+1)\cdots(y+m-1)$.
\end{theorem}

For one time this gives $\varepsilon_jN^{-1-jp}$ and $(3p+1)/(2p)$ for $\mu_1^2$ (Theorem~\ref{thm:N6}(ii)); the exact linearity of the
$n^{-1}$ coefficient in $\mu$ (Lemma~\ref{lem:E37}) is the case $k=1$ of (i). So, at leading order in $1/n$ at each order in the times,
\emph{the band-factor times act as a shift of the index}, $n\mapsto x$; the shift is implicit and depends on the local value of $b$.

\begin{proof}[Outline]
Grade the term $\hat\mu^{\br}n^{-k-S}$ of a layer by its \emph{defect} $k-|\br|$. Defects add under products, a shift $n\mapsto n+i$
does not lower the defect, and a difference $(1-S_{-1})$, $S_{-1}\colon n\mapsto n-1$, raises it by one and acts on the lowest-defect part
as $\partial_n$. Both equations of Theorem~\ref{thm:toda}(ii) are in this conservation form, so by induction on the order in the times
every term has defect $\ge0$, which is (i), and the defect-$0$ parts $(a^0,b^0)$ of $(\hat a,\hat b)$, written as $a_n,b_n$, solve the
dispersionless Toda hierarchy
\[
  \partial_{t_j}b^0=\tfrac12b^0\partial_nH_j(a^0,b^0),\qquad\partial_{t_j}a^0=\partial_nG_j(a^0,b^0),
\]
where $H_j(a,b)$ and $G_j(a,b)$ are $(\mathcal J^j)_{nn}$ and $b_{n+1}(\mathcal J^j)_{n+1,n}$ for the constant Jacobi matrix with entries
$a,b$; the solution is unique given the initial data. These are $a^0=2cn^{-p}$, $b^0=cn^{-p}$ (the $n^0$ coefficients of $A$ and $B$
are $1$), which lie on the line $a=2b$ of a hard edge at $0$. Since $(1+x)^{2j}=(1+2x+x^2)^j$, on this line
$H_j=\binom{2j}jb^j$ and $G_j=\frac j{j+1}\binom{2j}jb^{j+1}$, so $\partial_t(a^0-2b^0)=0$: the line is invariant, and on it every flow is
the Hopf equation $\partial_{t_j}b=V_j(b)\partial_nb$, $V_j(b)=\frac j2\binom{2j}jb^j$. Its solution by characteristics is $b=cx^{-p}$
with $x=n+\sum_jt_jV_j(cx^{-p})$, which is (ii) for $t=-\hat\mu$; Lagrange inversion of this equation gives the coefficients. For $\beta\ne\nu$ the same argument runs on the
identities of Theorem~\ref{thm:toda}(iii).
\end{proof}

The closed form agrees with every exact layer computed by the residue recursion: all monomials to fourth order for symbolic
exponents and to sixth order at rational parameters, with $\beta\ne\nu$ included, and, beyond the leading term, the Toda flow of the two
integer layers reproduces the computed dressings $g_j$ and $h_2$ of Theorem~\ref{thm:N6}(ii) (scripts \texttt{shift\_toda},
\texttt{shift\_existing}, \texttt{shift\_exact}, \texttt{shift\_dressing}). A ``dressed'' index shift $x=N-\frac1p\varepsilon x^{-p}(1+\gamma x^{-2})$
does not reproduce $h_2$: the corrections to the leading term are dispersive, not a reparametrisation of the index.

At leading order the mechanism is the genus-zero one of random-matrix map enumeration: the continuum limit of the Toda equations is an
inviscid Burgers (Hopf) equation with power-law speed, solved implicitly by characteristics, and the coefficients of the implicit
solution are obtained by Lagrange inversion (\cite[Section~2.3]{Ercolani}; \cite{EMP}; for Lagrange inversion in map counting, e.g.
\cite[Section~26.4]{Bouttier}). The endpoints of the support are Riemann invariants of the dispersionless Toda hierarchy
(cf.~\cite{KonopelchenkoMAM}); the hard edge fixes one of them at $0$, which is the invariance of the line $a=2b$. What differs here is the
initial profile $b=cn^{-p}$ of a fixed measure accumulating at a hard edge (instead of the Gaussian $b^2=n/N$), the reduction to the
line $a=2b$, arbitrary real and pole exponents, and zeta-valued times; the mechanism itself is not new.

\begin{remark}[non-integer exponents]\label{rem:indexshift-poles}
The layer coefficients of the decoupled model are rational functions of the exponents $j_a$, with denominators that are products of
$2j_a-(2i+1)$; a rational function vanishing at all integer points vanishes identically, so the identities of
Theorem~\ref{thm:indexshift} hold for the formal model at every exponent, in particular at the pole exponents $q_m=m-\frac1p$ of
Theorem~\ref{thm:kappa}. Theorem~\ref{thm:kappa} uses only the first-order layer at $j=q_1$. That the products of pole terms of the
shifted Hurwitz lattices are the continued layers of the model at second and higher order would need an extension of the proof of
Theorem~\ref{thm:kappa} to all orders in the poles, which we do not give here; for those measures the leading coefficients of the pole
products given by Theorem~\ref{thm:indexshift} are therefore conjectural.
\end{remark}

\begin{proposition}[the diagonal layer to order four]\label{prop:diag4}
Let $M:=n+\frac\beta2+1$. Identically in $(p,\beta,\nu)$, the diagonal integer layer satisfies
\[
  \Bigl(\frac Mn\Bigr)^{\!p}A(n;\beta,\nu)=1+\frac{a_2}{M^2}+\frac{a_4}{M^4}+O(M^{-5}),\qquad a_2=-\frac{p+1}{24}+\frac{(p+1)(p+2)\nu^2}8,
\]
\[
  a_4=\frac{(p+1)(p+3)\bigl[45p^2(p+2)^2\nu^4+30p(p^2+6p+8)\nu^2-21p^2-38p+4\bigr]}{5760\,p\,(3p+2)},
\]
with no $M^{-1}$ and no $M^{-3}$ term and no dependence on $\beta$ to this order. For the lattice ($\beta=\nu=\alpha$), $p\ne1$ and
$1<s<1+2p$, these are the coefficients of the integer layer of $n^pa_n/(2c)$ in Theorem~\ref{thm:toda}(i).
\end{proposition}

\begin{proof}
Exact computer algebra on the second moment of the residue recursion, as in Lemma~\ref{lem:N6d}, identically in $(p,\nu,\beta)$ to order
$n^{-4}$; the $n^{-1}$ coefficient is $-p(\beta+2)/2$, the first term of $(n/M)^p$ (script \texttt{toda\_layers}).
\end{proof}

At the Bessel point $(p,\nu)=(2,-\frac12)$, $a_2=\frac14$ and $a_4=\frac1{16}$, the first coefficients of $(1-1/(4M^2))^{-1}$, the
diagonal of the Lommel measures. The coefficient $a_2$ also follows from the first equation of Theorem~\ref{thm:toda}(ii) and the
first-order layer $\mu_1$, without computing $A$. Beyond order four we have only numerical evidence:

\begin{conjecture}[diagonal shift]\label{conj:diag}
For the decoupled model, $(M/n)^pA(n;\beta,\nu)=G_p(M;\nu^2)$ with $G_p$ even in $1/M$ and independent of $\beta$ to all orders.
At $(p,\nu)=(2,-\frac12)$, $G_2=(1-1/(4M^2))^{-1}$.
\end{conjecture}

The evidence beyond Proposition~\ref{prop:diag4} is numerical: the recursion at $40$ digits to order $8$, with $\beta$ and $\nu$
decoupled ($19$ checks), and $a_2$ reconstructed from $p\in\{\frac12,\dots,6\}$ before it was derived exactly (scripts \texttt{m4\_toda},
\texttt{m4\_diag}). So $b_n$ is centred at $N$ and $a_n$ at $M=N+\frac12$, the arrangement of
\cite[Theorem~5.2]{BleherIts} for varying weights, and Theorem~\ref{thm:reflect} corresponds to the reversal symmetry
$b_n\leftrightarrow b_{-n}$, $a_n\leftrightarrow a_{-n-1}$ of the Toda lattice. The proof of Theorem~\ref{thm:shift} should apply to
$a_n$ through the same invariance of the extraction once the order-$K$ expansion of $a_n$ is written for $M(\beta,\nu)$.

\begin{remark}[the parity of the zeta layers]\label{rem:parity}
By Theorem~\ref{thm:toda}, the parity statements of Conjecture~\ref{conj:S1Nprime} reduce to the integer layers. If $B$ is
even in $N$ and $A$ is even in $M=N+\frac12$, then $a_n-a_{n-1}=A(N+\frac12)-A(N-\frac12)$ (in natural form) is odd in $N$, so each
first-order zeta term is $\varepsilon_jN^{-1-jp}$ times an even series; second differences $b_{n+1}^2-2b_n^2+b_{n-1}^2$ are even, so the
$\varepsilon_1^2$ layer is $N^{-2-2p}$ times an even series, as found in Theorem~\ref{thm:N6}(ii). Higher Toda derivatives are built
from such differences and products, and we expect the full monomial structure to follow; we have not written this out in general,
and exact computer algebra confirms it to third order in the times for generic even integer layers (script \texttt{shift\_dressing}). By
Corollary~\ref{cor:even} the evenness of $B$ is proved; that of $A$ is Proposition~\ref{prop:diag4} to order $M^{-4}$ and
Conjecture~\ref{conj:diag} beyond. The $1/(2j-1)$ in $g_j$ has a
heuristic reading in the same picture: for a slowly varying Jacobi matrix, $(\mathcal J^j)_{nn}$ is approximately
$\frac1{2\pi}\int_0^{2\pi}(a+2b\cos\theta)^j\diff\theta$ (the dispersionless symbol), whose first variation off the point $a=2b$ involves
$\binom{2j-2}{j-1}/\binom{2j}j=j/(2(2j-1))$.
\end{remark}

% Section 7: Hurwitz lattices (from research/G1/RH/spectral/JZC_M3a_H.md and its chain; the shifted family from
% JZC_W_outer.md, JZC_W_kappa.md, JZC_W_BT.md).
\section{Hurwitz lattices}\label{sec:hurwitz-thm}

For the Hurwitz lattice $\lambda_k=(k+\delta)^p$ the node zeta function is $Z_\Lambda(w)=\zeta(pw,a)$, $a:=1+\delta$, so the dictionary
of \S\ref{sec:dictionary} predicts the natural index $n+\frac\alpha2-\zeta(0,a)=n+\frac\alpha2+\frac12+\delta$ and the zeta terms with
$\zeta(-jp,a)$ in place of $\zeta(-jp)$. This section proves it. The mechanism is an index transfer: the offset $\delta$ does not
change the Bessel order of the hard edge, it changes only the Szeg\H{o} exponent, $\alpha\mapsto\beta'=\alpha+2\delta$; the Hurwitz
problem is the decoupled model $M(\alpha+2\delta,\alpha)$ of \S\ref{sec:decoupled}, and Theorem~\ref{thm:shift} converts the shifted
Szeg\H{o} exponent into the shifted natural index.

\subsection{Statement}\label{sec:H-statement}
\begin{theorem}[Hurwitz lattices]\label{thm:H}
Let $p>0$, $p\ne1$, $-1<\delta\le0$ (with $\delta<0$ if $p<1$), $a:=1+\delta$, $1<s<1+2p$, and
\[
  \sigma^\delta:=\sum_{k\ge1}(k+\delta)^{-s}\,\delta_{(k+\delta)^{-p}} ,
\]
the measure $\sigma_\Lambda$ of~\eqref{eq:sigmaL} for $\lambda_k=(k+\delta)^p$. Put $\beta':=\alpha+2\delta$,
$\mu^\delta_{n,j}:=2\zeta(-jp,a)/(jn^{jp})$, $N_\delta:=n+\frac\alpha2+\frac12+\delta$ and $\varepsilon_j(\delta):=\binom{2j}jp\,c^j\zeta(-jp,a)$.
\begin{enumerate}
\item[(i)] \emph{Reduction.} For all $K,J,B\ge1$, the recurrence coefficients of $\sigma^\delta$ satisfy
\[
  \frac{n^pb_n}c=\sum_{|\br|\le B}(\mu^\delta_n)^{\br}\,\cD_{\br,K}(n;\beta',\alpha)+O\bigl(n^{-K-1}+n^{-2-(B+1)p}+n^{-1-(J+1)p}\bigr),
\]
with the layers $\cD_{\br,K}(n;\beta,\nu)$ of the decoupled model~\eqref{eq:decoupled}, truncated at order $n^{-K}$.
\item[(ii)] \emph{Natural form.} In $N_\delta$, the integer layer is that of the lattice to all orders,
$\cD_0(n;\beta',\alpha)=(n/N_\delta)^pF_p(N_\delta;\alpha)$, and the layers $e_1$, $e_2$, $2e_1$ are those of the lattice to all orders, with
$\zeta(-jp)$ replaced by $\zeta(-jp,a)$. For every $j\ge1$ the layer $e_j$ is $(n/N_\delta)^p\varepsilon_j(\delta)N_\delta^{-1-jp}(1+g_jN_\delta^{-2}+O(N_\delta^{-4}))$.
In particular, expansion~\eqref{eq:thmA} of Theorem~\ref{thm:A} holds for $\sigma^\delta$ with $N$ replaced by $N_\delta$ and
$\varepsilon_j$ by $\varepsilon_j(\delta)$, with the same $f_2,f_4,f_6,g_j,g_{1,4}$ and the same remainder.
\end{enumerate}
\end{theorem}

At $\delta=0$ this is Theorem~\ref{thm:N6}. The $1/n$ coefficient is $-p(\alpha+1)/2-p\delta=-\frac{s-1}2-p\delta$, the first prediction tested
in Table~\ref{tab:hurwitz}, and the $n^{-1-p}$ coefficient is $\varepsilon_1(\delta)=2pc\,\zeta(-p,a)$. The layers of weight $\ge3$ ($\mu_3$,
$\mu_1\mu_2$, $\mu_1^3$) are given by~(i), but their natural form beyond the orders of Lemma~\ref{lem:N6d} is open, as for the lattice.
Nothing is claimed for $\delta>0$: then $t=\delta$, in the index variable $t=nz^{-1/p}$ of \cite{PaperII}, is a zero of the
interpolating function that carries no mass (a ``phantom node''), and the pole removal of \cite[Section~5]{PaperII} fails there.
Nothing is claimed at $p=1$ with $\delta\ne0$.

\emph{The constants.} The proof produces the band-factor moments as Fourier sums (step (H-c) below),
\begin{equation}\label{eq:fourier}
  m_k(\delta)=-\frac p\pi\Gamma(kp)(2\pi n)^{-kp}S_k(\delta),\qquad S_k(\delta):=\sum_{m\ge1}\frac{\sin(\pi kp/2-2\pi m\delta)}{m^{1+kp}} ,
\end{equation}
and $\mu^\delta_{n,k}=2m_k(\delta)$. Their identification with Hurwitz zeta values is Hurwitz's formula, which we use in the form
stated in \cite[Theorem~1.1]{DixitKumar} (where it is quoted from Titchmarsh): \emph{for $0<a\le1$ and $\Re s<0$,}
\[
  \zeta(s,a)=\frac{2\Gamma(1-s)}{(2\pi)^{1-s}}\Bigl\{\sin\frac{\pi s}2\sum_{m\ge1}\frac{\cos(2\pi ma)}{m^{1-s}}
  +\cos\frac{\pi s}2\sum_{m\ge1}\frac{\sin(2\pi ma)}{m^{1-s}}\Bigr\}.
\]
We apply it at $s=-kp$, where $\Re s=-kp<0$, and $a=1+\delta\in(0,1]$. Since $\cos(2\pi ma)=\cos(2\pi m\delta)$,
$\sin(2\pi ma)=\sin(2\pi m\delta)$ and $\sin(-\pi kp/2)=-\sin(\pi kp/2)$, the bracket is $-S_k(\delta)$, and
$2\Gamma(1+kp)(2\pi)^{-1-kp}=\frac{kp}\pi\Gamma(kp)(2\pi)^{-kp}$. Hence $\zeta(-kp,a)=-\frac{kp}\pi\Gamma(kp)(2\pi)^{-kp}S_k(\delta)=kn^{kp}m_k(\delta)$,
which is the normalisation of Theorem~\ref{thm:H}. (The formula also agrees with direct evaluation of $\zeta(-kp,a)$ to $5\cdot10^{-42}$
at $48$ parameter points, script \texttt{m3a\_hurwitz\_formula}.)

\subsection{The index transfer and the \texorpdfstring{$\alpha$}{alpha}-frame}\label{sec:H-mechanism}
In the variable $t=nz^{-1/p}$ the Hurwitz nodes sit at $t\in a+\N_0$, and the mass at a node is still $n^{p}t^{-s}$ up to the
constant of \cite{PaperII}; so the weight on the band, its power $z^\alpha$ and the Bessel order are unchanged, and only the phase
changes, $\theta=\pi t\mapsto\theta_\delta=\pi(t-\delta)$. Interpolation is done with $\sin\pi(t-\delta)$ and $1/\Gamma(t-\delta)$, so the
outer scalar of \cite{PaperII} is built from
\[
  j_\delta(t):=\operatorname{Ln}\Gamma(t-\delta)-(t-\delta-\tfrac12)\log t+t-\tfrac12\log2\pi
  =j(t-\delta)+(t-\delta-\tfrac12)\operatorname{Log}(1-\delta/t)+\delta ,
\]
where $j=j_0$ is the Binet function of \cite{PaperII}. Nothing changes at the hard edge ($t\to\infty$), but at $z=\infty$ ($t\to0$) the
outer scalar $F_\delta$ now behaves like $C_\delta(-z)^{\delta}$ with $C_\delta\ne0$. Normalising by the gauge
\[
  F^\natural:=F_\delta\,(-z)^{-\delta}
\]
restores the normalisation at infinity, at fixed degree, scale and $g$-function, and moves the power $(-z)^{-2\delta}$ into the
band weight: the Szeg\H{o} exponent becomes $\beta'=\alpha+2\delta=\alpha-2Z_\Lambda(0)-1$. On the lips in the hard-edge disc the
gauge leaves the phase $e^{\pm2\pi i\delta}$, so the $(2,1)$ entry of the lip jump is $\omega^\natural e^{-4\zeta^{1/2}}$ with
$\omega^\natural=\omega^{(\beta')}e^{\pm2\pi i\delta}$, whose boundary values on the band are $e^{\mp i\pi\beta'}e^{\pm2\pi i\delta}=e^{\mp i\pi\alpha}$.
This is the data of a Bessel model of order $\alpha$ glued to the Szeg\H{o} function of $x^{\beta'}$: the decoupled model
$M(\beta',\alpha)$.

The hard-edge comparison must therefore be made in the \emph{$\alpha$-frame}: the local parametrix is the Bessel model of
order $\alpha$, and $\beta'$ enters only through the quotient $\rho^\natural=\rho^{(\alpha)}/G$, with
$G:=(-z)^{-\delta}D_{x^{2\delta}}$, $D$ the Szeg\H{o} function on $[0,4c]$. By the formula $D_{x^{b}}(z)^2=(4cz/\Upsilon(z))^{b}$ of \cite{PaperII},
$|G|=|4c/\Upsilon|^\delta$ with $\Upsilon=(z^{1/2}+(z-4c)^{1/2})^2$, $\Upsilon(0)=-4c$, so $G$ is bounded above and below near $0$
uniformly in $n$, for every real $\beta'$. A comparison in a $\beta'$-frame (Bessel order $\beta'$) would fail: the a-priori bound
$|z|^{-|\beta'|/2}$ is not square integrable at $0$ once $|\beta'|\ge1$, which happens for $\delta$ near $-1$.

\subsection{Outline of the proof}\label{sec:H-proof}
The proof follows \cite[Sections~5--10]{PaperII} and Section~\ref{sec:natural}, in six steps.
\begin{enumerate}
\item[(H-a)] \emph{Interpolation.} The lemmas of \cite[Section~5]{PaperII} with $t\mapsto t-\delta$: node masses as residues, and the top
node $X_1=(n/a)^p>n^p$, so that the expansion at infinity holds for $|z|>X_1$ and the contour and the far-arc bounds are taken
with constants that depend also on $a$.
\item[(H-b)] \emph{Outer scalar and a-priori bound.} By the shift identity for $j_\delta$, every bound on $j_\delta$ reduces to the Binet
bound of \cite{PaperII} on $j$; this gives the sector bounds of the outer scalar ($p>1$: Binet cut; $p<1$: seam version) and its
behaviour $C_\delta(-z)^\delta$ at infinity. The a-priori bound $|Q|\le C_{n,\eta}|z|^{-|\alpha|/2-\eta}$ near $0$ and the $L^2$ boundary
values follow by the column factorisation of \cite[Appendix~C]{PaperII}, with the exponent of the lattice, $|\alpha|/2$, because $G$
is bounded above and below.
\item[(H-c)] \emph{Band factor.} The jump densities of the Binet cut ($p>1$) or of the seams ($p<1$) are
$2\Im\operatorname{Log}(1-e^{-2\pi i(t_+-\delta)})$; their moments are~\eqref{eq:fourier}, by \cite[(5.9.1), (5.2.1), (25.2.1)]{DLMF} and
absolute convergence, with the Fubini bounds of Lemma~\ref{lem:Fprime}, uniformly in $\delta$. The gauge of
Lemma~\ref{lem:gauge} then holds with $\mu^\delta_n$ on both circles.
\item[(H-d)] \emph{Index transfer.} With $F^\natural$ in place of the outer scalar, the jumps are those of $M(\beta',\alpha)$ with the
Hurwitz band factor, as in \S\ref{sec:H-mechanism}. This is exact algebra on the jumps of \cite[Section~5]{PaperII}.
\item[(H-e)] \emph{Hard edge and small norm.} In the $\alpha$-frame the comparison of \cite[Section~7]{PaperII} holds verbatim, and the
order-$K$ small-norm expansion of Proposition~\ref{prop:N6c} holds with the circle jumps $\Delta_k(\beta',\alpha)$; Lemmas~\ref{lem:N6a}
and~\ref{lem:N6b} hold for $M(\beta,\nu)$, since $\beta$ enters only through $\rho^2=\phi^{2(\beta+1)}$ and $\phi^{-2(\beta+1)}$, bounded above
and below on the circles. For $p<1$ the glued hard-edge parametrix of \cite[Section~10]{PaperII} is rebuilt with the gauged weight,
using a Cauchy-difference bound on the seam band factor that is uniform in $\delta$. The step to~(i) also needs the linear
response of Lemma~\ref{lem:E37} for $M(\beta,\nu)$; it holds with $d_1=-p(\beta+1)/2$, by the same computation with the residue at $0$
depending on $\nu$ only and the shift at $4c$ on $\beta$ only (verified by exact computer algebra, script \texttt{m3a\_e37\_decoupled}).
\item[(H-f)] \emph{Natural form.} Exact computer algebra for $M(\beta,\nu)$, identically in $(p,\nu,\beta)$: the integer layer is
$\beta$-free in $N=n+\frac{\beta+1}2$ to order $N^{-4}$, and the first-order zeta layer, for $j=1,2,3$ and for symbolic $j$, is
$\varepsilon_jN^{-1-jp}(1+g_jN^{-2}+O(N^{-4}))$ with $\beta$-free $g_j=A_j+B_j\nu^2$ (scripts \texttt{m4\_shift\_exact}, \texttt{m3a\_g5\_exact}).
\end{enumerate}
Part~(i) is (H-a)--(H-e) followed by the proof of Theorem~\ref{thm:N6}(i). Part~(ii) is~(i) with~(H-f) and Theorem~\ref{thm:shift}
at $(\beta,\nu)=(\alpha+2\delta,\alpha)$: by Theorem~\ref{thm:shift}(i) the integer layer in $N_\delta=n+\frac{\beta'+1}2$ is the $\beta$-free series
$F_p(N_\delta;\alpha)$, the same as for the lattice in its own natural index, and by Theorem~\ref{thm:shift}(ii) the same holds for the
layers $e_1,e_2,2e_1$. The estimates of steps (H-a)--(H-e) are written out in the notes of the repository named in the data availability
statement; here we only list them. Numerically, the model $M(\alpha+2\delta,\alpha)$ with the Hurwitz $\mu^\delta$ reproduces all $36$
stored fits of Hurwitz lattices of \S\ref{sec:hurwitz}, and excludes the alternatives (Bessel order moved with $\beta'$, no Szeg\H{o}
shift, lattice $\mu$) (script \texttt{m3a\_hd\_model}).

\subsection{Shifted Hurwitz lattices}\label{sec:kappa}
The first node sets with poles of $Z_\Lambda$ that we can treat are the Hurwitz lattices shifted by a constant,
$\lambda_k=(k+\delta)^p+\kappa$. For $p>1$, $Z_\Lambda(0)=-\frac12-\delta$ does not depend on $\kappa$, and $w\mapsto Z_\Lambda(-w)$ has the
moved Weyl pole $q=1-\frac1p$, with residue $R_q=-\kappa q/p$, and further poles at $q_m=m-\frac1p$, $m\ge2$, whose terms are
$O(n^{-2p})$. The pole law of Conjecture~\ref{conj:natural}(c) predicts the term $pR_qW(q)\pi\cot(\pi q)n^{-p}=-2(2-p)c\kappa\,n^{-p}$
(\S\ref{sec:num-pole}).

\begin{theorem}[shifted Hurwitz lattices]\label{thm:kappa}
Let $p>1$, $-1<\delta\le0$, $a=1+\delta$, $\kappa\ne0$ with $\kappa>-a^p$, $1<s<1+2p$, and
$\sigma^{\delta,\kappa}:=\sum_k\lambda_k^{-s/p}\delta_{1/\lambda_k}$ with $\lambda_k=(k+\delta)^p+\kappa$. Then, with $\beta'=\alpha+2\delta$,
\[
  \frac{n^pb_n}c=\sum_{0\le k<1+p}d_k(p,\beta',\alpha)\,n^{-k}-2(2-p)\,c\,\kappa\,n^{-p}+O(n^{-1-p}).
\]
Equivalently, the $n^{-k}$ coefficients with $k<p$ do not depend on $\kappa$ and are those of Theorem~\ref{thm:H}, and the
coefficient of $n^{-p}$ is that of the Hurwitz lattice plus the pole term $pR_qW(q)\pi\cot(\pi q)$ of
Conjecture~\ref{conj:natural}(c).\footnote{The last step of the proof, the first-order response of the model to the branched tilt
(the splitting of the gauge of a tilt $e^{-\mu z^\rho}$ into an analytic part with the Taylor data of the symbolic-$j$ layer at $j=\rho$
and a remainder $\pm i\tan(\pi\rho)z^\rho$), and the proposition that uses it, were re-read independently just before this
draft was assembled; no mathematical error was found, and the minor corrections from that reading were still being incorporated
into the underlying notes at the time of writing.}
\end{theorem}

By Theorem~\ref{thm:shift}, $\sum_kd_k(p,\beta',\alpha)n^{-k}=(n/N_\delta)^p(1+f_2N_\delta^{-2}+f_4N_\delta^{-4}+\dotsb)$, so the integer part is
the lattice's natural form in $N_\delta$, and the pole term sits at the same natural index. For $1<p<2$ the pole term lies between
$d_1/n$ and $d_2/n^2$; at $p=2$ its coefficient vanishes.

\begin{proof}[Outline]
The index map is a M\"obius change of variable: $\kappa$-nodes satisfy $t_\kappa(z)=k+\delta$ with $t_\kappa=t\circ\zeta_\kappa$,
$\zeta_\kappa(z)=z/(1-\kappa z/n^p)$. Composing the Hurwitz outer scalar with $\zeta_\kappa$ gives a scalar with the right jumps, which,
unlike the Hurwitz one, behaves like $C(-z)^{-1/2}$ at infinity: the exponent at infinity of an unnormalised outer scalar is not an
invariant. The invariant is the exponent of the normalising gauge at the hard edge. We take the gauge
$G_\kappa=(-z)^{-\delta}\exp\mathcal C[\rho_1]$, where $\rho_1=t-t_\kappa$ is minus the index offset and $\mathcal C$ a Cauchy transform on
$(-\infty,0)$; then the normalised scalar tends to a nonzero constant at infinity, and at the hard edge
$G_\kappa=(-z)^{-\delta}E_\kappa(0)(1+O(n^{1-p}|z|^{1-1/p}))$, so the Szeg\H{o} shift is $\beta'=\alpha+2\delta$, independent of $\kappa$.
The jumps are those of $M(\beta',\alpha)$ with a band factor $H_n^\kappa$ and an extra lip factor $e^{\pm2\pi i\rho_1}$. The integer powers of
$\log H_n^\kappa$ give the Toda parameters $2Z_\Lambda(-j)/j$ minus the density coefficients (the gauge supplies the term
$\kappa^j\zeta(0,a)$), and the branched part is the tilt $\sum_m2\pi\cot(\pi q_m)a_m(z/n^p)^{q_m}$, $a_m=-R_{q_m}/q_m$. The branched tilt
and the lip factor come from one function, analytic across $(-\delta_0,0)$, so the circle jumps stay meromorphic (Lemma~\ref{lem:N6b})
and no further hard-edge term arises. The a-priori bound, the $\alpha$-frame comparison and the small-norm expansion are transferred
from steps (H-b)--(H-e); for $1<p<2$ the H\"older exponent of the band factor at $0$ drops below $\frac12$, but every use of it
survives because the lip coefficient carries the gauge only through a factor with exponent $\frac12$. The first-order response of
the model to the tilt $z^{q}$ is the symbolic-$j$ layer of Lemma~\ref{lem:N6d}(b) continued to $j=q$, so that a tilt $e^{\tau x^q}$ of the
weight contributes $-\frac{pqW(q)}2\tau(n/N_\delta)^pN_\delta^{-1-pq}$ at leading order; for $\tau=2\pi\cot(\pi q)a_1$ this is exactly $pR_qW(q)\pi\cot(\pi q)(n/N_\delta)^pN_\delta^{-1-pq}$.
\end{proof}

The proof gives more: up to $O(n^{-2p})$, $n^pb_n/c$ is the full double series of $M(\beta',\alpha)$ with the Toda parameters of
$\Lambda$, plus for each pole the first-order layer of Lemma~\ref{lem:N6d}(b) continued to $j=q_m$, evaluated at the tilt
$\tau_m=2\pi\cot(\pi q_m)a_m$. It says nothing about products of pole terms. Fits at $(p,\delta,\kappa)=(2.5,-0.3,\pm0.3)$ and $(1.5,-0.4,0.4)$ confirm $d_1$, $d_2$ of $M(\alpha+2\delta,\alpha)$ to
$10^{-10}$ and $5\cdot10^{-8}$, and the $n^{-p}$ term to $4\cdot10^{-7}$; the Szeg\H{o} exponent read from the behaviour of the
unnormalised scalar at infinity is excluded (script \texttt{w\_outer\_e2e}). For $p<1$ the route
fails: the gauge density $\rho_1$ is of size $|\kappa|n^{(1-p)(2+p)/p}$ on the Bessel scale, so the lattice $g$-function cannot carry
it; the rows $p=0.6$ and $p=0.3$ of \S\ref{sec:num-pole} are therefore evidence, not instances of Theorem~\ref{thm:kappa}.

% Section 6: Examples and numerical evidence.
\section{Examples and numerical evidence}\label{sec:examples}

This section tests Theorem~\ref{thm:N6} on the lattice and collects the evidence for the conjectures of
Section~\ref{sec:conj} on other node sets. No proof uses it. The computations are floating-point; the scripts and their
stored outputs are listed in Appendix~\ref{app:repro}.

\emph{Method.} The coefficients $b_n$, $n\le1600$, are computed in quadruple precision by adding the atoms
$1/\lambda_m$, $m\le K$, one at a time and restoring tridiagonal form by plane rotations, with a Gauss--Jacobi rule for the
continuous tail beyond $K$; $K$ is $100$ to $200$ times the largest $n$. Expansion coefficients are obtained by
least-squares fits of $n$-windows ($n\ge200$, $400$, $800$ or mid-range windows) after subtracting the terms that are known
exactly, with a pool of nuisance exponents; the number in parentheses is the spread of the fitted value over windows and
pools. When a prediction was formed after looking at some of the data, we say which rows were used to form it; those
rows are not evidence for it.

\subsection{The lattice in natural form}\label{sec:num-lattice}
Table~\ref{tab:natural} maps the lattice coefficients to $F(N):=(N/n)^pn^pb_n/c$, subtracts $1+f_2N^{-2}+\sum_j\varepsilon_jN^{-1-jp}$
and fits the remainder. The $N^{-1}$ and $N^{-2}$ residuals re-confirm $d_1$ and $d_2$. The $N^{-3}$ residual is zero
to the precision of the tail windows, against lattice $n^{-3}$ coefficients of size $0.1$ to $1$ (for example $-0.267$ at
$p=1.5$, $s=2$); its sign follows $\alpha$ at $p=1.5$, but the large-$n$ end of the data is noise-dominated at this scale,
and mid-range windows $n\in[10,400]$ give $N^{-3}$ coefficients below $4\cdot10^{-8}$ in all cases of
Table~\ref{tab:f4}. Theorem~\ref{thm:N6} proves the absence of the $N^{-3}$ term for $p>1$.

\begin{table}[ht]
\centering
\small
\begin{tabular}{rrrrrr}
\toprule
$p$ & $s$ & $\alpha$ & $N^{-1}$ residual & $N^{-2}$ residual & $N^{-3}$ residual\\
\midrule
1.5 & 2 & $-\frac13$ & $-4\cdot10^{-11}$ & $+1.5\cdot10^{-7}$ & $-3.0\cdot10^{-4}$ $(2\cdot10^{-4})$\\
1.5 & 3 & $+\frac13$ & $+5\cdot10^{-11}$ & $-1.6\cdot10^{-7}$ & $+3.2\cdot10^{-4}$ $(2\cdot10^{-4})$\\
1.25 & 2 & $-0.2$ & $+4\cdot10^{-11}$ & $-1.5\cdot10^{-7}$ & $+5.2\cdot10^{-4}$ $(4\cdot10^{-4})$\\
2.5 & 2 & $-0.6$ & $-2\cdot10^{-11}$ & $+5\cdot10^{-8}$ & $-6\cdot10^{-5}$ $(6\cdot10^{-5})$\\
3 & 2 & $-\frac23$ & $+2\cdot10^{-12}$ & $-4\cdot10^{-9}$ & $+2.3\cdot10^{-6}$ $(5\cdot10^{-7})$\\
1.5 & 2.5 & $0$ & $+2\cdot10^{-12}$ & $-7\cdot10^{-9}$ & $+1.3\cdot10^{-5}$ $(2\cdot10^{-5})$\\
2 & 3 & $0$ & $-2\cdot10^{-13}$ & $+4\cdot10^{-10}$ & $-3\cdot10^{-7}$ $(3\cdot10^{-7})$\\
\bottomrule
\end{tabular}
\caption{Residual coefficients of $F(N)-1-f_2N^{-2}-\sum_j\varepsilon_jN^{-1-jp}$ for the lattice, tail windows
$n\ge200,400$ (script \texttt{m4\_natural\_index}).}
\label{tab:natural}
\end{table}

\begin{table}[ht]
\centering
\small
\begin{tabular}{rrrrr}
\toprule
$p$ & $\alpha$ & $f_4$ fitted & $f_4$, \eqref{eq:f4} & difference\\
\midrule
2 & 0 & $0.044921952$ $(1\cdot10^{-6})$ & $0.044921875$ & $7.7\cdot10^{-8}$\\
2 & 1 & $-0.064452710$ $(2\cdot10^{-6})$ & $-0.064453125$ & $4.2\cdot10^{-7}$\\
4 & 0 & $0.083271330$ $(3\cdot10^{-7})$ & $0.083271329$ & $2.9\cdot10^{-10}$\\
4 & 1 & $-0.108692879$ $(9\cdot10^{-7})$ & $-0.108692956$ & $7.7\cdot10^{-8}$\\
2.5 & 0.5 & $0.029254295$ $(1\cdot10^{-4})$ & $0.029206272$ & $4.8\cdot10^{-5}$\\
\bottomrule
\end{tabular}
\caption{The $N^{-4}$ coefficient of the lattice, mid-range windows (script \texttt{m4\_f4\_lattice}). The polynomial
$f_4(2,\alpha)=(23-40\alpha^2-16\alpha^4)/512$ was first read off fits of this kind; the exact computation of
Lemma~\ref{lem:N6d} was built independently.}
\label{tab:f4}
\end{table}

For the zeta layer, after subtracting the integer layer and the bare terms $\varepsilon_jN^{-1-jp}$, mid-range fits at
$p\in\{1.5,2.5\}$ give an $N^{-2-p}$ coefficient below $4\cdot10^{-7}$ and an $N^{-3-p}$ coefficient equal to $\varepsilon_1g_1$
within $4\cdot10^{-6}$ to $6\cdot10^{-5}$, with spreads $8\cdot10^{-4}$ to $10^{-2}$ (script \texttt{m4\_zeta\_lattice}); at
$p=1.5$ the values $g_1=95/288$ ($\alpha=0$) and $22/9$ ($\alpha=1$) were first seen in these fits.

\subsection{Hurwitz lattices}\label{sec:hurwitz}
For $-1<\delta\le0$ (and $\delta<0$ if $p<1$) and $1<s<1+2p$ the statements tested here are Theorem~\ref{thm:H}; the rows with $a\ge\frac32$
($\delta>0$) or $s\ge1+2p$ are outside it, and all rows predate the theorem.
For $\lambda_m=(m+\delta)^p$, $\delta>-1$, one has $Z_\Lambda(w)=\zeta(pw,a)$ with $a=1+\delta$, so $Z_\Lambda(-j)=\zeta(-jp,a)$ and
$Z_\Lambda(0)=\frac12-a$. The natural-index law of Section~\ref{sec:conj} predicts the $1/n$ coefficient
$d_1=-\frac{s-1}2-p\delta$ and the $n^{-1-p}$ coefficient $2pc\,\zeta(-p,a)$ in place of $\kp$. Table~\ref{tab:hurwitz}
compares both with fits; the lattice value $\kp$ is excluded in every row.

\begin{table}[ht]
\centering
\small
\begin{tabular}{rrrrrrr}
\toprule
$p$ & $s$ & $a$ & $d_1$ fitted & $d_1$ predicted & $n^{-1-p}$ fitted & $2pc\,\zeta(-p,a)$\\
\midrule
0.6 & 2 & $\frac12$ & $-0.19999996$ & $-0.2$ & $+0.0306156$ $(8\cdot10^{-5})$ & $+0.0306142$\\
0.6 & 2 & $\frac32$ & $-0.80000024$ & $-0.8$ & $-0.309374$ $(1.4\cdot10^{-4})$ & $-0.309386$\\
0.6 & 2 & $2$ & $-1.10000058$ & $-1.1$ & $-0.605309$ $(3\cdot10^{-4})$ & $-0.605320$\\
0.75 & 3 & $2$ & $-1.75000033$ & $-1.75$ & $-0.779282$ $(1.3\cdot10^{-3})$ & $-0.779293$\\
1.5 & 2 & $\frac32$ & $-1.25000000$ & $-1.25$ & $-0.57244$ $(5.5\cdot10^{-3})$ & $-0.57263$\\
1.5 & 2 & $2$ & $-2.00000000$ & $-2$ & $-1.74157$ $(5.6\cdot10^{-3})$ & $-1.74211$\\
\bottomrule
\end{tabular}
\caption{Hurwitz lattices (script \texttt{spec\_hur}). The lattice values of the $n^{-1-p}$ coefficient are $-0.0899767$
($p=0.6$), $-0.0918692$ ($p=0.75$) and $-0.0432947$ ($p=1.5$).}
\label{tab:hurwitz}
\end{table}

The $n^{-2}$ coefficients of Hurwitz lattices at $p=0.6$ ($a=\frac12,\frac32,2$), $0.75$ and $1.5$ agree with the
index-shifted lattice value $d_2+\frac{p(p+1)}2\delta^2-(p+1)\delta d_1$ to between $2\cdot10^{-6}$ and $3\cdot10^{-5}$
(spreads $3\cdot10^{-5}$ to $9\cdot10^{-4}$), for shifts of size $0.3$ to $3$ (script \texttt{spec\_hur\_d2}). At $p=1.5$ the
rule ``the lattice expansion as a function of $n+\delta$, multiplied by $(n/(n+\delta))^p$, with every $\zeta(-jp)$ replaced by
$\zeta(-jp,a)$'' also predicts the $n^{-3}$ coefficients ($\delta=-\frac12,\frac12,1$; agreement within the fit ranges,
$1\cdot10^{-4}$ to $1.3\cdot10^{-2}$, relative $3\cdot10^{-4}$ to $2\cdot10^{-3}$ of the shifts; this uses the fitted lattice
$d_3$) and the $n^{-3.5}$ coefficients, where the alternative ``zeta terms unshifted'' is excluded at $\delta=\frac12,1$
(script \texttt{spec\_shift\_hi}). At $p=\frac13$ and $\frac14$ ($\delta=-\frac12$) the first two, resp.\ three, coefficients
$2pc\,\zeta(-p,a)$, $6pc^2\zeta(-2p,a)$, $20pc^3\zeta(-3p,a)$ agree to between $4\cdot10^{-7}$ and $3\cdot10^{-4}$, and at
$p=\frac13$ the merged $n^{-2}$ coefficient requires the term $\frac13\zeta(-1,a)$: without it the fit is off by seven
spreads (script \texttt{spec\_hur\_small}).

\subsection{The pole law on synthetic lattices}\label{sec:num-pole}
For $\lambda_k=(k+\gamma k^{-m})^p$ the counting function is unchanged at order $k^0$, and to first order in $\gamma$,
$Z_\Lambda(-w)=\zeta(-pw)+pw\gamma\zeta(m+1-pw)+O(\gamma^2)$ has a new pole at $q=m/p$ with residue $R=-m\gamma/p$. The pole
law (Conjecture~\ref{conj:natural}(c)) predicts the term $pRW(q)\pi\cot(\pi q)n^{-1-pq}=-m\gamma W(m/p)\pi\cot(\pi m/p)n^{-1-m}$ in
$n^pb_n/c$. Table~\ref{tab:pole} compares it with fits of $n(y_n-y_n^{\mathrm{lat}})$ after subtraction of the first-order
zeta terms $pW(j)\delta Z(-j)n^{-jp}$. For $m=1$, $W(1/p)=1$, so the rows test the cotangent; the rows $m=\frac12$ test
$W$ as well.

\begin{table}[ht]
\centering
\small
\begin{tabular}{rrrrrl}
\toprule
$p$ & $m$ & $\cot(\pi m/p)$ & predicted & fitted, $n\ge200$ / $400$ / $800$ & role\\
\midrule
0.6, 1.5 & 1 & $-0.5774$ & $+0.906900$ & $+0.906900$ & identification\\
1.25 & 1 & $-1.3764$ & $+2.162016$ & $+2.162016$ / $+2.162017$ / $+2.162022$ & prediction\\
1.8 & 1 & $-0.1763$ & $+0.276974$ & $+0.276974$ / $+0.276974$ / $+0.276973$ & prediction\\
0.8 & 1 & $+1$ & $-1.570796$ & $-1.570783$ / $-1.570792$ / $-1.570813$ & prediction\\
0.75 & 1 & $+0.5774$ & $-0.906900$ & $-0.906866$ / $-0.906885$ / $-0.906954$ & prediction\\
1.5 & $\frac12$ & $+0.5774$ & $-0.424704$ & $-0.424703$ / $-0.424697$ / $-0.424630$ & prediction\\
1.25 & $\frac12$ & $+0.3249$ & $-0.235183$ & $-0.235177$ / $-0.235112$ / $-0.232695$ & prediction\\
\bottomrule
\end{tabular}
\caption{The pole law, $s=2$, $\gamma=\frac12$ (scripts \texttt{spec\_pole}, \texttt{spec\_gamma}). The kernel $\pi\cot(\pi w)$ was identified from
the first row; the predictions of the other rows were fixed before the corresponding measures were computed. For $m=\frac12$,
$W(m/p)=0.9366$ ($p=1.5$) and $0.9216$ ($p=1.25$). The $n\ge800$ windows drift where few points remain.}
\label{tab:pole}
\end{table}

\emph{A constant shift.} For $\lambda_k=k^p+\kappa$ the index offset grows like $k^{1-p}$, and the expansion changes at
relative order $n^{-p}$, not $n^{-1-p}$. To first order, $Z_\Lambda(-w)=\zeta(-pw)+\kappa w\zeta(p(1-w))$ has its Weyl pole moved
to $q=1-1/p$, with residue $-\kappa q/p$, and the pole law gives $-2(2-p)c\kappa\,n^{-p}$ in $n^pb_n/c$, using
$W(-1/p)=p\tan(\pi/p)/\pi$. Fits of the $n^{-p}$ coefficient give $-1.20342$ $(0.13)$, $-0.60154$ $(0.04)$, $-1.22311$ $(0.74)$
at $(p,\kappa)=(0.6,1),(0.6,\frac12),(0.3,1)$, against $-1.20247$, $-0.60123$, $-1.22262$ (script \texttt{spec\_kappa\_fit}); the
first row was the identification point, and at $p=0.3$ the spread is $60\%$ of the value, so that row is not evidence.
For $p>1$ the law is now Theorem~\ref{thm:kappa} (with $\delta=0$); the rows here have $p<1$, where it is not proved.
The Hankel determinants of the shifted measure are those of the old nodes with an $m$-dependent weight: since
$\varphi(y)=y/(1+\kappa y)$ satisfies $\varphi(a)-\varphi(b)=(a-b)/((1+\kappa a)(1+\kappa b))$, Heine's formula gives
$D_m(\sigma_\kappa)=D_m((1+\kappa y)^{-2(m-1)}\sigma')$ with $\sigma'=(1+\kappa y)^{-s/p}\sigma$, and a first-order expansion of
$\log D_m$ in $\kappa$ reproduces the law, with a relative gap that decays like $\kappa n^{-p}$ (script \texttt{spec\_mobius\_pred}).

\subsection{The subleading pole term}\label{sec:num-subpole}
The pole law gives the leading power $n^{-1-pq}$ of a pole. The same pole contributes at $n^{-2-pq}$, and this term is
not removed when $\cot(\pi q)=0$ (Remark~\ref{rem:sublead}). For the family of \S\ref{sec:num-pole}, the part linear in
$\gamma$ of the $n^{-2-m}$ coefficient of $n^pb_n/c$, isolated by antisymmetrising in $\gamma\in\{\pm\frac14,\pm\frac12\}$ and
Richardson extrapolation, is compared in Table~\ref{tab:subpole} with
\begin{equation}\label{eq:S4}
  S=-\frac{mpq_1}2W(q_1)+\mathcal L\,\frac{(1+m+p)\,d_1}p,\qquad q_1:=\frac{m+1}p,\qquad
  \mathcal L:=-mW(q)\pi\cot(\pi q),
\end{equation}
per unit $\gamma$, with $d_1=-(s-1)/2$. The first term is the response to the change of node density $\log|t'|$, a tilt
$\propto z^{q_1}$ branched at the hard edge; the second is the leading pole term evaluated at the shifted index
$n+(s-1)/(2p)=n+\delta_N$.

\begin{table}[ht]
\centering
\small
\begin{tabular}{rrrrrrl}
\toprule
$p$ & $s$ & $m$ & $q_1$ & $S$ fitted & \eqref{eq:S4} & role\\
\midrule
2 & 2 & 1 & 1 & $-1.233702$ $(1\cdot10^{-6})$ & $-1.233701$ & identification\\
2.5 & 2 & 1 & 0.8 & $-0.258693$ $(2\cdot10^{-5})$ & $-0.258696$ & identification\\
3 & 2 & 1 & 0.667 & $+0.371552$ $(1\cdot10^{-6})$ & $+0.371551$ & identification\\
2.2 & 2 & 1 & 0.909 & $-0.777016$ $(7\cdot10^{-5})$ & $-0.776976$ & prediction\\
4 & 2 & 1 & 0.5 & $+1.261975$ $(4\cdot10^{-7})$ & $+1.261975$ & prediction\\
3 & 2 & 2 & 1 & $-8.270517$ $(3\cdot10^{-5})$ & $-8.270552$ & prediction\\
3 & 3 & 1 & 0.667 & $+1.883052$ $(1\cdot10^{-6})$ & $+1.883050$ & prediction\\
2.5 & 3 & 1 & 0.8 & $+0.659993$ $(2\cdot10^{-5})$ & $+0.659993$ & prediction\\
1.8 & 2 & 1 & 1.111 & $-1.847995$ $(3\cdot10^{-4})$ & $-1.849477$ & see text\\
1.5 & 2 & 1 & 1.333 & $-3.39874$ $(1.4\cdot10^{-2})$ & $-3.44204$ & see text\\
1.25 & 2 & 1 & 1.6 & $-6.54$ $(0.26)$ & $-7.02$ & see text\\
\bottomrule
\end{tabular}
\caption{The subleading pole coefficient (script \texttt{spec\_subpole}). At $p=2$, $\mathcal L=0$ and $S=-W(1)=-\pi^2/8$.}
\label{tab:subpole}
\end{table}

For $q_1\le1$ the formula holds to the fit precision in the eight rows; three of them were used to form it. For $q_1>1$
the exponents crowd near $p=1$ and the fits depend on the nuisance pool: refits of the stored data over twenty pools
(windows $n\ge150$ to $600$, three to seven nuisance exponents) give medians minus prediction of $-2\cdot10^{-3}$ ($p=1.5$),
$+3\cdot10^{-3}$ ($p=1.8$) and $-9\cdot10^{-2}$ ($p=2.5$, $m=2$), with signs depending on the pool. The same decomposition, at
$p=2$, $m=1$, gives a nonzero $n^{-3}$ coefficient $-3.70\gamma$ for this family; for Bessel zeros (\S\ref{sec:bessel}),
whose index correction at $p=2$ is tied to the offset, the corresponding terms cancel, and the expansion is the pure
index shift of Proposition~\ref{prop:bessel}.

\subsection{Conjecture~\ref{conj:W}: predictions before fits}\label{sec:num-W}
To test Conjecture~\ref{conj:W} on node sets with poles, with $Z(0)\ne-\frac12$, and with masses that are not pure powers, every
coefficient was predicted from the spectral data alone (regularised $Z(0)$ and $Z(-j)$, residues, density and mass factors) and the
recursion of \S\ref{sec:recursion} at $40$ digits (Toda responses by central differences), and the predictions were stored before
any of the nine measures was computed; the test script asserts the time stamps. The conjecture has no fitted constant. The
alternatives are the conjecture with one piece switched off: no Szeg\H{o} shift ($\beta=\alpha$), no pole term, no $Z(-j)$, no density or
mass tilt, no natural-index shift of the pole terms. A coefficient passes if $|\text{fit}-\text{prediction}|\le\max(3\cdot\text{spread},
2\cdot10^{-6}(1+|\text{prediction}|))$. Of the $41$ coefficients tested, $40$ pass; the exception is the sphere at $n^{-4}$, outside the
claimed precision (Remark~\ref{rem:Wprime}). Table~\ref{tab:W} shows a selection.

The families are: (A) $\lambda_k=k^p+ak^{p-1}$, $s=2$, which is a Hurwitz offset $\delta=a/p$ plus decaying corrections, so
$Z(0)=-\frac12-\frac ap$, with poles at $q_i=(i-1)/p$; (B) Bessel zeros $\lambda_k=(j_{\nu',k}/\pi)^P$, $s=2$, with $Z(0)=-\frac{\nu'}2-\frac14$
and poles at $q=(2i-1)/P$ (the residues use McMahon's expansion, checked against computed zeros to $5\cdot10^{-26}$); (M2) the nodes
$l^{3/2}$ with masses $(2l+1)l^{-s}$, $s=3$, a mass factor with the fractional time $\rho=\frac23$; (M1) the sphere, nodes $l(l+1)$ with
multiplicity $2l+1$, $s=4$.

\begin{table}[ht]
\centering
\footnotesize
\setlength{\tabcolsep}{4pt}
\begin{tabular}{llrrrr}
\toprule
measure & coeff. & prediction & fit (spread) & difference & next alternative\\
\midrule
A: $p=1.5$, $a=-\frac34$ & $n^{-2}$ & $+0.0984431$ & $+0.0984431$ $(3\cdot10^{-5})$ & $-3\cdot10^{-9}$ & no pole: $1392$\\
 & $n^{-2.5}$ & $+0.2215758$ & $+0.2215749$ $(4\cdot10^{-5})$ & $-9\cdot10^{-7}$ & no $Z(-j)$: $1912$\\
 & $n^{-3}$ & $+0.0699377$ & $+0.0699393$ $(1.1\cdot10^{-3})$ & $+1.6\cdot10^{-6}$ & no tilt: $25$\\
A: $p=1.5$, $a=0.6$ & $n^{-2}$ & $+1.0955036$ & $+1.0955031$ $(1.6\cdot10^{-4})$ & $-5\cdot10^{-7}$ & no pole: $150$\\
A: $p=3$, $a=-0.6$ & $n^{-2}$ & $+0.1347742$ & $+0.1347742$ $(5\cdot10^{-9})$ & $+5\cdot10^{-10}$ & no pole: $31968$\\
 & $n^{-3}$ & $+0.0452648$ & $+0.0452649$ $(2.7\cdot10^{-6})$ & $+1.0\cdot10^{-7}$ & no shift: $1466$\\
A: $p=0.6$, $a=-0.3$ & $n^{-1.6}$ & $+0.0854664$ & $+0.0854637$ $(3.7\cdot10^{-4})$ & $-3\cdot10^{-6}$ & no $Z(-j)$: $77$\\
 & $n^{-2}$ & $+0.1329122$ & $+0.1327116$ $(1.1\cdot10^{-3})$ & $-2.0\cdot10^{-4}$ & no pole: $27$\\
B: $\nu'=0$, $P=1.5$ & $n^{-2}$ & $+0.1957151$ & $+0.1957150$ $(1.4\cdot10^{-5})$ & $-1.3\cdot10^{-7}$ & no pole: $562$\\
B: $\nu'=1$, $P=1.5$ & $n^{-2.5}$ & $-0.0664936$ & $-0.0664957$ $(5\cdot10^{-4})$ & $-2\cdot10^{-6}$ & no $Z(-j)$: $42$\\
B: $\nu'=0$, $P=3$ & $n^{-2}$ & $+0.0742502$ & $+0.0742502$ $(3\cdot10^{-9})$ & $+5\cdot10^{-10}$ & no pole: $10692$\\
 & $n^{-3}$ & $+0.0049259$ & $+0.0049264$ $(1.2\cdot10^{-5})$ & $+4.8\cdot10^{-7}$ & no shift: $276$\\
M2 & $n^{-2}$ & $+0.1180556$ & $+0.1180558$ $(3.8\cdot10^{-6})$ & $+2.7\cdot10^{-7}$ & no mass tilt: $21830$\\
 & $n^{-3}$ & $+0.1900476$ & $+0.1900481$ $(1.1\cdot10^{-3})$ & $+4.8\cdot10^{-7}$ & no shift: $90$\\
M1 (sphere) & $n^{-3}$ & $-5$ & $-4.99999992$ $(3\cdot10^{-6})$ & $8\cdot10^{-8}$ & one tilt only: $12851$\\
 & $n^{-4}$ & $+7.5449219$ & $+7.776241$ $(3.3\cdot10^{-4})$ & $+0.23132$ & not claimed\\
\bottomrule
\end{tabular}
\caption{Conjecture~\ref{conj:W} against fits (scripts \texttt{w\_predict}, \texttt{w\_test}). The last column gives the miss of the
closest alternative in units of the tolerance. The row $a=0.6$ has $\delta=0.4>0$, outside Theorem~\ref{thm:H}; the rows $p=0.6$ are the
weakest tests. Fit-free residuals $n^e(y_n-\sum_{e'<e}W_{e'}n^{-e'})-W_e$ at $n=200,400,800,1600$ decay in the rows shown except the sphere's $n^{-4}$ (script
\texttt{w\_resid}).}
\label{tab:W}
\end{table}

The $n^{-1}$ coefficients test the Szeg\H{o} shift $\beta=\alpha-2Z(0)-1$ at $Z(0)\ne-\frac12$: they agree to $10^{-8}$ or better, and $\beta=\alpha$
is excluded by $10^3$--$10^5$ tolerances. The $n^{-3}$ rows of the families A and B test, beyond Conjecture~\ref{conj:natural}, the natural-index
shift of the pole terms with $Z(0)\ne-\frac12$ and the density tilt as a fractional Toda time. At $q_1=\frac43$ (family A, $p=1.5$) the agreement
is $2\cdot10^{-6}$--$3\cdot10^{-5}$ against spreads of $10^{-3}$: consistent, not precise. For the Bessel anchor at $P=2$
(Proposition~\ref{prop:bessel}) the prediction of Conjecture~\ref{conj:W}, the bare model $M(\nu'-1,-\frac12)$, reproduces the coefficients
$n^{-1},\dots,n^{-4}$ of the exact closed form to $2\cdot10^{-21}$. The probes of Remark~\ref{rem:Wprime} were designed after the sphere's
$n^{-4}$ failure (script \texttt{w\_probe}); the two measures $k^2+\frac14$ at $s=2$ and $s=3.5$ were then computed after the predictions of the
refinement had been stored (script \texttt{w\_probe2}).

\subsection{Drums}\label{sec:drums}
Let $\lambda$ run over the Dirichlet eigenvalues of a planar domain of area $\mathcal A$ and perimeter $\mathcal P$,
rescaled to $\tilde\lambda=\mathcal A\lambda/4\pi$, so that the counting function is
$\tilde\lambda-\beta\tilde\lambda^{1/2}+c_0+(\text{fluctuations})$ with $\beta=\mathcal P/\sqrt{4\pi\mathcal A}$; the
constant $c_0=Z(0)$ carries corners and curvature ($\frac14$ for rectangles, $\frac13$ for the equilateral triangle,
$\frac16$ for smooth simply connected domains). We take $p=1$, $s=2$, which is outside the range of the theorems. A
\emph{smoothed drum} is the node set defined by the counting function without fluctuations, $\mathcal N(\lambda_k)=k-\frac12-c_0$.
The perimeter is a pole of $Z(-w)$ at $q=-\frac12$ with residue $\beta/2$, where $\cot(\pi q)=0$ and $W$ has a pole; the
limit $W(q)\pi\cot(\pi q)\to-\pi(4c)^{-1/2}$ of \eqref{eq:Lt-W} gives the prediction
$n(y_n-y_n^{\mathrm{lat}})\supset-\frac{\pi\beta}{2\sqrt2}n^{1/2}$, which is the first term of Remark~\ref{rem:drum-root}.

\begin{table}[ht]
\centering
\small
\begin{tabular}{lrrr}
\toprule
smoothed drum & $\beta$ & fitted $n^{1/2}$ coefficient & $-\pi\beta/(2\sqrt2)$\\
\midrule
square & $1.128379$ & $-1.253317$ $(1\cdot10^{-5})$ & $-1.253314$\\
equilateral triangle & $1.286074$ & $-1.428473$ $(1\cdot10^{-5})$ & $-1.428469$\\
disk & $1$ & $-1.110722$ $(7\cdot10^{-6})$ & $-1.110721$\\
golden rectangle & $1.161199$ & $-1.289770$ $(1\cdot10^{-5})$ & $-1.289767$\\
\bottomrule
\end{tabular}
\caption{The perimeter term at $p=1$ (script \texttt{spec\_corners}). The prediction, made before the run, was that no
$n^{1/2}$ term occurs, which forgot the pole of $W$; the value in the last column is the $0\cdot\infty$ limit of the same
formula, with no free parameter.}
\label{tab:drums}
\end{table}

The corner term is isolated by comparing two smoothed drums with the same $\beta$ (the square's) and $c_0=\frac14$,
$\frac13$: the $n^0$ coefficient of the difference is $0.0833487$ $(2\cdot10^{-5})$, against $\Delta c_0=\frac1{12}$, as the term
$pZ(0)$ of the dictionary requires. The absolute $n^0$ coefficient has median $c_0+\frac12+\frac{\pi^2}{16}\beta^2$ over fits
in all smoothed drums, to $2\cdot10^{-4}$, but single fits of it are unstable; the $\beta^2$ part is the second term of
Remark~\ref{rem:drum-root}, i.e.\ it is produced by solving the continuum soft-edge condition exactly rather than to first
order. For the actual eigenvalues of the golden rectangle, the deviation from the smoothed drum has mean near $0$ and a
standard deviation growing like $n^{0.55}$: the counting fluctuations of the spectrum dominate everything below the
perimeter term (\S\ref{sec:noise}).

\subsection{Anharmonic oscillators}\label{sec:quartic}
For the eigenvalues $E_K$ ($K\ge1$) of $-\frac{\diff^2}{\diff q^2}+q^N$ on $\R$, $N\ge4$ even, the semiclassical
quantisation reads $K-\frac12=AE^{a}+BE^{-a}+O(E^{-3a})$, $a=\frac12+\frac1N$, with
$A=\frac1\pi\int_{-1}^1\sqrt{1-t^N}\diff t=\frac2{\pi N}B(\frac1N,\frac32)$. With $p:=1/a=2N/(N+2)$ and $\lambda_K:=A^pE_K$,
\[
  \lambda_K^{1/p}=k'+\gamma/k'+O(k'^{-3}),\qquad k'=K-\tfrac12,\qquad\gamma=-AB,
\]
a Hurwitz lattice with $a=\frac12$ plus an index correction with $m=1$, whose pole $q=1/p$ has $W(q)=1$ and
$\cot(\pi/p)=-\tan(\pi/N)\ne0$. (For Bessel zeros, and for linear potentials, the poles of $Z(-w)$ sit at half-integers,
where the cotangent vanishes; here it does not.) The predictions are $d_1=-\frac{s-1}2+\frac p2$ ($Z(0)=\zeta(0,\frac12)=0$), no $n^{-1-p}$ term (the
heat-trace exponents of $q^N$ contain no integer, so $Z(-1)=0$), and an $n^{-2}$ coefficient equal to the Hurwitz-shifted
lattice value plus the pole term $-\gamma\pi\cot(\pi/p)=AB\pi\cot(\pi/p)$.

\emph{The pole term is rational.} If the second-order quantisation condition
\[
  \frac1\pi\int_{-q_+}^{q_+}\sqrt{E-V}\diff q-\frac1{24\pi}\,\partial_E^2\int_{-q_+}^{q_+}V'(q)^2(E-V)^{-1/2}\diff q=K-\tfrac12+O(E^{-3a})
\]
holds, then for $V=q^N$ the substitutions $q=E^{1/N}t$, $u=t^N$ give $B=-\frac N{12\pi}e(e-1)B(2-\frac1N,\frac12)$ with
$e=\frac32-\frac1N$, and the reflection formulas collapse the product to
\[
  AB=-\frac{(N-1)\cot(\pi/N)}{6\pi(N+2)},\qquad\text{hence}\qquad AB\,\pi\cot\frac\pi p=\frac{N-1}{6(N+2)},
\]
since $\pi/p=\pi/2+\pi/N$. The cotangent of the kernel cancels the $\cot(\pi/N)$ of the semiclassical coefficient. The
values are $\frac1{12},\frac5{48},\frac7{60},\frac18,\frac{11}{84}$ for $N=4,\dots,12$; with $B$ fitted from exact eigenvalues they
are reproduced to between $10^{-10}$ and $7\cdot10^{-8}$ (script \texttt{spec\_homog}); at $N=4$ the fitted
$B=-0.0476724854$ equals $-1/(12\pi A)$ to $4\cdot10^{-10}$.
% TODO (citation): the second-order Bohr--Sommerfeld (Dunham) term is not quoted in the project's citation record;
% its coefficient -1/(24 pi) is checked numerically at N = 4 (B/(5 I_6) = -1/(24 pi)).

\begin{table}[ht]
\centering
\small
\begin{tabular}{rrrrrr}
\toprule
$N$ & $d_1$ fitted (predicted) & $n^{-1-p}$ fitted & $n^{-2}$ fitted & predicted & without pole term\\
\midrule
4 & $+0.16666667$ $(\frac16)$ & $1\cdot10^{-5}$ $(3\cdot10^{-3})$ & $+0.276010$ $(1\cdot10^{-5})$ & $+0.276042$ & $+0.192708$\\
6 & $+0.25000000$ $(\frac14)$ & $5\cdot10^{-5}$ $(2\cdot10^{-3})$ & $+0.315918$ $(3\cdot10^{-5})$ & $+0.315972$ & $+0.211806$\\
\bottomrule
\end{tabular}
\caption{Oscillators $-\frac{\diff^2}{\diff q^2}+q^N$, $s=2$, eigenvalues from a banded harmonic-oscillator basis converged to
$3\cdot10^{-16}$ for $K\le1200$, with the semiclassical series beyond (scripts \texttt{spec\_quartic}, \texttt{spec\_qN}). With the
sign of the pole term reversed the predictions would be $0.109375$ and $0.107639$.}
\label{tab:quartic}
\end{table}

\subsection{Irregular spectra}\label{sec:noise}
If the nodes deviate from a smooth node function by a fixed fraction of the spacing, the expansion is replaced by an
irregular component. For iid jitter $\lambda_m=(m+\xi_m)^p$, $\xi_m$ uniform on $(-0.3,0.3)$, the residual of $y_n$ after a smooth
fit has root mean square of order $n^{-1/2}$ (for example $1.3\cdot10^{-2}$ on $[100,200)$ and $4.2\cdot10^{-3}$ on $[800,1600)$ at
$p=0.6$), and the fitted $d_1$ is meaningless. Jitter confined to nodes deep inside the band for all $n$ considered is
invisible (root mean square $10^{-9}$ to $3\cdot10^{-7}$), and jitter confined to the first $300$ nodes stops producing noise once
the soft-edge threshold $\lambda\approx n^p/(4c)$ has passed them (scripts \texttt{spec\_jitter}, \texttt{spec\_window}).

\emph{Single-node kernel.} Moving node $k_0\in\{100,200,400,600\}$ by $\pm0.05$ and following
$F(n;k_0):=[D(+h)-D(-h)]/(2h)$, $D=n(y_n-y_n^{\mathrm{lat}})$, in $\tau=(n-n_*)/n_*^{1/3}$, where $n_*$ is the index at which
the node crosses the edge, gives one curve for all $k_0$: outside the band ($\tau\ge2$) it is the derivative of the per-node
law $-p(1-u)^{-1/2}$ to $1$--$3\%$; near $\tau=0$ it stays finite; inside the band it oscillates with a period of about ten
nodes, and its sums over $n$ nearly cancel (script \texttt{spec\_edgekernel}).

\emph{Two ensembles.} Replacing the first $3000$ nodes of the lattice ($p=0.6$, $s=2$) by unfolded eigenphases of random
unitary matrices or by a Poisson bridge gives the root mean square of $D$ in Table~\ref{tab:noise}. Both grow like $n^{1/2}$
(exponents $0.50$ and $0.52$), and the Poisson amplitude is larger by the stable factor $3.3$--$4.0$. The noise is not a
linear statistic realisation by realisation (correlations between $-0.04$ and $0.47$ between $D$ and its linearisation), but the
linearisation has the same root mean square in every window (script \texttt{spec\_noiselin}). The power spectrum of the
linear kernel, measured with sinusoidal displacements $\epsilon\sin(2\pi k/\Lambda_0)$, scales as $n\kappa(\omega)$ with
$\kappa(\omega)\propto\omega$ for small $\omega$ (script \texttt{spec\_kernelfft}). The formula
$\Var D\approx n\int_0^\pi\kappa(\omega)S_h(\omega)\diff\omega/\pi$, with displacement spectrum
$S_h=K(\omega)/(4\sin^2(\omega/2))$, $K=1$ (Poisson) and $K=\min(\omega/2\pi,1)$ (unitary ensemble), predicts a ratio $3.95$ and
root mean squares $5.67$, $8.02$, $11.36$ (unitary) and $22.5$, $31.7$, $44.9$ (Poisson) on the last three windows, against the
measured $5.38$, $7.46$, $11.56$ and $20.8$, $28.3$, $41.8$ for one seed.
% NOTE: the predicted values come from a post-processing step described in spec_kernelfft.py; the stored output
% spec_kernelfft.out lists the kernel spectrum only. Store the post-processing output before submission.

\begin{table}[ht]
\centering
\small
\begin{tabular}{lrrr}
\toprule
$n$-window & unitary eigenphases & Poisson bridge & ratio\\
\midrule
$[50,100)$ & $2.67\pm0.36$ & $8.7\pm2.5$ & $3.3$\\
$[100,200)$ & $3.76\pm0.34$ & $15.0\pm4.8$ & $4.0$\\
$[200,400)$ & $5.38\pm0.18$ & $19.3\pm3.4$ & $3.6$\\
$[400,800)$ & $7.40\pm0.33$ & $26.3\pm2.1$ & $3.6$\\
$[800,1600)$ & $10.68\pm0.50$ & $39.8\pm5.1$ & $3.7$\\
\bottomrule
\end{tabular}
\caption{Root mean square of $n(y_n-y_n^{\mathrm{lat}})$, mean and standard deviation over five seeds (scripts
\texttt{spec\_stats}, \texttt{spec\_stats\_summary}).}
\label{tab:noise}
\end{table}

% Section 7: Conjectures and open problems.
\section{Conjectures and open problems}\label{sec:conj}

The theorems of Sections~\ref{sec:equilibrium}--\ref{sec:hurwitz-thm} concern the lattice, Hurwitz lattices and shifted
Hurwitz lattices, apart from the potential theory of Section~\ref{sec:equilibrium}, which holds for general densities. This section
states what they and the evidence of Section~\ref{sec:examples} suggest for other node sets, and what is missing for proofs.

\subsection{The natural-index law}\label{sec:conj-law}
Call $\Lambda$ \emph{regular} if $\lambda_m^{1/p}$ has an asymptotic expansion $m+a-1+\sum_{i\ge1}\gamma_im^{-r_i}$ with
exponents $0<r_1<r_2<\dotsb\to\infty$, or if $\Lambda$ differs from such a set in finitely many elements. Then $Z_\Lambda$ is
meromorphic, $Z_\Lambda(0)$ and $Z_\Lambda(-j)$ are defined by continuation, and $w\mapsto Z_\Lambda(-w)-\zeta(-pw,a)$ has
poles only at points $q$ produced by the corrections $\gamma_im^{-r_i}$ (at $q=r_i/p$ to first order).

\begin{conjecture}[natural-index law]\label{conj:natural}
Let $\Lambda$ be regular and $N_\Lambda:=n+\frac\alpha2-Z_\Lambda(0)$. Then $n^pb_n/c=(n/N_\Lambda)^pF$, where $F$ has an
asymptotic expansion in powers of $1/N_\Lambda$ and of the spectral monomials below, with
\begin{enumerate}
\item[(a)] no term $N_\Lambda^{-1}$, and the term $f_2N_\Lambda^{-2}$, with $f_2$ as in~\eqref{eq:f2};
\item[(b)] for $j\ge1$ the term $pW(j)Z_\Lambda(-j)N_\Lambda^{-1-jp}$;
\item[(c)] for every pole $q\ne-1/p$ of $Z_\Lambda(-w)-\zeta(-pw)$, with residue $R_q$, the term
$pR_qW(q)\pi\cot(\pi q)N_\Lambda^{-1-pq}$, with the value $-(4c)^q\Lt(q)$ at $q\in-\frac12-\N_0$.
\end{enumerate}
Coefficients of coincident exponents add. The statement concerns the leading power of each datum in~(b) and~(c); it
says nothing about the terms of order $N_\Lambda^{-2-pq}$ (Remark~\ref{rem:sublead}).
\end{conjecture}

For the lattice, (a) and (b) are Theorem~\ref{thm:N6} (with $Z_\Lambda(0)=-\frac12$, $W(j)p\zeta(-jp)=\varepsilon_j$), and (c)
is empty. For Bessel zeros the whole expansion is Proposition~\ref{prop:bessel}. For Hurwitz lattices, synthetic
corrections, smoothed drums and the oscillators $q^N$ the evidence is in Section~\ref{sec:examples}. In the form~(b)--(c) the
two kinds of spectral data are the residues of one integrand,
\[
  \frac pn\cdot\frac1{2\pi i}\int_{\uparrow}-W(w)\,\pi\cot(\pi w)\,Z_\Lambda(-w)\,n^{-wp}\diff w ,
\]
at the integers (Theorem~\ref{thm:K}) and at the poles of $Z_\Lambda$ (Proposition~\ref{prop:M12}). The Weyl pole of
$Z_\Lambda$ is excluded: it is the leading term $n$ of the soft-edge condition, not a correction. A moved Weyl pole, as for
$\lambda_k=k^p+\kappa$, is included (\S\ref{sec:num-pole}). If a pole meets an integer, the two residues diverge separately
and their sum produces a term $n^{-1-jp}\log n$; one weak test at $p=m=\frac12$ is consistent with the predicted coefficient.

Two consequences make the law a spectral dictionary. Through $Z_\Lambda(0)$, the $1/n$ coefficient $d_1=-p\alpha/2+pZ_\Lambda(0)$
reads off the index offset of the spectrum, such as the constant $a$ of a Hurwitz lattice (\S\ref{sec:hurwitz}) or the
corner and curvature constant of a drum (\S\ref{sec:drums}). Through $Z_\Lambda(-j)$ and the poles, it reads off the
coefficients of the heat trace of the operator whose spectrum is $\Lambda$, weighted by $W(q)\pi\cot(\pi q)$.

\begin{remark}[subleading pole terms]\label{rem:sublead}
Each pole also contributes at order $n^{-2-pq}$, one order below its leading term, and this contribution is not removed
by $\cot(\pi q)=0$. For Bessel zeros at $p=2$ the McMahon correction has $q=\frac12$; its leading term vanishes, and its
subleading term at $n^{-3}$ cancels the zeta term $2pcZ_\Lambda(-1)$ exactly, which is why Proposition~\ref{prop:bessel} is a
pure index shift; a pure Hurwitz lattice with the same offset does carry the zeta term at $n^{-3}$. For the synthetic
family the subleading term is~\eqref{eq:S4} when $q_1\le1$:
\end{remark}

\begin{conjecture}\label{conj:S4}
For $\lambda_k=(k+\gamma k^{-m})^p$, $m>0$, $p\ne1$, $q_1=(m+1)/p\le1$, the part linear in $\gamma$ of the $n^{-2-m}$ coefficient
of $n^pb_n/c$ is $\gamma S$ with $S$ as in~\eqref{eq:S4}.
\end{conjecture}

Both terms of~\eqref{eq:S4} are of the type that Conjecture~\ref{conj:natural} does not describe: a change of node density
on the band, and an index shift of a first-order term. In Theorem~\ref{thm:N6} the same shift is visible as the factor
$(n/N)^p$ in front of every layer. This suggests that every first-order term is to be evaluated at the natural index,
plus a density term.

\subsection{The natural form to all orders}\label{sec:conj-even}
\begin{conjecture}\label{conj:S1Nprime}
For the lattice and $p\ne1$, $F:=(N/n)^pn^pb_n/c$ depends on $n$ only through $N^{-2}$ and the monomials
$X_j:=\varepsilon_jN^{-1-jp}$: $F=\sum_{\br}c_{\br}(N^{-2})\prod_jX_j^{r_j}$, with each $c_{\br}$ a power series.
\end{conjecture}

Theorem~\ref{thm:N6} proves the instances $\br=e_j$ to relative order $N^{-3}$ for every $j$ (and $N^{-4}$ for $j=1$) and
$\br=2e_1$ to relative order $N^{-3}$, with $c_{2e_1}(0)=(3p+1)/(2p)$, and Corollary~\ref{cor:even} proves the instance $\br=0$ to all
orders. Theorem~\ref{thm:indexshift} gives the leading coefficient of every layer, $c_{\br}(0)=(p+1+S)_{|\br|-1}/(\br!\,p^{|\br|-1})$,
$S=p|\br|_w$, and shows that each layer is $\prod_jX_j^{r_j}$ times a power series in $N^{-1}$; that this series is even in $N$
beyond the orders above is open. Theorem~\ref{thm:shift}(ii) shows that the layers $e_1$, $e_2$, $2e_1$ are natural in $N$ to all orders (their coefficients in
$N$ do not depend on the Szeg\H{o} exponent); their evenness beyond the orders of Theorem~\ref{thm:N6}(ii), and the layers of weight
$\ge3$, are open. By Remark~\ref{rem:parity} the parity part of the conjecture would follow from Conjecture~\ref{conj:diag}
(proved to order $M^{-4}$ in Proposition~\ref{prop:diag4}). The numbers
$(3p+1)/(2p)=1+(p+1)/(2p)$, like $\varepsilon_j$, do not depend on $\alpha$. A closed form for $c_0=F_p$ is known only at
$(p,\alpha)=(2,-\frac12)$, where it is $(1-x/4)^{-1/2}$ in $x=N^{-2}$. The hard-edge denominators $\prod_k((2k+1)p+2)$ of
Remark~\ref{rem:N6}(b) suggested a hypergeometric structure in $a'=\frac1p+\frac12$, but the series $F_p$ appears to be divergent:
evaluating the recursion to order $40$ at $120$--$170$ digits, at $(p,\nu)=(3,0.3)$, the coefficients are small up to $f_{16}\approx-0.031$
and then grow factorially with oscillating sign ($f_{30}\approx-9.0\cdot10^4$, $f_{40}\approx-1.9\cdot10^{12}$), and Borel--Pad\'e
transforms at six values of $(p,\nu)$ show a complex-conjugate pair of singularities (scripts in \texttt{checks/large\_order}). These
statements are numerical; they do not affect Theorem~\ref{thm:N6}, which concerns finite truncations.

\subsection{A general form: weighted Weyl-type spectra}\label{sec:conj-W}
The theorems of Sections~\ref{sec:shiftsec} and~\ref{sec:hurwitz-thm} suggest a single statement for general node sets, in which
the Szeg\H{o}-shifted model, the Toda flow and the pole terms of Conjecture~\ref{conj:natural} appear together. Let
$\sigma:=\sum_kw_k\delta_{1/\lambda_k}$ with:
\begin{enumerate}
\item[(W1)] distinct nodes $0<\lambda_1<\lambda_2<\dotsb\to\infty$ (eigenvalues of an operator with multiplicity are merged into one
node, and the multiplicity goes into the mass);
\item[(W2)] an index function $\mathfrak n$ with $\mathfrak n(\lambda_k)=k$ and an expansion, differentiable termwise,
$\mathfrak n(\lambda)=\lambda^{1/p}+\sum_{r\in\mathcal R}A_r\lambda^r+\mathfrak n_0+O(\lambda^{-M})$, with $\mathcal R\subset(-\infty,\frac1p)\setminus\{0\}$
finite for each $M$ (this excludes fluctuating counting functions, \S\ref{sec:noise});
\item[(W3)] masses $w_k=\lambda_k^{-s/p}v(\lambda_k)$ with $v(\lambda)=v_0\lambda^\tau(1+\sum_\rho b_\rho\lambda^{-\rho})$, $\rho>0$; we put
$s_{\mathrm{eff}}:=s-p\tau>1$ and, in this subsection, $\alpha:=(s_{\mathrm{eff}}-1-p)/p$;
\item[(W4)] $p\ne1$, and (W5) no pole $q$ below is an integer (otherwise a term $n^{-1-pq}\log n$ appears and nothing is claimed from
that order on).
\end{enumerate}
The spectral data are read from $(\Lambda,v)$: the node zeta function $Z(w)=\sum_k\lambda_k^{-w}$, continued by (W2), with
$Z(0)=\mathfrak n_0-\frac12$; the poles $q=-r$ of $Z(-w)$, $r\in\mathcal R$, with residues $R_q=qA_{-q}$ (the Weyl pole $q=-\frac1p$ is
excluded); the \emph{density factor} $D(x):=(\mathrm d\mathfrak n/\mathrm d\lambda)/(\frac1p\lambda^{1/p-1})$ at $x=1/\lambda$, which is exactly
$1-p\sum_qR_qx^{q+1/p}$ up to the error term of (W2); the \emph{mass factor} $V(x):=v(1/x)x^\tau/v_0$; the \emph{band tilt}
$\log V+\log D=:\sum_\rho B_\rho x^\rho$; and the \emph{Toda times}
\[
  t_j:=-\frac{2Z(-j)}j+B_j\quad(j\in\N),\qquad t_\rho:=B_\rho\quad(\rho\notin\N).
\]
For the lattice, $Z(0)=-\frac12$, $Z(-j)=\zeta(-jp)$, there are no poles, $D=V=1$, and $t_j$ are the times~\eqref{eq:todatimes}.

\begin{conjecture}\label{conj:W}
Under (W1)--(W5), with $\beta:=\alpha-2Z(0)-1$, $\nu:=\alpha$ and $N:=n+\frac{\beta+1}2=n+\frac\alpha2-Z(0)$,
\[
  \frac{n^pb_n}c\simeq\sum_{\br}\mu_n^{\br}\cD_{\br}(n;\beta,\nu)
  +\sum_{q\notin\N}pR_qW(q)\pi\cot(\pi q)\Bigl(\frac nN\Bigr)^p\!N^{-1-pq}
  -\sum_{\rho\notin\N}\frac{p\rho W(\rho)}2\,t_\rho\Bigl(\frac nN\Bigr)^p\!N^{-1-p\rho},
\]
with $\mu_{n,j}:=-t_jn^{-jp}$ and the layers~\eqref{eq:decoupled} of the decoupled model $M(\beta,\nu)$. The claim is: the integer and Toda
layers to every order, each pole and fractional term to relative order $O(N^{-2})$; products of two pole or fractional terms, and of
such a term with a Toda time, are not claimed.
\end{conjecture}

The first sum is the decoupled model flowed along the Toda hierarchy (Theorem~\ref{thm:toda}); for $\rho=j\in\N$ the first-order
Toda response is $\varepsilon_j=pW(j)Z(-j)=-\frac{pjW(j)}2t_j$ at $B_j=0$, so there is one linear response $-\frac{p\rho W(\rho)}2t_\rho$ for
every time, integer or not. The lattice is Theorem~\ref{thm:N6}; Hurwitz lattices are Theorem~\ref{thm:H} ($\beta=\alpha+2\delta$);
the shifted Hurwitz lattices are Theorem~\ref{thm:kappa}, at first order in the poles; and Conjecture~\ref{conj:natural} is the first
order of Conjecture~\ref{conj:W} in each datum. For the Bessel zeros $\lambda_k=(j_{\nu',k}/\pi)^2$, $s=2$, the poles are half-integers,
where $\cot(\pi q)=0$, and the Toda times computed from the spectral data vanish ($|t_1|<10^{-21}$, $|t_2|<10^{-16}$ at $\nu'=0,1$, the
precision of the regularised zero sums); so Conjecture~\ref{conj:W} predicts that the measure is the bare model $M(\nu'-1,-\frac12)$,
which is what Proposition~\ref{prop:bessel} says. In this reading the cancellation at $n^{-3}$ of Remark~\ref{rem:sublead} is the
statement $t_1=0$. The tests of Conjecture~\ref{conj:W}, with predictions fixed before the measures were computed, are in
\S\ref{sec:num-W}.

\begin{remark}[multiplicities]\label{rem:mult}
For an operator spectrum with multiplicities the zeta function in Conjecture~\ref{conj:W} is that of the \emph{distinct} nodes, and the
multiplicity is a mass factor. With multiplicities counted, already the Weyl exponent is wrong: for the eigenvalues $l(l+1)$ of the
sphere, with multiplicity $2l+1$, the distinct nodes have $p=2$, and the fits give $n^{-2}b_n\to c(2)$ with $d_1=-2.00000000$ at $s=4$,
while the spectral zeta function with multiplicities has Weyl exponent $1$. For these nodes $Z(-j)=0$ for every $j\ge1$: with
$t=l+\frac12\in\frac32+\N_0$, $(t^2-\frac14)^j=\sum_mc_mt^{2m}$, and $\zeta(-2m,\frac32)=\zeta(-2m,\frac12)-2^{-2m}=-2^{-2m}$ for $m\ge0$
(since $\zeta(s,\frac12)=(2^s-1)\zeta(s)$ and $\zeta(-2m)=0$ for $m\ge1$), so $Z(-j)=-(\frac14-\frac14)^j=0$. The mass and density factors
cancel, so Conjecture~\ref{conj:W} predicts the bare model $M(1,0)$ at $p=2$; the fits agree at $n^{-1}$, $n^{-2}$, $n^{-3}$ to
$8\cdot10^{-8}$, and the alternatives ``mass tilt only'' and ``density tilt only'' are excluded.
\end{remark}

\begin{remark}[the dressing of pole terms]\label{rem:Wprime}
At $n^{-4}$ the sphere misses $M(1,0)$ by $+0.23132$ (spread $3\cdot10^{-4}$), outside the precision claimed in
Conjecture~\ref{conj:W}. The probes $k^2\pm\frac14$ (moved Weyl pole $q=\frac12$, residue $R=\mp\frac1{16}$) miss by $\mp0.23131$: the term is
odd in the residue, so it is of first order in the pole, one order below its leading power $n^{-1-pq}$. It is explained by the following
refinement, found after this failure: each pole term carries the dressing of the zeta layer continued in the index,
\[
  pR_qW(q)\pi\cot(\pi q)\Bigl(\frac nN\Bigr)^pN^{-1-pq}\bigl[1+g(q)N^{-2}+O(N^{-4})\bigr],\qquad g(w):=A(w)+B(w)\nu^2,
\]
with $A_j$, $B_j$ of Theorem~\ref{thm:A} read as functions of $w$ (and the same for the fractional times). By Remark~\ref{rem:N6}(d),
$g$ has a simple pole at $w=\frac12$, where $\cot(\pi w)$ vanishes, with $\Res_{w=1/2}g=-(p+2)^2(p+4)(1-4\nu^2)/(256p)$. So a pole at
$q=\frac12$, whose leading term vanishes, contributes the finite term
\[
  R_{1/2}\,W(\tfrac12)\,\frac{\pi^2(p+2)^2(p+4)(1-4\nu^2)}{256}\,n^{-3-p/2},
\]
which is $\frac{3\pi^2}8R_{1/2}=\mp0.2313188$ for the probes at $p=2$, $\nu=0$ (fits $-0.2313059$ with spread $7\cdot10^{-4}$, and $7.7762408$
against $7.7762407$ for the sphere's $n^{-4}$ coefficient). It vanishes at $\nu=\pm\frac12$ for every $p$, which is why the Lommel measures
($\nu=-\frac12$) and the earlier tests with $q=\frac12$ poles (all at $\nu=\pm\frac12$) show nothing. The refinement has one test made after
it was formulated and before the measures were computed: for $k^2+\frac14$ at $s=3.5$ ($\nu=\frac14$) it predicts $+2.9296023$ at $n^{-4}$, the
fit gives $+2.9296226$ (spread $3.6\cdot10^{-3}$), and the unrefined value $+3.1030915$ is excluded; at $s=2$ ($\nu=-\frac12$) it predicts no
extra term, as found. The dressing at poles $q\ne\frac12$ is untested. Its premise, that pole terms are the analytically continued
zeta layer and nothing else, holds in the decoupled model with the pole structure of \S\ref{sec:kappa} (the tilt and the lip factor come
from one function), and Theorem~\ref{thm:kappa} proves it at leading order.
\end{remark}

\begin{remark}[a hard-edge companion for branched tilts]\label{rem:companion}
For a branched tilt $e^{-\mu z^\rho}$, $\rho\notin\mathbb Z$, that does not come from a pole (for example a mass factor), the gauge of
Lemma~\ref{lem:gauge} splits near $0$ into an analytic part, whose Taylor data are those of the symbolic-$j$ layer of
Lemma~\ref{lem:N6d}(b) at $j=\rho$, and a remainder $\pm i\tan(\pi\rho)z^\rho$ in $\C_\pm$. In the formal model the remainder acts only through
the lip jumps at the Bessel scale, and a Mellin computation with the products of the Bessel functions $K_\nu$, $K_\nu'$ gives, besides the
symbolic-$j$ layer at $j=\rho$, a first-order \emph{companion} term proportional to
\[
  (p+2)^{2\rho+1}\Phi_\nu(\rho)\Bigl(\frac nN\Bigr)^{(\rho+1)p}N^{-2\rho-2}\sum_{k\ge0}c_kN^{-k},\qquad
  \Phi_\nu(\rho)=\frac{\Gamma(-1-2\rho)\Gamma(\nu+\rho+1)}{\Gamma(-\rho)^2\Gamma(\nu-\rho)},
\]
with $c_0=(p\rho+2\rho+1)/2$, $c_1=c_3=0$, and each $c_k$ given by a finite exact computation; it is $\beta$-free in $N$. At
$p=-1$, where $M(\nu-1,\nu)$ is the Laguerre weight, this reproduces the exact branched Laguerre series through $N^{-4}$ identically in
$(\nu,\rho)$, and pre-registered residues at order $N^{-7}$ at $p=\frac32$ and $3$ hold. The companion is of relative order $N^{-2\rho-1}$; it is
beyond the precision claimed in Conjecture~\ref{conj:W} for $\rho>\frac12$, which covers all the tests of \S\ref{sec:num-W}, and would
have to be added for $\rho<\frac12$. That these formal layers are asymptotic expansions for actual measures with branched mass factors is
not proved. Pole terms have no companion: for them the tilt and the lip factor combine to an analytic gauge.
\end{remark}

\begin{problem}\label{prob:kernel2}
Prove that the $N^{-2}$ dressing of a pole term at $q$ is $g(q)$ with $g(w)=A(w)+B(w)\nu^2$ the continuation of $g_j$ (Remark~\ref{rem:Wprime}),
or find the correct second-order edge kernel. By Remark~\ref{rem:N6}(d) it should combine a polynomial in $\theta=u\frac{\diff}{\diff u}$
acting on $(1-u)^{-1/2}$ with a multiple of $(1-u)^{1/2}$. The data of Table~\ref{tab:subpole} and of drums are test cases.
\end{problem}

\subsection{What a proof needs}\label{sec:conj-proof}
The derivation of Conjecture~\ref{conj:natural}(c) from the analysis of \cite{PaperII} needs that analysis for a node set
other than the lattice, with three ingredients: the $g$-function of the constrained equilibrium of the $n$-dependent
continuum node density (Theorem~\ref{thm:Q}); the band factor given by the Cauchy integral of the jump of
$\log(1-e^{-2\pi it_\kappa})$ for the index function $t_\kappa(z)=\kappa(n^p/z)$ (Proposition~\ref{prop:pert}); and parametrix
corrections whose first-order dependence on the perturbation is $O(n^{-2-m})$. Under these, the first-order change of
$n^pb_n/c$ at order $n^{-1-m}$ is the edge shift of~\eqref{eq:Q} for the perturbed density, and
Proposition~\ref{prop:M12} gives its value. For Hurwitz lattices the index normalisation of the outer function is the gauge
$(-z)^{-\delta}$ of \S\ref{sec:H-mechanism}, and the shift $-p\delta$ of $d_1$ arises as the Szeg\H{o} shift $\beta'=\alpha+2\delta$
(Theorem~\ref{thm:H}); for $\lambda_k=(k+\delta)^p+\kappa$, $p>1$, the lattice $g$-function still suffices, and the pole term comes from a
branched band tilt rather than from the equilibrium (Theorem~\ref{thm:kappa}). For a general Weyl-type node set the template of
Theorem~\ref{thm:H} needs four pieces: an outer scalar built from a regularised determinant of $\Lambda$ composed with an analytic
index map, whose normalising gauge has the exponent $-Z_\Lambda(0)-\frac12$ at the hard edge; band-factor moments equal to
$Z_\Lambda(-j)/(jn^{jp})$; the $n$-dependent equilibrium above; and the linear response of Lemma~\ref{lem:E37} for branched tilts,
including the companion of Remark~\ref{rem:companion}. For $p<1$ the lattice $g$-function cannot carry the index offset of a shifted
lattice (\S\ref{sec:kappa}), and the $n$-dependent equilibrium is needed already there.

\begin{remark}[a proved instance of the per-node law]\label{rem:S2}
Conjecture~\ref{conj:natural}, read at leading order per node, says that one node with $u=4c\lambda n^{-p}<1$ outside the band
contributes $-p(1-u)^{-1/2}/n$ to $y_n$ when it is removed: this is the derivative of the soft-edge condition of
Remark~\ref{rem:Q}(c). For one atom $k_0^{-p}$ of $\sigma_{p,s}$ this is proved, uniformly for $u_n=4c(k_0/n)^p\le1-\eta$, with
error $O(n^{-1-\epsilon})$, $\epsilon=\min(2p,1)$ for $p<1$ and $\epsilon=1$ for $p\ge1$, from the rank-one update of the Hankel
matrices and a contraction argument for the minimal solution of the recurrence at the deleted node; the proof uses only the
first-order expansions of $a_n$ and $b_n$.
% TODO: this is Theorem S2' of research/G1/RH/spectral/SPEC.md Sec. 3.2 (reviewed), not in any paper; outside the
% sources requested for this paper. Either include it with proof or cite it once it is written up.
\end{remark}

\subsection{Further problems}\label{sec:conj-further}
(a) \emph{The case $p=1$.} Drums and the measures $\sum_kk^{-s}\delta_{1/k}$ have $p=1$, where the two cotangents of
Theorem~\ref{thm:K} merge and \cite{PaperII} uses a different hard-edge analysis; Section~\ref{sec:drums} indicates that the
dictionary persists. (b) \emph{Higher layers.} The layers $\mu_1\mu_2$, $\mu_2^2$, $\mu_1^3$ control the coefficients below
$n^{-2-3p}$ and are not computed, and their shift (Theorem~\ref{thm:shift}(ii) for weight $\ge3$) is open; it would follow from the
linear independence of the products $\prod_j\zeta(-jp)^{r_j}$ over $\mathbb Q(p)$, or from a version of Lemma~\ref{lem:frame} for the
tilted measures of Theorem~\ref{thm:toda}, in which the $\mu_j$ are free parameters. (c) \emph{Irregular spectra.} For spectra with displacements of the order of the spacing
there is no deterministic expansion; Section~\ref{sec:noise} suggests that the variance of the irregular part is
$n\int\kappa(\omega)S_h(\omega)\diff\omega/\pi$ with a universal kernel spectrum $\kappa$. A closed form for $\kappa$, presumably from
the Airy kernel of the soft edge, is open, as is a proof that rigidity at wavelengths of two to forty spacings sets the
noise.

% Section 10: Relation to the literature (from research/G1/RH/spectral/NOVELTY_JZC.md; the bibliographic data of the entries
% added for this section were relayed second-hand and must be checked against the sources before submission; DixitKumar:
% statement quoted in research/G1/RH/CITATIONS.md (Hurwitz-DK); metadata from the arXiv abstract page via WebFetch, 2026-10-22).
% Rule R review B of JZC_M1 (B6, 2026-10-23): Rakhmanov, DragnevSaff and VanAsscheVanBaelen are quoted with hypotheses in
% research/G2/CITATIONS.md, entries (R96), (DS97), (VAVB); BeckermannKuijlaars2002 is not quoted in any CITATIONS.md (hence
% "according to"); BeckermannKuijlaars2001 has not been read.
% Rule L, 2026-10-23: rewritten from the drafted passages of NOVELTY_Q.md Sec. 3 (constrained equilibrium), NOVELTY_TODA.md Sec. 5
% (shifted index, Toda) and NOVELTY_ZETA.md Sec. 6.2 (spectral data, large order). Every work added here was read only through a
% summarising fetch tool (WebFetch); none is quoted in a CITATIONS.md, hence "according to" and "we have not found" throughout.
% Drafted citations not used (no bib entry made): Dickinson, Grosjean, Carlitz/Goh-Wimp, Koelink-Van Assche, Parker et al.,
% Avdoshkin-Dymarsky-Smolkin, Naber-Bruck-Costello, Ercolani-Waters.
\section{Relation to the literature}\label{sec:literature}

\emph{Constrained equilibrium.} Existence, uniqueness and the structure (band, saturated region, balayage) of equilibrium
problems with an upper constraint are classical \cite{Rakhmanov,DragnevSaff}; these results assume a finite constraint
(Dragnev--Saff) or a compact interval (Rakhmanov) and are not used in the proof of Theorem~\ref{thm:Q}, which establishes
existence and uniqueness directly. For constraints of finite mass with compact support and no external field, the structure of
Theorem~\ref{thm:Q} appears in the analysis of conjugate-gradient convergence by Beckermann and Kuijlaars
\cite{BeckermannKuijlaars2001}. Let $\mu_t$ denote the constrained equilibrium measure of mass $t$ for a constraint $\sigma$ on $[A,B]$.
According to \cite[Lemma~1.5]{BeckermannHelart}, which follows \cite[Lemma~3.1(a)]{BeckermannKuijlaars2001}, if
$\sqrt{(x-A)(B-x)}\,\sigma'(x)$ vanishes at $A$ and increases on $[A,B]$, then the support of $\sigma-\mu_t$ is a single interval
$[a(t),B]$, and $a(t)$ solves an equation that becomes~\eqref{eq:Q} after the reflection $y=B-x$; for $Y<\infty$ this criterion is
implied by the hypothesis of Remark~\ref{rem:Q}(a). According to \cite{BeckermannGuettel,BeckermannKuijlaars2002},
\cite[Theorem~2.1]{BeckermannKuijlaars2001} and its proof also give a constrained analogue of the Buyarov--Rakhmanov formula
\cite{BuyarovRakhmanov}, in which $\mu_t$ is a superposition of the equilibrium measures of the decreasing sets
$\operatorname{supp}(\sigma-\mu_\tau)$, $\tau<t$. In the present setting this superposition is
$\nu(x)=\int_a^Y|F'(b)|\,\omega_{[0,b]}(x)\diff b$, with $F$ as in~\eqref{eq:Q} and $\omega_{[0,b]}(x)=1/(\pi\sqrt{x(b-x)})$ the
equilibrium density of $[0,b]$; we have checked numerically, but not proved, that this agrees with~\eqref{eq:nu}, which is then a closed
form of it. What is new here, as far as we know, is the constraint of infinite mass at a hard edge, with $Y=\infty$ allowed. In this
case the finite-mass theory \cite{Rakhmanov,DragnevSaff} does not apply, and existence, uniqueness and the strict outer inequality are
proved directly. Also new, as far as we know, are the single-integral form~\eqref{eq:nu} and the edge constant of
Remark~\ref{rem:Q}(b). Corollary~\ref{cor:Q1} is an elementary evaluation of an endpoint equation of this known type, for the lattice
problem of \cite[Theorem~1.3]{PaperI}. We have not been able to consult \cite{BeckermannKuijlaars2001} itself. The fully saturated
equilibrium of \cite{VanAsscheVanBaelen} ($y\sigma$ constant on a $q$-lattice) is the boundary case of hypothesis~(H2). We have not found
Proposition~\ref{prop:M12}, the expansion of the soft-edge root along the poles of the node zeta function; its termwise Mellin
asymptotics is a standard technique.

\emph{Spectral data and Jacobi parameters.} Orthogonal polynomials whose moments are spectral zeta values at positive integers are
classical for Bessel zeros: the Lommel polynomials are orthogonal on the reciprocal zeros \cite{VA} (Proposition~\ref{prop:bessel}),
and an analogue for the zeros of $J_\nu'$ is given in \cite{ChungLeePark}. The Tricomi--Carlitz polynomials, orthogonal on the points
$\pm(k+\alpha)^{-1/2}$ \cite{WuLinXuZhao}, give a further family with nodes accumulating at $0$ and explicit recurrence coefficients;
from the three-term recurrence as stated there, the lattice offset enters only through the index $n+\alpha$, as it does in
Theorem~\ref{thm:H} (at $p=1$ in the variable $x^2$, outside the range of our theorems). Positive Jacobi matrices with compact inverse
are studied in \cite{StovicekSwiderski}, where the asymptotics of the polynomials at fixed argument are expressed through a regularised
characteristic function; according to our reading, asymptotics of the Jacobi parameters are not treated there. We know of no expansion
of recurrence coefficients in terms of spectral zeta values at non-positive arguments, as in Theorems~\ref{thm:K} and~\ref{thm:N6}. For
continuous weights it is known that singularities of the weight produce identifiable power-law or oscillating terms in the recurrence
coefficients \cite{Magnus,FoulquieMoreno}; the pole terms of Proposition~\ref{prop:M12}, Theorem~\ref{thm:kappa} and
Conjecture~\ref{conj:W} may be viewed as an instance of this principle at the accumulation point of a discrete spectrum, with the poles
of $Z(-w)$ as singularity data. Subleading corrections to Lanczos coefficients have recently been derived from the endpoint data of
continuous equilibrium densities, including a hard-edge $n^{-2}$ term \cite{MuruganVanZylWatanabe}, which we have not compared with
$f_2$. We have not found a reading of Weyl-law or heat-trace invariants (perimeter, corners) from Jacobi or Lanczos coefficients.

\emph{Expansions in a shifted index.} For varying weights $e^{-nV}$ the recurrence coefficients have expansions in even
powers of $1/n$ after a half-index shift; the odd terms are removed by a reindexing symmetry of the string equations
\cite{BleherIts}, by the structure of the local parametrices \cite{KuijlaarsTibboel}, and for even weights the coefficients are
generating functions of maps \cite{EMP}. For Jacobi-type weights on $[-1,1]$ the first correction vanishes \cite{KMVV}, and an exact pairing of
$n+\alpha/2$ with $\alpha^2$ occurs for semiclassical Laguerre weights \cite{FilipukVanAsscheZhang} (Section~\ref{sec:shift-thm}). The present setting is a
fixed discrete measure without a string equation: the shift involves the spectral constant $Z(0)$, and the evenness comes with the zeta
layer and its own even dressing (Theorem~\ref{thm:N6}, Conjecture~\ref{conj:S1Nprime}). The grading argument in the proof of
Theorem~\ref{thm:reflect} is that of \cite[Lemma~4.1]{KuijlaarsTibboel}, and the theorem plays the role of the reindexing symmetry of
\cite{BleherIts}. What we have not found in the literature is the decoupling of the Szeg\H{o} exponent from the Bessel order, the frame
identity~\eqref{eq:FI}, and its proof from a second steepest-descent frame at a non-integer index (Lemma~\ref{lem:frame}), a gauge
$z^{\gamma\sigma_3}$ of the exact Riemann--Hilbert problem combined with a change of frame. The centring of $b_n$ at $N$ and of $a_n$ at
$N+\frac12$ (Proposition~\ref{prop:diag4} and Conjecture~\ref{conj:diag}) is the arrangement of \cite[Theorem~5.2]{BleherIts}.

\emph{Toda flows.} The Toda hierarchy for exponential tilts (Lemma~\ref{lem:todatilt}) is classical \cite{VanAsscheToda}. Re-expanding
the Toda equations with shifted indices in $1/n$ to generate the coefficients of an asymptotic expansion is the method of
\cite[Section~4]{EMP} for varying weights; Theorem~\ref{thm:toda} applies it to a fixed measure, with the zeta values as base times,
and transfers the identities from the exact flows by a category argument. The mechanism of Theorem~\ref{thm:indexshift} is the
genus-zero one of random-matrix map enumeration \cite{Ercolani,EMP,Bouttier}, and the role of the endpoints as Riemann invariants of
the dispersionless Toda hierarchy is known \cite{KonopelchenkoMAM}; Section~\ref{sec:indexshift} gives the references in place. The
new ingredients are the setting (a fixed measure accumulating at a hard edge, real and pole exponents, zeta-valued times), not the
mechanism.

\emph{Hurwitz zeta values.} The band-factor moments of a Hurwitz lattice come out of the analysis as Fourier series, and their
identification with $\zeta(-jp,a)$ is Hurwitz's formula, which we quote from \cite[Theorem~1.1]{DixitKumar}; there it is attributed to
Hurwitz and quoted from Titchmarsh's monograph, which we have not consulted. We have not found the index transfer of
Section~\ref{sec:hurwitz-thm}, a change of the Szeg\H{o} exponent at fixed Bessel order, or the statements of Theorems~\ref{thm:H}
and~\ref{thm:kappa} in the literature; the closest exact instances are the Lommel and Tricomi--Carlitz recurrences above, in which an
offset of the nodes enters only through the index.

\emph{Large order.} Divergent $1/N$ expansions whose large-order behaviour is governed by an instanton action are familiar for matrix
models with varying weights \cite{MarinoSchiappaWeiss}, and Euler--Maclaurin remainders are resurgent \cite{CostinGaroufalidis}; we have
not found a large-order analysis for a fixed discrete measure, and the singularities found numerically in
Section~\ref{sec:conj-even} are not identified.

\appendix
% Appendix A: Reproducibility.
\section{Reproducibility}\label{app:repro}

All scripts are in the directory \texttt{research/G1/RH/spectral/checks} of the repository named in the data
availability statement, next to their stored outputs (\texttt{.out}). They use the project's toolkit (quadruple-precision
recurrence coefficients, cached under names that include a hash of the node array and of $(p,s,K)$, so that a changed
measure never reuses a cache), \texttt{mpmath} and \texttt{sympy}. The scripts that test a prediction end with assertions
on it and print a line of the form ``ALL $n$ CHECKS PASSED'' or ``ALL IDENTITIES HOLD''.

\subsection*{Exact computer algebra (Lemma~\ref{lem:N6d})}
The outputs begin with the SHA-256 hashes of the scripts that produced them:
\begin{center}
\small
\begin{tabular}{ll}
\toprule
script & SHA-256\\
\midrule
\texttt{m4\_rh\_exact.py} & \texttt{1e2904089fa3053cd6a7249dd4cd6149}\\
 & \texttt{2be3bb9c90b46c0932ff2507c9d9ffcc}\\
\texttt{m4\_zeta\_exact.py} & \texttt{3df62305bba0a4801a767c252eeee000}\\
 & \texttt{31e499f7c80a36946542fea6cd834041}\\
\texttt{m4\_rh\_exact2.py} & \texttt{143a308b8744bfccf3110c3a3f4a60a8}\\
 & \texttt{f192268b0f65f78b64bf8bd76c116521}\\
\texttt{m4\_zeta\_exact2.py} & \texttt{badda6a534fbf621c2bbef3eca6c7822}\\
 & \texttt{b26d13493a655a1655860449a8bab9f4}\\
\bottomrule
\end{tabular}
\end{center}
The runs and what they assert, identically in the symbols:
\begin{center}
\small
\begin{tabular}{lp{4.2cm}p{5.6cm}}
\toprule
output & configuration & content\\
\midrule
\texttt{m4\_rh\_exact\_k4.out} & $K=4$ & Lemma~\ref{lem:N6d}(a) to $f_4$\\
\texttt{m4\_rh\_exact\_k6.out} & $K=6$ & Lemma~\ref{lem:N6d}(a): $d_1$, $d_2$, $f_1=f_3=f_5=0$, $f_2,f_4,f_6$\\
\texttt{m4\_zeta\_exact.out} & $K=5$, $j=1,\dots,5$ & Lemma~\ref{lem:N6d}(b) for integer $j$; $g_{1,4}$\\
\texttt{m4\_zeta\_exact2\_A.out} & $K=4$, $j$ symbolic & Lemma~\ref{lem:N6d}(b) for all $j$\\
\texttt{m4\_zeta\_exact2\_B.out} & $K=5$, second order in $\mu_1$ & Lemma~\ref{lem:N6d}(c)\\
\bottomrule
\end{tabular}
\end{center}
The first two runs also assert the coefficient identities $b_k=a_k+(k-\frac12)a_{k-1}$ (the termwise derivative of
\cite[(10.40.2)]{DLMF}) and the recursion for $u_k$, $v_k$ of \cite[(9.7.2)]{DLMF} for $k\le8$; the symbolic run \texttt{m4\_zeta\_exact2\_A} asserts the
Taylor data of $Q_j/\binom{2j}j$ against the integer polynomials for $j\le8$.

\subsection*{Floating-point scripts}
\begin{center}
\small
\begin{tabular}{lp{10cm}}
\toprule
used in & scripts\\
\midrule
Remark~\ref{rem:Q-checks} & \texttt{m1\_theoremQ}, \texttt{m1\_edge}\\
\S\ref{sec:kernel} & \texttt{spec\_kernel} (parts K0, K, P, Q), \texttt{m2\_thmK\_plt1}\\
\S\ref{sec:bessel} & \texttt{spec\_lommel}\\
Remark~\ref{rem:N6}(b) & \texttt{m4\_rh4}, \texttt{m4\_f8\_denominator}\\
\S\ref{sec:num-lattice} & \texttt{m4\_natural\_index}, \texttt{m4\_f4\_lattice}, \texttt{m4\_zeta\_lattice}\\
\S\ref{sec:hurwitz} & \texttt{spec\_hur}, \texttt{spec\_hur\_d2}, \texttt{spec\_shift\_hi}, \texttt{spec\_hur\_small}\\
\S\ref{sec:num-pole} & \texttt{spec\_pole}, \texttt{spec\_gamma}, \texttt{spec\_kappa}, \texttt{spec\_kappa\_fit},
  \texttt{spec\_mobius}, \texttt{spec\_mobius\_pred}, \texttt{spec\_resonance}\\
\S\ref{sec:num-subpole} & \texttt{spec\_subpole}\\
\S\ref{sec:drums} & \texttt{spec\_corners}\\
\S\ref{sec:quartic} & \texttt{spec\_quartic}, \texttt{spec\_qN}, \texttt{spec\_homog}\\
\S\ref{sec:shift-proof} & \texttt{m4\_shift\_exact}, \texttt{sigma\_frame\_exact} (exact), \texttt{m4\_shift\_num},
  \texttt{sigma\_frame\_num}\\
\S\ref{sec:toda} & \texttt{m4\_toda}, \texttt{m4\_diag}\\
\S\ref{sec:conj-even} & \texttt{large\_order/large\_order}, \texttt{large\_order/borel}\\
Theorem~\ref{thm:H} & \texttt{m3a\_hurwitz\_formula}, \texttt{m3a\_e37\_decoupled}, \texttt{m3a\_g5\_exact} (exact),
  \texttt{m3a\_hd\_model}\\
Theorem~\ref{thm:kappa} & \texttt{w\_outer\_e2e}\\
\S\ref{sec:num-W}, \S\ref{sec:conj-W} & \texttt{w\_predict}, \texttt{w\_test}, \texttt{w\_resid}, \texttt{w\_probe}, \texttt{w\_probe2}\\
\S\ref{sec:noise} & \texttt{spec\_jitter}, \texttt{spec\_window}, \texttt{spec\_edgekernel}, \texttt{spec\_stats},
  \texttt{spec\_stats\_summary}, \texttt{spec\_noiselin}, \texttt{spec\_kernelfft}\\
\bottomrule
\end{tabular}
\end{center}

The shell script \texttt{run\_all.sh} in the same directory regenerates the outputs from scratch in dependency order
(readers of other scripts' caches last), including the five exact runs (about $35$ minutes) and the coefficient
computations (one to two and a half minutes per measure with $n\le1600$ on two cores). The scripts \texttt{spec\_quartic},
\texttt{spec\_qN}, \texttt{spec\_corners}, \texttt{spec\_edgekernel}, \texttt{spec\_noiselin} and \texttt{spec\_kernelfft} are not in
\texttt{run\_all.sh} and are run individually.
The scripts listed for Sections~\ref{sec:shiftsec} and~\ref{sec:hurwitz-thm}, \S\ref{sec:conj-even} and \S\ref{sec:num-W} are
not in \texttt{run\_all.sh} either and are run individually; the measures of \S\ref{sec:num-W} are cached outside the repository.
% TODO: add the scripts of Sections 6, 7, 8.5, 9.2 to run_all.sh (the W caches live in tmp/w_cache, gitignored).
% TODO: add the six scripts above to run_all.sh, and store the post-processing of spec_kernelfft (the variance
% integral quoted in Section 6.8) as an output, before the tag is made.


\section*{Acknowledgements}
The research and the exposition were developed with the assistance of an AI system (Claude, Anthropic); the author
takes responsibility for the content.

\section*{Data and code availability}
{\raggedright The scripts that produced the exact identities of Section~\ref{sec:natural} and the numbers of
Section~\ref{sec:examples}, with their outputs, are available in the repository
\url{https://github.com/speagleg/jacobi-asymptotics}; Appendix~\ref{app:repro} lists them.
% TODO: tag (v1.0-paper4) to be fixed when the paper is frozen; the branch at drafting time is plan-e2-d2.
\par}

\bibliographystyle{amsplain}
\bibliography{refs}

\end{document}
